% Organisation de la cellule et échanges cellulaires — SVTEEHB Tle C — Cours signé OnBuch+
% Programme officiel MINESEC. Compiler avec : tectonic N01-echanges-cellulaires.tex
\newcommand{\DOCMATIERE}{SVTEEHB}
\newcommand{\DOCNIVEAU}{Tle C}
\newcommand{\DOCMODULE}{Module I : Le Monde Vivant}
\newcommand{\DOCLECON}{N01}
\newcommand{\DOCTITRE}{Organisation de la cellule et échanges cellulaires}
% OnBuch+ — préambule « cours signés » ENRICHI (compilé avec tectonic / XeLaTeX)
% Système visuel « Pop » d'OnBuch+ : fond crème (jamais blanc), bordures encre
% épaisses, ombres « dures » décalées, titres Archivo Black en capitales.
% Version MATHÉMATIQUES : graphes (pgfplots/tikz), boîtes pédagogiques
% (définition, propriété, méthode, attention, le savais-tu, exercice résolu,
% exercice, TP, bilan), page de garde et signature OnBuch+ ancrée en bas.
%
% Macros attendues dans main.tex AVANT \input{preamble} :
%   \DOCMATIERE  \DOCNIVEAU  \DOCMODULE  \DOCLECON (numéro)  \DOCTITRE
\documentclass[11pt]{article}

\usepackage{fontspec}
\usepackage[french]{babel}
\usepackage[a4paper,margin=2cm,top=3.4cm,bottom=2.8cm,headheight=1.4cm,footskip=1.3cm]{geometry}
\usepackage{amsmath,amssymb,amsfonts}
\usepackage{mathtools} % \xmapsto, \coloneqq... souvent utilisés par les modèles
\usepackage{cancel} % simplifications barrées (bilans de forces)
% \abs{z} (module, valeur absolue) et \norm{u} (norme), utilisés par les modèles
\providecommand{\abs}{}\let\abs\relax\DeclarePairedDelimiter{\abs}{\lvert}{\rvert}
\providecommand{\norm}{}\let\norm\relax\DeclarePairedDelimiter{\norm}{\lVert}{\rVert}
\usepackage[table]{xcolor} % [table] : fournit aussi \rowcolors (alternance de lignes)
\usepackage{pagecolor}
\usepackage{tikz}
\usetikzlibrary{arrows.meta,positioning,calc,shapes.geometric,shapes.misc,shapes.arrows,decorations.pathmorphing,decorations.pathreplacing,decorations.markings,patterns,shadows,mindmap,trees,fit,backgrounds,matrix,intersections,angles,quotes}
\usepackage{pgfplots}
\usepackage[version=4]{mhchem} % équations nucléaires \ce{^{235}_{92}U}
\usepackage[european,siunitx]{circuitikz} % schémas électriques
\pgfplotsset{compat=1.18}
\usepgfplotslibrary{fillbetween,groupplots} % aires entre courbes (calcul intégral)
\usepackage{enumitem}
\usepackage{fancyhdr}
% [most] charge toutes les bibliothèques tcolorbox (listings, minted, raster,
% documentation...) et gonfle inutilement la mémoire TeX sur un document long
% (~300 boîtes) : on ne charge que ce qui est réellement utilisé.
\usepackage[skins,breakable]{tcolorbox}
\usepackage{graphicx}
\usepackage{colortbl}
\usepackage{booktabs}
\usepackage{tabularx}
\usepackage{multirow}
\usepackage{diagbox} % case d'en-tête diagonale des tableaux à double entrée
% Colonnes extensibles centrée / gauche / droite (C, L, R), souvent utilisées par les modèles
\newcolumntype{C}{>{\centering\arraybackslash}X}
\newcolumntype{L}{>{\raggedright\arraybackslash}X}
\newcolumntype{R}{>{\raggedleft\arraybackslash}X}
\usepackage{array}
\usepackage{multicol}
\usepackage{siunitx}
% parse-numbers=false : accepte aussi \SI{2,5 \times 10^{-3}}{...}
\sisetup{locale=FR,output-decimal-marker={,},group-minimum-digits=4,parse-numbers=false}
\DeclareSIUnit\ohm{\ensuremath{\symup{\Omega}}} % U+2126 absent de la police
\DeclareSIUnit\division{div} % réglages d'oscilloscope (ms/div, V/div)
\DeclareSIUnit\tr{tr} % tours (vitesse de rotation, ex. \SI{1500}{\tr\per\minute})
\DeclareSIUnit\molar{M} % concentration molaire (ex. \SI{3}{\molar}, \SI{1}{\micro\molar})
\DeclareSIUnit\an{an} % année (le modèle confond parfois avec \year, un compteur TeX) — ex. \SI{5}{\kilo\meter\per\an}
\DeclareSIUnit\FCFA{FCFA} % franc CFA (ex. \SI{500}{\FCFA\per\kilo\gram})
\DeclareSIUnit\degreeBrix{\textdegree Brix} % teneur en sucre d'un sirop/jus (réfractométrie)
\DeclareSIUnit\month{mois} % le modèle utilise parfois \month, un compteur TeX, comme unité de durée
\DeclareSIUnit\microliter{\micro\liter} % le modèle écrit parfois \microliter en un seul token
\DeclareSIUnit\million{millions} % dénombrement (ex. \SI{15}{\million\per\mL} spermatozoïdes)
\usepackage{qrcode}
% \degree : commande fréquemment utilisée par les modèles pour un angle ou une
% température en dehors de \SI{}{} (non fournie par les paquets chargés).
\providecommand{\degree}{^{\circ}}
% \qed (fin de démonstration), utilisé par les modèles sans amsthm
\providecommand{\qed}{\hfill$\square$}
% \barre : double barre des valeurs interdites dans un tableau de variations (tkz-tab)
\providecommand{\barre}{\ensuremath{\|}}
% environnement proof (amsthm non chargé) utilisé par les modèles
\ifdefined\proof\else\newenvironment{proof}[1][Démonstration]{\par\noindent\textit{#1.}\ }{\qed\par}\fi
\usepackage{lastpage}
\usepackage{hyperref}
\hypersetup{hidelinks}

% ============================================================
% Palette « Pop » OnBuch+ — mêmes hex que la classe P (plus_ui.dart)
% ============================================================
\definecolor{poporange}{HTML}{F59321}
\definecolor{popdark}{HTML}{E07A0C}
\definecolor{popcream}{HTML}{FFF6E8}
\definecolor{popink}{HTML}{1C1714}
\definecolor{poppurple}{HTML}{7A5AE0}
\definecolor{popgreen}{HTML}{2E9E6A}
\definecolor{poppink}{HTML}{E0457B}
\definecolor{popblue}{HTML}{2F7DE1}
\definecolor{popgold}{HTML}{C98A12}
\definecolor{popmuted}{HTML}{8A7F76}
\definecolor{popline}{HTML}{EADFD2}
\definecolor{popgray}{HTML}{8A7F76}
\colorlet{popgrey}{popgray}
% Alias tolérants (le modèle nomme parfois les couleurs différemment)
\colorlet{popwhite}{white}
\colorlet{popblack}{popink}
\colorlet{popred}{poppink}
\colorlet{popyellow}{popgold}
% Teintes claires (fonds de tableaux/encadrés) — le modèle nomme parfois
% les couleurs avec un suffixe L (ex. popblueL) sans les avoir définies.
\colorlet{poporangeL}{poporange!20}
\colorlet{popdarkL}{popdark!20}
\colorlet{poppurpleL}{poppurple!20}
\colorlet{popgreenL}{popgreen!20}
\colorlet{poppinkL}{poppink!20}
\colorlet{popblueL}{popblue!20}
\colorlet{popgoldL}{popgold!20}
\colorlet{popredL}{popred!20}
\colorlet{popgrayL}{popgray!20}
% Teintes claires pour fonds de tableaux / zones
\colorlet{popblueL}{popblue!10!white}
\colorlet{poporangeL}{poporange!14!white}
\colorlet{popgreenL}{popgreen!12!white}
\colorlet{poppurpleL}{poppurple!12!white}
\colorlet{poppinkL}{poppink!12!white}
\colorlet{popgoldL}{popgold!14!white}

\pagecolor{popcream}

% ============================================================
% Typographie
% ============================================================
\defaultfontfeatures{Ligatures=TeX}
\setmainfont{PlusJakartaSans-Medium.ttf}[Path=../../../fonts/,BoldFont=PlusJakartaSans-Bold.ttf]
% Maths en sans-serif (Fira Math) avec chiffres et lettres droites de Plus
% Jakarta, pour que les formules se fondent dans le texte.
\usepackage[math-style=ISO,bold-style=ISO]{unicode-math}
\setmathfont{FiraMath-Regular.otf}[Path=../../../fonts/]
\setmathfont{PlusJakartaSans-Medium.ttf}[Path=../../../fonts/,range={up/num,up/{Latin,latin},\mathup}]
\usepackage{icomma} % virgule décimale sans espace en mode math
% Caractères mathématiques Unicode absents des polices de texte (Plus Jakarta,
% Archivo Black) : rendus par leur équivalent mathématique, en texte comme en
% formule — sinon ils s'affichent en carré vide (ex. « Arithmétique dans ℤ »).
\usepackage{newunicodechar}
\newunicodechar{ℝ}{\ensuremath{\mathbb{R}}}
\newunicodechar{ℤ}{\ensuremath{\mathbb{Z}}}
\newunicodechar{ℂ}{\ensuremath{\mathbb{C}}}
\newunicodechar{ℕ}{\ensuremath{\mathbb{N}}}
\newunicodechar{ℚ}{\ensuremath{\mathbb{Q}}}
\newunicodechar{𝔻}{\ensuremath{\mathbb{D}}}
\newunicodechar{α}{\ensuremath{\alpha}}
\newunicodechar{β}{\ensuremath{\beta}}
\newunicodechar{γ}{\ensuremath{\gamma}}
\newunicodechar{δ}{\ensuremath{\delta}}
\newunicodechar{ε}{\ensuremath{\varepsilon}}
\newunicodechar{θ}{\ensuremath{\theta}}
\newunicodechar{λ}{\ensuremath{\lambda}}
\newunicodechar{μ}{\ensuremath{\mu}}
\newunicodechar{σ}{\ensuremath{\sigma}}
\newunicodechar{τ}{\ensuremath{\tau}}
\newunicodechar{φ}{\ensuremath{\varphi}}
\newunicodechar{ψ}{\ensuremath{\psi}}
\newunicodechar{ω}{\ensuremath{\omega}}
\newunicodechar{Σ}{\ensuremath{\Sigma}}
% majuscules produites par \MakeUppercase dans les titres (α-aminés -> Α)
\newunicodechar{Α}{\ensuremath{\alpha}}
\newunicodechar{Β}{\ensuremath{\beta}}
\newunicodechar{Γ}{\ensuremath{\Gamma}}
\newunicodechar{Δ}{\ensuremath{\Delta}}
\newunicodechar{Ε}{\ensuremath{\varepsilon}}
\newunicodechar{Θ}{\ensuremath{\Theta}}
\newunicodechar{Λ}{\ensuremath{\Lambda}}
\newunicodechar{Μ}{\ensuremath{\mu}}
\newunicodechar{Τ}{\ensuremath{\tau}}
\newunicodechar{Φ}{\ensuremath{\Phi}}
\newunicodechar{Ψ}{\ensuremath{\Psi}}
\newunicodechar{Ω}{\ensuremath{\Omega}}
\newunicodechar{↦}{\ensuremath{\mapsto}}
\newunicodechar{∈}{\ensuremath{\in}}
\newunicodechar{∉}{\ensuremath{\notin}}
\newunicodechar{∧}{\ensuremath{\wedge}}
\newunicodechar{⇔}{\ensuremath{\Leftrightarrow}}
\newunicodechar{⟺}{\ensuremath{\Longleftrightarrow}}
\newunicodechar{⇒}{\ensuremath{\Rightarrow}}
\newunicodechar{⟹}{\ensuremath{\Longrightarrow}}
\newunicodechar{∘}{\ensuremath{\circ}}
\newunicodechar{∪}{\ensuremath{\cup}}
\newunicodechar{∩}{\ensuremath{\cap}}
\newunicodechar{≡}{\ensuremath{\equiv}}
\newunicodechar{⊂}{\ensuremath{\subset}}
\newunicodechar{⊥}{\ensuremath{\perp}}
\newunicodechar{∃}{\ensuremath{\exists}}
\newunicodechar{∀}{\ensuremath{\forall}}
\newunicodechar{∛}{\ensuremath{\sqrt[3]{\,}}}
\newunicodechar{∜}{\ensuremath{\sqrt[4]{\,}}}
\newunicodechar{⃗}{\ensuremath{{}^{\to}}}
\newunicodechar{ⁿ}{\ifmmode^{n}\else\textsuperscript{\ensuremath{n}}\fi}
\newunicodechar{ˣ}{\ifmmode^{x}\else\textsuperscript{\ensuremath{x}}\fi}
\newunicodechar{ᵇ}{\ifmmode^{b}\else\textsuperscript{\ensuremath{b}}\fi}
\newunicodechar{ʳ}{\ifmmode^{r}\else\textsuperscript{\ensuremath{r}}\fi}
\newunicodechar{ᵘ}{\ifmmode^{u}\else\textsuperscript{\ensuremath{u}}\fi}
\newunicodechar{ʸ}{\ifmmode^{y}\else\textsuperscript{\ensuremath{y}}\fi}
\newunicodechar{ᵃ}{\ifmmode^{a}\else\textsuperscript{\ensuremath{a}}\fi}
\newunicodechar{ᵉ}{\ifmmode^{e}\else\textsuperscript{\ensuremath{e}}\fi}
\newunicodechar{ᵏ}{\ifmmode^{k}\else\textsuperscript{\ensuremath{k}}\fi}
\newunicodechar{ᵖ}{\ifmmode^{p}\else\textsuperscript{\ensuremath{p}}\fi}
\newunicodechar{ᴺ}{\ifmmode^{N}\else\textsuperscript{\ensuremath{N}}\fi}
\newunicodechar{ᶜ}{\ifmmode^{c}\else\textsuperscript{\ensuremath{c}}\fi}
\newunicodechar{ʰ}{\ifmmode^{h}\else\textsuperscript{\ensuremath{h}}\fi}
\newunicodechar{ᵠ}{\ifmmode^{q}\else\textsuperscript{\ensuremath{q}}\fi}
\newunicodechar{ₙ}{\ifmmode_{n}\else\textsubscript{\ensuremath{n}}\fi}
\newunicodechar{ₐ}{\ifmmode_{a}\else\textsubscript{\ensuremath{a}}\fi}
\newunicodechar{ᵢ}{\ifmmode_{i}\else\textsubscript{\ensuremath{i}}\fi}
\newunicodechar{ᵣ}{\ifmmode_{r}\else\textsubscript{\ensuremath{r}}\fi}
\newunicodechar{ₖ}{\ifmmode_{k}\else\textsubscript{\ensuremath{k}}\fi}
\newunicodechar{ᵦ}{\ifmmode_{\beta}\else\textsubscript{\ensuremath{\beta}}\fi}
\newunicodechar{ₓ}{\ifmmode_{x}\else\textsubscript{\ensuremath{x}}\fi}
\newfontfamily\popdisplay{ArchivoBlack-Regular.ttf}[Path=../../../fonts/]
\newfontfamily\poplogo{SpaceGrotesk-Bold.ttf}[Path=../../../fonts/]
\newfontfamily\popsemi{PlusJakartaSans-SemiBold.ttf}[Path=../../../fonts/]
\newfontfamily\popxbold{PlusJakartaSans-ExtraBold.ttf}[Path=../../../fonts/]
\newfontfamily\popxboldls{PlusJakartaSans-ExtraBold.ttf}[Path=../../../fonts/,LetterSpace=3.0]

\setlist[enumerate]{leftmargin=1.7em,itemsep=5pt,topsep=4pt}
\setlist[itemize]{leftmargin=1.7em,itemsep=4pt,topsep=4pt,label={\color{poporange}\textbullet}}
\setlength{\parindent}{0pt}
\setlength{\parskip}{5pt}
\renewcommand{\arraystretch}{1.25}

% pgfplots : style Pop par défaut
\pgfplotsset{
  every axis/.append style={axis line style={popink,line width=1pt},tick style={popink},
    grid style={popline,line width=0.5pt},label style={font=\small},tick label style={font=\footnotesize}},
  popcurve/.style={poporange,line width=1.8pt,no markers,smooth},
}
% Même style disponible dans un \draw TikZ simple (hors pgfplots)
\tikzset{popcurve/.style={poporange,line width=1.8pt}}
\tikzset{popgrid/.style={popline,very thin}}
\colorlet{popcurve}{poporange} % le nom du style est parfois utilisé comme couleur

% ============================================================
% Pill / squiggle
% ============================================================
\newcommand{\poppill}[3]{%
  \tikz[baseline=-0.55ex]\node[draw=popink,line width=1.2pt,rounded corners=2.6pt,%
    fill=#2,inner xsep=6.5pt,inner ysep=3pt]{%
    \popxboldls\fontsize{7}{7}\selectfont\color{#3}\MakeUppercase{#1}};%
}
\newcommand{\popsquiggle}[1]{%
  \tikz[baseline=-0.4ex]{
    \draw[#1,line width=2.6pt,line cap=round]
      plot [smooth] coordinates {(0,0.09) (0.34,0.01) (0.68,0.21) (1.02,0.01) (1.36,0.10)};
  }%
}
% Mot-clé surligné (marqueur orange sous le texte)
\newcommand{\cle}[1]{{\popxbold\color{popdark}#1}}

% ============================================================
% En-tête / pied
% ============================================================
\pagestyle{fancy}
\fancyhf{}
\renewcommand{\headrulewidth}{0pt}
\renewcommand{\footrulewidth}{0pt}
\lhead{\raisebox{-0.15ex}{\poplogo\fontsize{13}{13}\selectfont\color{popink}OnBuch\color{poporange}+}}
\rhead{\poppill{\DOCMATIERE\ \textbullet\ \DOCNIVEAU\ \textbullet\ Leçon \DOCLECON}{white}{popink}}
\lfoot{\popxbold\fontsize{7.6}{7.6}\selectfont\color{popmuted}\MakeUppercase{Cours signé OnBuch+}\ \textbullet\ {\popsemi\DOCTITRE}}
\rfoot{\tikz[baseline=-0.6ex]\node[rounded corners=8.5pt,fill=popink,minimum height=17pt,minimum width=17pt,inner xsep=5pt,inner sep=0]{\popxbold\fontsize{8}{8}\selectfont\color{popcream}\thepage\,/\,\pageref*{LastPage}};}

% ============================================================
% Titres
% ============================================================
\newcommand{\chaptitre}[2]{%
  \begin{tcolorbox}[enhanced,colback=poporange,colframe=popink,boxrule=2.6pt,arc=11pt,
      drop shadow=popink,left=15pt,right=15pt,top=12pt,bottom=11pt,before skip=3pt,after skip=18pt]
    \poppill{#2}{popink}{popcream}\par\vspace{9pt}
    {\raggedright\hyphenpenalty=10000\exhyphenpenalty=10000\popdisplay\fontsize{22}{24}\selectfont\color{popcream}\MakeUppercase{#1}\par}
  \end{tcolorbox}
}
% Section numérotée : pastille orange avec le numéro + titre Archivo
\newcounter{popsec}
\newcommand{\coursec}[1]{%
  \stepcounter{popsec}\setcounter{popsub}{0}%
  \par\vspace{16pt}\noindent
  \begin{minipage}{\linewidth}
  \tikz[baseline=-0.7ex]\node[draw=popink,line width=1.6pt,fill=poporange,rounded corners=4pt,minimum size=22pt,inner sep=0]{\popdisplay\fontsize{12}{12}\selectfont\color{popcream}\thepopsec};%
  \hspace{8pt}{\popdisplay\fontsize{15}{17}\selectfont\color{popink}\MakeUppercase{#1}}\par
  \vspace{4pt}\noindent{\color{poporange}\rule{36pt}{3.6pt}}
  \end{minipage}\par\nopagebreak\vspace{8pt}
}
\newcounter{popsub}
\newcommand{\courssub}[1]{%
  \stepcounter{popsub}%
  \par\vspace{9pt}\noindent{\popxbold\fontsize{11.5}{13}\selectfont\color{popink}{\color{poporange}\thepopsec.\thepopsub}\hspace{6pt}#1}\par\nopagebreak\vspace{4pt}
}

% ============================================================
% Boîtes pédagogiques — même recette : fond blanc, bordure épaisse colorée,
% ombre dure encre, badge plein. Titre optionnel [#1].
% ============================================================
\tcbset{popbase/.style={breakable,enhanced,colback=white,boxrule=2.3pt,arc=9pt,
  shadow={3pt}{-3pt}{0pt}{popink},left=14pt,right=14pt,top=11pt,bottom=11pt,before skip=11pt,after skip=15pt,
  fontupper=\popsemi\fontsize{10.6}{14.5}\selectfont\color{popink},
  fontlower=\popsemi\fontsize{10.6}{14.5}\selectfont\color{popink}}}
% Coche dessinée (le glyphe ✓ n'existe pas dans les polices du document)
\newcommand{\popcheck}[1]{\tikz[baseline=-0.5ex]\draw[#1,line width=1.6pt,line cap=round,line join=round](-0.13,0.0)--(-0.04,-0.09)--(0.13,0.1);}
\providecommand{\coche}{\popcheck{popgreen}}
\providecommand{\resultat}[1]{\par\smallskip\noindent\fcolorbox{popgreen}{popgreenL}{\parbox{\dimexpr\linewidth-2\fboxsep-2\fboxrule\relax}{\textbf{Résultat.} #1}}\par\smallskip}
\newcommand{\popboxhead}[3]{\poppill{#1}{#2}{white}\ifx\relax#3\relax\else\hspace{6pt}{\popxbold\fontsize{10.6}{12}\selectfont\color{popink}#3}\fi\par\vspace{7pt}}

\newtcolorbox{aretenir}[1][]{popbase,colframe=popgreen,before upper={\popboxhead{À retenir}{popgreen}{#1}}}
\newtcolorbox{exemplebox}[1][]{popbase,colframe=poporange,before upper={\popboxhead{Exemple}{poporange}{#1}}}
\newtcolorbox{definition}[1][]{popbase,colframe=popblue,before upper={\popboxhead{Définition}{popblue}{#1}}}
\newtcolorbox{propriete}[1][]{popbase,colframe=poppurple,before upper={\popboxhead{Propriété}{poppurple}{#1}}}
\newtcolorbox{methode}[1][]{popbase,colframe=popgold,before upper={\popboxhead{Méthode}{popgold}{#1}}}
\newtcolorbox{attention}[1][]{popbase,colframe=poppink,before upper={\popboxhead{Attention}{poppink}{#1}}}
\newtcolorbox{experience}[1][]{popbase,colframe=popblue,colback=popblueL,before upper={\popboxhead{Expérience}{popblue}{#1}}}
% « Le savais-tu ? » : carte violette pleine (contexte, culture, Cameroun)
\newtcolorbox{savaistu}[1][]{popbase,colback=poppurple,colframe=popink,
  fontupper=\popsemi\fontsize{10.4}{14.2}\selectfont\color{white},
  before upper={\poppill{Le savais-tu\,?}{white}{poppurple}\ifx\relax#1\relax\else\hspace{6pt}{\popxbold\color{white}#1}\fi\par\vspace{7pt}}}
% Exercice résolu : énoncé en haut, \tcblower puis la solution sur fond crème
\newcounter{popexo}\newcounter{popexr}
\newtcolorbox{exoresolu}[1][]{popbase,colframe=poporange,colbacklower=popcream,
  segmentation style={popink,line width=1.2pt,dashed},
  before upper={\stepcounter{popexr}\popboxhead{Exercice résolu \thepopexr}{poporange}{#1}},
  before lower={\poppill{Solution}{popink}{popcream}\par\vspace{6pt}}}
% Exercice (non résolu en place) — la correction est dans la partie Corrigés
\newtcolorbox{exercice}[1][]{popbase,colframe=popink,
  before upper={\stepcounter{popexo}\popboxhead{Exercice \thepopexo}{popink}{#1}}}
% Corrigé
\newtcolorbox{corrige}[1][]{popbase,colframe=popgreen,colback=popgreenL,
  before upper={\popboxhead{Corrigé}{popgreen}{#1}}}
% Objectifs / prérequis / bilan
\newtcolorbox{objectifs}{popbase,colframe=popink,colback=poporangeL,
  before upper={\popboxhead{Objectifs de la leçon}{popink}{}}}
\newtcolorbox{prerequis}{popbase,colframe=popink,colback=popblueL,
  before upper={\popboxhead{Prérequis}{popblue}{}}}
\newtcolorbox{bilan}{popbase,colframe=popink,colback=white,boxrule=2.8pt,
  before upper={\popboxhead{Fiche bilan}{poporange}{L'essentiel à connaître pour l'examen}}}
% Figure encadrée avec légende
\newtcolorbox{popfigure}[1][]{enhanced,breakable=false,colback=white,colframe=popline,boxrule=1.4pt,arc=8pt,
  left=10pt,right=10pt,top=10pt,bottom=8pt,before skip=10pt,after skip=12pt,halign=center,
  lower separated=false,halign lower=center,
  fontlower=\popsemi\fontsize{9}{11.5}\selectfont\color{popmuted},
  #1}
\newcommand{\legende}[1]{\par\vspace{6pt}{\popsemi\fontsize{9}{11.5}\selectfont\color{popmuted}#1}}

% ============================================================
% Signature OnBuch+
% ============================================================
\newcommand{\poplogoinline}[1]{{\poplogo\fontsize{#1}{#1}\selectfont\color{popink}OnBuch\color{poporange}+}}
\newcommand{\signatureonbuch}{%
  \begin{tcolorbox}[enhanced,colback=popink,colframe=popink,boxrule=2.2pt,arc=10pt,
      drop shadow=poporange,left=14pt,right=14pt,top=10pt,bottom=10pt,before skip=0pt,after skip=0pt]
    \tikz[baseline=-0.7ex]\node[circle,fill=popgreen,draw=popcream,line width=1.3pt,minimum size=22pt,inner sep=0]{\popcheck{white}};%
    \hspace{9pt}\begin{minipage}[c]{0.62\linewidth}
      {\popxbold\fontsize{12}{13}\selectfont\color{popcream}Signé\ \textbullet\ {\poplogo OnBuch\color{poporange}+}}\par\vspace{2pt}
      {\raggedright\popsemi\fontsize{8}{10.5}\selectfont\color{popline}\MakeUppercase{Équipe pédagogique OnBuch+}\\\MakeUppercase{Conforme au programme officiel MINESEC}\par}
    \end{minipage}\hfill
    {\poplogo\fontsize{22}{22}\selectfont\color{popcream}OnBuch\color{poporange}+}
  \end{tcolorbox}%
}

% ============================================================
% Page de garde
% #1 titre · #2 matière · #3 niveau · #4 module · #5 n° leçon · #6 durée ·
% #7 description / objectifs (texte ou itemize)
% ============================================================
\newcommand{\pagegarde}[7]{%
  \thispagestyle{empty}
  \newgeometry{a4paper,margin=1.7cm,top=1.6cm,bottom=1.4cm}%
  \begin{tikzpicture}[remember picture,overlay]
    \fill[poporange!18] ([xshift=-3.2cm,yshift=-3.6cm]current page.north east) circle (4.6cm);
    \fill[poppurple!14] ([xshift=2.4cm,yshift=7.6cm]current page.south west) circle (3.8cm);
  \end{tikzpicture}%
  {\poplogo\fontsize{30}{30}\selectfont\color{popink}OnBuch\color{poporange}+}\hfill
  \raisebox{0.6ex}{\poppill{Cours signé \textbullet\ Premium}{popink}{popcream}}\par
  \vspace{1pt}\popsquiggle{poppurple}\par
  \vspace{8pt}
  \begin{tcolorbox}[enhanced,colback=poporange,colframe=popink,boxrule=2.8pt,arc=13pt,
      drop shadow=popink,left=18pt,right=18pt,top=12pt,bottom=13pt,after skip=10pt]
    \poppill{#2 \textbullet\ #3}{popink}{popcream}\hspace{5pt}\poppill{#4}{white}{popink}\par\vspace{8pt}
    {\popdisplay\fontsize{14}{14}\selectfont\color{popink}LEÇON #5}\par\vspace{4pt}
    {\raggedright\hyphenpenalty=10000\exhyphenpenalty=10000\popdisplay\fontsize{25}{27}\selectfont\color{popcream}\MakeUppercase{#1}\par}
    \vspace{8pt}
    {\popxbold\fontsize{9.5}{9.5}\selectfont\color{popink}\MakeUppercase{Durée conseillée : #6}}
  \end{tcolorbox}
  \begin{tcolorbox}[enhanced,colback=white,colframe=popink,boxrule=2.2pt,arc=10pt,
      drop shadow=popink,left=16pt,right=16pt,top=10pt,bottom=10pt,after skip=8pt]
    \poppill{Dans cette leçon}{poppurple}{white}\par\vspace{5pt}
    \popsemi\fontsize{9.3}{12.6}\selectfont\color{popink}#7
  \end{tcolorbox}
  \noindent\begin{minipage}[t]{0.487\linewidth}
    \begin{tcolorbox}[enhanced,colback=popgreen,colframe=popink,boxrule=2.2pt,arc=10pt,
        drop shadow=popink,left=11pt,right=11pt,top=10pt,bottom=10pt,height=3.1cm,valign=center]
      \tcbox[colback=white,colframe=popink,boxrule=1.3pt,arc=5pt,boxsep=0pt,nobeforeafter,
        left=3pt,right=3pt,top=3pt,bottom=3pt,valign=center]{\qrcode[height=1.9cm]{https://play.google.com/store/apps/details?id=com.luvvix.onbuch&hl=fr}}%
      \hspace{8pt}\begin{minipage}[c]{0.5\linewidth}
        {\popdisplay\fontsize{11}{12.5}\selectfont\color{white}\MakeUppercase{Télécharge OnBuch}}\par\vspace{4pt}
        {\raggedright\popsemi\fontsize{8.4}{11}\selectfont\color{white}Gratuit \textbullet\ Google Play\\Tuteur IA, annales, concours, résultats\par}
      \end{minipage}
    \end{tcolorbox}
  \end{minipage}\hfill
  \begin{minipage}[t]{0.487\linewidth}
    \begin{tcolorbox}[enhanced,colback=poppurple,colframe=popink,boxrule=2.2pt,arc=10pt,
        drop shadow=popink,left=11pt,right=11pt,top=10pt,bottom=10pt,height=3.1cm,valign=center]
      \tikz[baseline=-0.7ex]\node[circle,fill=white,draw=popink,line width=1.3pt,minimum size=1.6cm,inner sep=0]{\popdisplay\fontsize{24}{24}\selectfont\color{poppurple}+};%
      \hspace{8pt}\begin{minipage}[c]{0.6\linewidth}
        {\popdisplay\fontsize{11}{12.5}\selectfont\color{white}\MakeUppercase{Passe à OnBuch+}}\par\vspace{4pt}
        {\raggedright\popsemi\fontsize{8.4}{11}\selectfont\color{white}Tous les cours signés, Tuteur IA illimité, fiches PDF — onglet OnBuch+ de l'app\par}
      \end{minipage}
    \end{tcolorbox}
  \end{minipage}
  \vspace{8pt}
  \popcontenu
  \vfill
  \signatureonbuch
  \restoregeometry
  \newpage
}

% Rangée « Ce cours contient » (page de garde)
\newcommand{\poptile}[3]{% #1 couleur #2 gros texte #3 légende
  \begin{tcolorbox}[enhanced,colback=#1,colframe=popink,boxrule=2pt,arc=9pt,shadow={2.5pt}{-2.5pt}{0pt}{popink},
      left=6pt,right=6pt,top=9pt,bottom=9pt,halign=center,valign=center,height=1.75cm,width=\linewidth,nobeforeafter]
    {\popdisplay\fontsize{12.5}{13}\selectfont\color{white}\MakeUppercase{#2}}\par\vspace{3pt}
    {\centering\popsemi\fontsize{7.8}{9.5}\selectfont\color{white}#3\par}
  \end{tcolorbox}}
\newcommand{\popcontenu}{%
  \par\noindent{\popxboldls\fontsize{7.5}{7.5}\selectfont\color{popmuted}CE COURS SIGNÉ CONTIENT}\par\vspace{7pt}
  \noindent\begin{minipage}[t]{0.235\linewidth}\poptile{popblue}{Cours}{complet, illustré, pas à pas}\end{minipage}\hfill
  \begin{minipage}[t]{0.235\linewidth}\poptile{poporange}{Méthodes}{et exercices résolus}\end{minipage}\hfill
  \begin{minipage}[t]{0.235\linewidth}\poptile{poppink}{Exercices}{type examen + corrigés}\end{minipage}\hfill
  \begin{minipage}[t]{0.235\linewidth}\poptile{popgreen}{Fiche bilan}{l'essentiel à retenir}\end{minipage}\par}

% Sommaire
\newtcolorbox{sommaire}{enhanced,colback=white,colframe=popink,boxrule=2.2pt,arc=10pt,
  drop shadow=popink,left=18pt,right=18pt,top=13pt,bottom=13pt,before skip=6pt,after skip=20pt,
  before upper={\poppill{Sommaire}{poppurple}{white}\par\vspace{9pt}\popsemi\fontsize{10.6}{14.5}\selectfont\color{popink}}}

% Signature de fin de cours — poussée tout en bas de la dernière page
\newcommand{\findecours}{%
  \par\vfill\nopagebreak
  \begin{center}\popsquiggle{poporange}\end{center}\vspace{4pt}
  \signatureonbuch
}
\AtBeginDocument{\def\llbracket{\mathopen{[\mkern-3.5mu[}}\def\rrbracket{\mathclose{]\mkern-3.5mu]}}} % intervalles d'entiers (glyphe ⟦ absent de la police)
% Glyphes absents de Fira Math : redéfinis à partir de glyphes disponibles
\AtBeginDocument{%
  \def\setminus{\mathbin{\backslash}}\let\smallsetminus\setminus
  \def\vdots{\mathord{\vbox{\baselineskip3.2pt\lineskiplimit0pt\kern4pt\hbox{.}\hbox{.}\hbox{.}}}}%
  \def\ddots{\mathinner{\mkern1mu\raise6.5pt\hbox{.}\mkern2mu\raise3.6pt\hbox{.}\mkern2mu\raise0.7pt\hbox{.}\mkern1mu}}%
  \def\triangle{\mathord{\scalebox{0.9}{$\symup{\Delta}$}}}%
  \def\nabla{\mathord{\raisebox{\depth}{\scalebox{1}[-1]{$\symup{\Delta}$}}}}%
  \def\checkmark{\popcheck{popdark}}%
  \def\star{\mathbin{\ast}}\def\bigstar{\mathord{\ast}}%
  \def\diamond{\mathbin{\scalebox{0.6}{$\Diamond$}}}%
  \def\bigcup{\mathop{\mathchoice{\raisebox{-0.3em}{\scalebox{1.7}{$\cup$}}}{\scalebox{1.25}{$\cup$}}{\cup}{\cup}}\displaylimits}%
  \def\bigcap{\mathop{\mathchoice{\raisebox{-0.3em}{\scalebox{1.7}{$\cap$}}}{\scalebox{1.25}{$\cap$}}{\cap}{\cap}}\displaylimits}%
  \def\lesseqqgtr{\mathrel{\lessgtr}}\def\bigoplus{\mathop{\oplus}\displaylimits}%
}
\providecommand{\ssi}{si et seulement si} % abréviation fréquente
\AtBeginDocument{\providecommand{\wideparen}{\overparen}} % arc de cercle

\begin{document}

\pagegarde{Organisation de la cellule et échanges cellulaires}{SVTEEHB}{Tle C}{Module I : Le Monde Vivant}{N01}{5 h}{Cette leçon présente l'organisation générale de la cellule animale observée en microscopie optique et électronique, ainsi que les rôles de ses principaux organites. Elle établit ensuite, à partir d'expériences historiques (Dutrochet, Pfeiffer) et de résultats expérimentaux, les mécanismes des échanges cellulaires : osmose, diffusion, transport actif, endocytose et exocytose.
\vspace{4pt}
\begin{multicols}{2}\small
\begin{itemize}
\item À la fin de cette leçon, je sais observer et schématiser une cellule animale en microscopie optique et en microscopie électronique
\item À la fin de cette leçon, je sais schématiser et annoter les principaux organites cellulaires et associer chacun à son rôle
\item À la fin de cette leçon, je sais définir l'osmose et la dialyse et les distinguer
\item À la fin de cette leçon, je sais interpréter les résultats de l'expérience de Dutrochet et de l'expérience de Pfeiffer
\item À la fin de cette leçon, je sais analyser, interpréter et conclure à partir de résultats d'expériences sur les échanges d'eau (hématies en milieux hypotonique, isotonique, hypertonique)
\item À la fin de cette leçon, je sais distinguer le transport passif (diffusion simple, diffusion facilitée) du transport actif à partir de courbes de cinétique d'absorption
\end{itemize}\end{multicols}}

\begin{sommaire}
\begin{multicols}{2}
\begin{enumerate}
\item Observation et organisation générale de la cellule animale
\item Les principaux organites cellulaires et leurs rôles
\item Les échanges d'eau : osmose et dialyse
\item Les échanges de substances dissoutes : transports passif et actif
\item Les échanges de particules : endocytose et exocytose
\item Importance des échanges cellulaires dans la vie des animaux
\item Activité d'intégration
\item Méthodes et exercices résolus
\item Exercices
\item Corrigés des exercices
\item Fiche bilan
\end{enumerate}
\end{multicols}
\end{sommaire}

\begin{objectifs}
\begin{itemize}
\item À la fin de cette leçon, je sais observer et schématiser une cellule animale en microscopie optique et en microscopie électronique
\item À la fin de cette leçon, je sais schématiser et annoter les principaux organites cellulaires et associer chacun à son rôle
\item À la fin de cette leçon, je sais définir l'osmose et la dialyse et les distinguer
\item À la fin de cette leçon, je sais interpréter les résultats de l'expérience de Dutrochet et de l'expérience de Pfeiffer
\item À la fin de cette leçon, je sais analyser, interpréter et conclure à partir de résultats d'expériences sur les échanges d'eau (hématies en milieux hypotonique, isotonique, hypertonique)
\item À la fin de cette leçon, je sais distinguer le transport passif (diffusion simple, diffusion facilitée) du transport actif à partir de courbes de cinétique d'absorption
\item À la fin de cette leçon, je sais décrire et schématiser l'endocytose et l'exocytose
\item À la fin de cette leçon, je sais expliquer l'importance des échanges cellulaires dans la vie des animaux (nutrition, respiration cellulaire, excrétion, immunité)
\end{itemize}
\end{objectifs}

\begin{prerequis}
\begin{itemize}
\item Notion de cellule comme unité structurale du vivant (classes de Seconde et Première)
\item Notion de solution, de concentration et de solvant (physique-chimie)
\item Notion de membrane et de milieu intérieur (début d'année de Terminale)
\item Lecture et exploitation de courbes et de tableaux de résultats expérimentaux
\item Notion de molécules organiques : glucides, lipides, protéines (Première)
\end{itemize}
\end{prerequis}

\newpage

\coursec{Observation et organisation générale de la cellule animale}

\textbf{Situation d'accroche.} Sur le marché de Douala, Aminatou vend du poisson fumé. Chaque matin, elle constate le même phénomène : les poissons généreusement salés rétrécissent, durcissent et perdent leur eau, tandis que les poissons frais laissés dans une bassine d'eau douce gonflent et deviennent fermes. Quelques kilomètres plus loin, à l'hôpital Laquintinie, l'infirmier administre aux patients déshydratés un sérum dit « physiologique » (solution de chlorure de sodium à \SI{9}{g/L}) et jamais de l'eau pure injectée directement dans le sang. Pourquoi cette précaution ? Que se passerait-il si l'on injectait de l'eau pure à un patient ?

Ces phénomènes quotidiens, en apparence sans rapport, ont une même explication : les \cle{échanges} qui se font à travers la \cle{membrane} des cellules. Mais avant de comprendre comment la cellule échange, il faut d'abord savoir comment elle est \cle{organisée}. \textbf{Problématique :} comment la cellule animale est-elle organisée, et que révèlent les différents outils d'observation dont dispose le biologiste ? C'est la question à laquelle cette première section va répondre.

\courssub{La cellule : unité structurale et fonctionnelle du vivant}

\textbf{Observation.} Prélève un peu de cellules de ta muqueuse buccale en grattant délicatement l'intérieur de ta joue avec un objet propre, étale le prélèvement sur une lame, colore-le au bleu de méthylène et observe au microscope optique : tu distingues de petites entités aplaties, ce sont des cellules épithéliales. Observe maintenant un frottis de sang : tu vois des millions de petits disques roses, les hématies. Tous les êtres vivants, du plus petit au plus grand, sont constitués de telles unités.

\begin{definition}[La cellule]
La \cle{cellule} est l'\textbf{unité structurale et fonctionnelle} de tous les êtres vivants : tout organisme est constitué d'une ou plusieurs cellules, et toute fonction vitale (nutrition, respiration, reproduction) s'accomplit au niveau cellulaire. Toute cellule provient d'une cellule préexistante, par division cellulaire.
\end{definition}

\textbf{La diversité des formes cellulaires animales.} Il n'existe pas une cellule animale « type » dans la réalité : les cellules animales présentent une grande diversité de formes et de tailles, \emph{adaptées à leur fonction}. C'est la notion de \cle{spécialisation cellulaire} :

\begin{itemize}
\item l'\textbf{hématie} (globule rouge) : petit disque biconcave d'environ \SI{7,5}{\micro m} de diamètre, dépourvue de noyau à l'état mature ; sa forme maximise la surface d'échange pour le transport du dioxygène par l'hémoglobine ;
\item le \textbf{neurone} : cellule étoilée portant un très long prolongement, l'axone (jusqu'à \SI{1}{m} chez l'Homme !), adapté à la conduction de l'influx nerveux sur de longues distances ;
\item le \textbf{spermatozoïde} : cellule d'environ \SI{50}{\micro m} de longueur totale, dotée d'un flagelle qui lui permet de se déplacer jusqu'à l'ovule ;
\item la \textbf{cellule épithéliale} : cellule aplatie ou pavimenteuse, jointive à ses voisines, formant des tissus de revêtement (peau, muqueuses).
\end{itemize}

\begin{aretenir}[Taille des cellules]
La plupart des cellules animales mesurent entre \SI{10}{\micro m} et \SI{100}{\micro m}. Elles sont donc \textbf{invisibles à l'œil nu} (dont le pouvoir de résolution est d'environ \SI{0,1}{mm}, soit \SI{100}{\micro m}) : leur observation exige des instruments d'optique.
\end{aretenir}

\courssub{L'observation de la cellule animale : microscopie optique et microscopie électronique}

\textbf{Le problème de la résolution.} Un instrument d'observation est caractérisé par son \cle{pouvoir de résolution} : la plus petite distance séparant deux points que l'instrument permet de distinguer comme distincts. En dessous de cette limite, deux points proches apparaissent confondus. C'est ce pouvoir de résolution — et non le grossissement — qui détermine ce que l'on peut voir.

\begin{propriete}[Pouvoirs de résolution]
\begin{itemize}
\item Œil humain : environ \SI{0,1}{mm} = \SI{100}{\micro m} ;
\item \textbf{Microscope optique (MO)} : environ \SI{0,2}{\micro m} (limite imposée par la longueur d'onde de la lumière visible) ; grossissement utile maximal : environ $\times 1000$ ;
\item \textbf{Microscope électronique (ME)} : environ \SI{0,2}{nm} (le faisceau d'électrons est associé à une longueur d'onde bien plus faible que celle de la lumière) ; grossissements supérieurs à $\times 100\,000$.
\end{itemize}
\end{propriete}

\textbf{Ce que révèle le microscope optique.} Sur une cellule animale colorée (par exemple au bleu de méthylène ou au carmin), le MO permet de distinguer \textbf{trois structures} :

\begin{enumerate}
\item la \textbf{membrane plasmique}, fine bordure qui délimite la cellule ;
\item le \textbf{cytoplasme}, contenu cellulaire d'aspect granuleux ;
\item le \textbf{noyau}, structure sphérique plus dense, souvent colorée plus intensément.
\end{enumerate}

\textbf{Ce que révèle le microscope électronique.} Le ME, grâce à son pouvoir de résolution environ mille fois supérieur, lève le voile sur l'\cle{ultrastructure} cellulaire : le cytoplasme n'est pas homogène, il contient de nombreux \textbf{organites} (mitochondries, réticulum endoplasmique, appareil de Golgi, ribosomes, lysosomes, centrioles...), et le noyau apparaît entouré d'une \textbf{double membrane nucléaire} percée de pores. La membrane plasmique elle-même, d'une épaisseur d'environ \SI{7,5}{nm}, n'est résolue qu'au ME.

\begin{attention}[Ne pas confondre grossissement et résolution]
Grossir une image floue ne la rend pas plus nette ! Au-delà du grossissement utile ($\times 1000$ en MO), l'image s'agrandit mais ne révèle aucun détail supplémentaire : c'est le \textbf{pouvoir de résolution} qui compte. C'est pourquoi les organites comme les mitochondries, bien que l'image puisse être agrandie à l'infini en MO, n'y sont pas résolus dans leurs détails internes.
\end{attention}

\begin{popfigure}
\begin{center}
\begin{tikzpicture}[scale=0.92, every node/.style={font=\small}]
% ===== Cellule MO =====
\begin{scope}
\fill[popblueL] (0,0) ellipse (2.6 and 1.9);
\draw[popblue, thick] (0,0) ellipse (2.6 and 1.9);
\fill[popgoldL] (0.2,0) circle (0.75);
\draw[popgold, thick] (0.2,0) circle (0.75);
\fill[popgold] (0.2,0) circle (0.18);
\draw[popdark] (-2.6,1.6) node[anchor=south west, font=\bfseries] {Vue en MO ($\times 1000$)};
% annotations MO
\draw[thin, popdark] (2.2,1.15) -- (3.4,1.6) node[anchor=west] {1. Membrane plasmique};
\draw[thin, popdark] (-1.3,-1.2) -- (-1.3,-2.5) node[anchor=north, text width=3cm, align=center] {2. Cytoplasme\\ (aspect granuleux)};
\draw[thin, popdark] (0.2,-0.75) -- (1.6,-1.9) node[anchor=north west] {3. Noyau};
\end{scope}
% ===== Cellule ME =====
\begin{scope}[xshift=9.5cm]
\fill[popcream] (0,0) ellipse (2.9 and 2.1);
\draw[poporange, thick] (0,0) ellipse (2.9 and 2.1);
\draw[poporange] (0,0) ellipse (2.82 and 2.02);
% noyau double membrane
\fill[popgoldL] (-0.7,0.3) circle (0.85);
\draw[popgold, thick] (-0.7,0.3) circle (0.85);
\draw[popgold] (-0.7,0.3) circle (0.78);
\fill[popgold] (-0.9,0.4) circle (0.2);
% mitochondries
\begin{scope}[shift={(1.3,0.9)}, rotate=20]
    \draw[popgreen, thick, fill=popgreenL] (0,0) ellipse (0.42 and 0.18);
    \draw[popgreen] (-0.3,0) .. controls (-0.2,0.1) and (-0.1,-0.1) .. (0,0)
      .. controls (0.1,0.1) and (0.2,-0.1) .. (0.3,0);
\end{scope}
\begin{scope}[shift={(1.7,-0.7)}, rotate=-30]
    \draw[popgreen, thick, fill=popgreenL] (0,0) ellipse (0.42 and 0.18);
    \draw[popgreen] (-0.3,0) .. controls (-0.2,0.1) and (-0.1,-0.1) .. (0,0)
      .. controls (0.1,0.1) and (0.2,-0.1) .. (0.3,0);
\end{scope}
\begin{scope}[shift={(0.3,-1.2)}, rotate=10]
    \draw[popgreen, thick, fill=popgreenL] (0,0) ellipse (0.42 and 0.18);
    \draw[popgreen] (-0.3,0) .. controls (-0.2,0.1) and (-0.1,-0.1) .. (0,0)
      .. controls (0.1,0.1) and (0.2,-0.1) .. (0.3,0);
\end{scope}
\begin{scope}[shift={(-1.9,-0.9)}, rotate=45]
    \draw[popgreen, thick, fill=popgreenL] (0,0) ellipse (0.42 and 0.18);
    \draw[popgreen] (-0.3,0) .. controls (-0.2,0.1) and (-0.1,-0.1) .. (0,0)
      .. controls (0.1,0.1) and (0.2,-0.1) .. (0.3,0);
\end{scope}
% réticulum
\draw[poppurple, thick] (0.4,1.1) .. controls (0.9,1.4) and (1.6,1.2) .. (2.1,1.5);
\draw[poppurple, thick] (0.4,0.95) .. controls (0.9,1.25) and (1.6,1.05) .. (2.05,1.35);
\draw[poppurple, thick] (0.35,0.8) .. controls (0.9,1.1) and (1.6,0.9) .. (2.0,1.2);
% Golgi
\foreach \y in {-1.55,-1.42,-1.29}{\draw[poppink, thick] (-0.4,\y) .. controls (0.1,\y+0.08) and (0.5,\y-0.08) .. (0.9,\y);}
% ribosomes
\foreach \p in {(0.6,1.35),(1.1,1.15),(1.7,1.3),(2.0,1.42)}{\fill[popdark] \p circle (0.045);}
% lysosome
\fill[popblue!40] (-2.1,0.9) circle (0.16); \draw[popblue] (-2.1,0.9) circle (0.16);
\draw[popdark] (-2.9,1.7) node[anchor=south west, font=\bfseries] {Vue en ME ($\times 20\,000$)};
% annotations ME
\draw[thin, popdark] (2.55,1.35) -- (3.5,2.0) node[anchor=west, text width=3.6cm] {4. Membrane plasmique (double couche)};
\draw[thin, popdark] (1.3,0.9) -- (3.2,0.6) node[anchor=west, text width=3.4cm] {5. Mitochondrie};
\draw[thin, popdark] (2.1,1.5) -- (3.3,1.2) node[anchor=west, text width=3.2cm] {6. Réticulum endoplasmique granuleux};
\draw[thin, popdark] (0.9,-1.42) -- (3.3,-1.5) node[anchor=west, text width=3.2cm] {7. Appareil de Golgi};
\draw[thin, popdark] (-0.7,-0.55) -- (-0.7,-2.6) node[anchor=north, text width=3.4cm, align=center] {8. Noyau à double membrane (nucléole visible)};
\draw[thin, popdark] (-2.1,0.9) -- (-3.4,1.3) node[anchor=east, text width=2.8cm, align=right] {9. Lysosome};
\end{scope}
\end{tikzpicture}
\end{center}
\legende{Comparaison d'une cellule animale observée en microscopie optique (à gauche : seules trois structures sont distinguées) et en microscopie électronique (à droite : l'ultrastructure des organites est révélée). Schéma d'interprétation, sans échelle.}
\end{popfigure}

\courssub{L'échelle des tailles : de la cellule à la molécule}

Pour situer les objets biologiques et les limites des instruments, il est commode d'utiliser une \textbf{échelle logarithmique} : chaque graduation correspond à une puissance de dix. Rappelons les unités : $\SI{1}{mm} = \SI{1000}{\micro m}$ ; $\SI{1}{\micro m} = \SI{1000}{nm}$.

\begin{popfigure}
\begin{center}
\begin{tikzpicture}[scale=1.0, every node/.style={font=\small}]
\draw[->, thick, popdark] (0,0) -- (13.2,0) node[anchor=west] {taille (échelle log.)};
% graduations : 1 nm -> 1 mm
\foreach \i/\lab in {0/{1 nm},1/{10 nm},2/{100 nm},3/{1 \textmu m},4/{10 \textmu m},5/{100 \textmu m},6/{1 mm}}{
  \draw[thick, popdark] ({2*\i},-0.12) -- ({2*\i},0.12);
  \node[below] at ({2*\i},-0.15) {\lab};
}
% objets
\node[above, popblue, align=center] at (6.4,0.15) {molécule d'hémoglobine};
\draw[popblue, thick] (6.4,0.1) -- (6.4,0.7);
\node[above, popgreen, align=center] at (5.0,0.85) {membrane plasmique (7,5 nm)};
\draw[popgreen, thick] (5.75,0.1) -- (5.0,0.8);
\node[above, poporange, align=center] at (7.6,0.85) {mitochondrie (1 \textmu m)};
\draw[poporange, thick] (6.0,0.1) -- (7.4,0.8);
\node[above, poppurple, align=center] at (9.6,0.15) {hématie (7,5 \textmu m)};
\draw[poppurple, thick] (7.75,0.1) -- (9.6,0.7);
\node[above, poppink, align=center] at (11.2,0.85) {ovule humain (100 \textmu m)};
\draw[poppink, thick] (10.0,0.1) -- (11.0,0.8);
% limites instruments (flèches vers le bas)
\draw[<-, very thick, popgold] (5.4,-0.9) -- (5.4,-0.25);
\node[below, popgold, align=center] at (5.4,-0.95) {limite du MO (0,2 \textmu m)};
\draw[<-, very thick, popgold] (1.4,-0.9) -- (1.4,-0.25);
\node[below, popgold, align=center] at (1.4,-0.95) {limite du ME (0,2 nm)};
\draw[<-, very thick, popgold] (10.0,-0.9) -- (10.0,-0.25);
\node[below, popgold, align=center] at (10.0,-0.95) {limite de l'œil (0,1 mm)};
\end{tikzpicture}
\end{center}
\legende{Échelle logarithmique des tailles biologiques et pouvoirs de résolution de l'œil nu, du microscope optique (MO) et du microscope électronique (ME). Tout objet situé à gauche de la limite d'un instrument lui est invisible.}
\end{popfigure}

\begin{exoresolu}[Convertir et situer sur l'échelle des tailles]
Énoncé. Une mitochondrie mesure environ \SI{1}{\micro m} de long. (a) Exprime cette taille en nanomètres. (b) Cette mitochondrie est-elle visible à l'œil nu ? Au microscope optique ? Justifie. (c) Combien de mitochondries mises bout à bout faudrait-il pour égaler le diamètre d'une hématie de \SI{7,5}{\micro m} ?
\tcblower
Solution. (a) $\SI{1}{\micro m} = \SI{1000}{nm}$, donc la mitochondrie mesure \SI{1000}{nm}. (b) À l'œil nu : non, car $\SI{1}{\micro m} < \SI{100}{\micro m}$ (limite de l'œil). Au MO : oui, car $\SI{1}{\micro m} > \SI{0,2}{\micro m}$ (limite de résolution du MO) ; la mitochondrie y apparaît comme un petit bâtonnet, mais ses crêtes internes (espacées de quelques dizaines de nanomètres) ne sont résolues qu'au ME. (c) Il faut $7{,}5 / 1 = 7{,}5$, soit environ \textbf{7 à 8 mitochondries} bout à bout.
\end{exoresolu}

\courssub{L'organisation générale de la cellule animale}

L'ensemble des observations en MO et en ME permet de dégager le plan d'organisation commun à toutes les cellules animales. Trois compartiments fondamentaux sont toujours présents :

\begin{definition}[Organisation générale de la cellule animale]
Toute cellule animale comprend :
\begin{enumerate}
\item une \textbf{membrane plasmique} (ou plasmalemme) : enveloppe fine (\SI{7,5}{nm}) qui délimite la cellule, contrôle les échanges avec le milieu extérieur et reçoit des signaux ;
\item un \textbf{cytoplasme} : constitué de l'\cle{hyaloplasme} (ou cytosol), phase aqueuse fondamentale dans laquelle baignent les \textbf{organites} (mitochondries, réticulum endoplasmique, appareil de Golgi, lysosomes, ribosomes, centrosome...) et le cytosquelette ;
\item un \textbf{noyau} : entouré d'une double membrane nucléaire, il contient l'ADN (matériel génétique) organisé en chromosomes, et un ou plusieurs nucléoles. C'est le centre de commandement de la cellule.
\end{enumerate}
\end{definition}

\begin{definition}[Cellule eucaryote]
Une cellule \cle{eucaryote} est une cellule possédant un \textbf{noyau individualisé} par une membrane nucléaire et des \textbf{organites membranaires}. La cellule animale est une cellule eucaryote \textbf{dépourvue de paroi cellulaire et de chloroplastes}.
\end{definition}

\begin{attention}[Erreur fréquente : membrane plasmique et paroi cellulaire]
Ne confonds jamais \textbf{membrane plasmique} et \textbf{paroi cellulaire} ! La membrane plasmique est une enveloppe fine, souple et vivante, présente chez \emph{toutes} les cellules. La paroi (cellulosique chez les végétaux) est une enveloppe rigide \emph{supplémentaire}, externe à la membrane. La \textbf{cellule animale n'a pas de paroi} : c'est pourquoi elle est déformable et fragile face aux variations de pression osmotique — point crucial pour la suite de la leçon (osmose).
\end{attention}

\begin{center}
\begin{tabularx}{\linewidth}{|l|X|X|}\hline
\rowcolor{popblueL} \textbf{Critère} & \textbf{Cellule animale} & \textbf{Cellule végétale chlorophyllienne} \\ \hline
Noyau individualisé & Oui (eucaryote) & Oui (eucaryote) \\ \hline
Membrane plasmique & Oui & Oui \\ \hline
Paroi cellulaire & \textbf{Non} & Oui (cellulose) \\ \hline
Chloroplastes & \textbf{Non} & Oui (photosynthèse) \\ \hline
Vacuole & Absente ou petites vacuoles & Grande vacuole centrale \\ \hline
Centrioles & Oui & Généralement absents \\ \hline
Forme & Variable, déformable & Régulière, rigide \\ \hline
\end{tabularx}
\end{center}

\courssub{Savoir-faire : schématiser correctement une cellule animale}

Au Baccalauréat, la qualité du schéma est notée. Un schéma biologique n'est pas un dessin artistique : c'est un document scientifique codifié.

\begin{methode}[Réaliser un schéma de cellule animale]
\begin{enumerate}
\item Dessine au crayon à papier, avec des \textbf{traits continus et fins}, sans hachures ni coloriage appuyé.
\item Représente uniquement les structures \textbf{réellement observées} (ou demandées), en respectant leurs proportions relatives (le noyau occupe environ 1/10 du volume cellulaire).
\item Écris les \textbf{annotations à l'horizontale}, à l'extérieur du schéma, en les reliant aux structures par des \textbf{traits fins} qui ne se croisent pas et se terminent par un point ou une flèche sur la structure.
\item Ajoute obligatoirement un \textbf{titre} sous le schéma et l'\textbf{échelle} (ou le grossissement).
\end{enumerate}
\end{methode}

\begin{attention}[Pièges de rédaction au Bac]
\begin{itemize}
\item Ne jamais écrire « paroi » pour une cellule animale : la bordure externe est la \emph{membrane plasmique}.
\item Ne pas dessiner de chloroplastes ni de grande vacuole dans une cellule animale.
\item Ne pas confondre « cytoplasme » (compartiment entier) et « hyaloplasme » (phase aqueuse du cytoplasme, hors organites).
\item Un schéma sans titre ni annotations horizontales perd des points, même s'il est exact.
\end{itemize}
\end{attention}

\begin{savaistu}[Un globule rouge sans noyau]
Chez l'Homme, l'hématie (globule rouge) \textbf{mature perd son noyau} au cours de sa différenciation dans la moelle osseuse : c'est une exception à la règle « toute cellule possède un noyau ». Cette perte libère de la place pour stocker davantage d'\textbf{hémoglobine} (environ 280 millions de molécules par hématie !), optimisant le transport du dioxygène. Contrepartie : incapable de synthétiser des protéines ni de se diviser, l'hématie ne vit que 120 jours environ et doit être renouvelée en permanence — environ 2 millions d'hématies produites chaque seconde dans ta moelle osseuse. C'est aussi pourquoi le paludisme, dont le parasite (\emph{Plasmodium}) se multiplie dans les hématies, est si redoutable au Cameroun.
\end{savaistu}

\begin{exercice}[Pour vérifier tes acquis]
\begin{enumerate}
\item Définis la cellule eucaryote et cite deux caractères qui distinguent la cellule animale de la cellule végétale.
\item Un organite mesure \SI{150}{nm}. Peut-on l'observer au microscope optique ? Justifie par une comparaison chiffrée.
\item Réalise le schéma annoté d'une cellule animale telle qu'elle apparaît en microscopie optique, en respectant les règles du schéma biologique.
\end{enumerate}
\end{exercice}

\begin{corrige}[Exercice]
\begin{enumerate}
\item Cellule eucaryote : cellule possédant un noyau individualisé par une membrane nucléaire et des organites membranaires. La cellule animale se distingue de la cellule végétale par l'absence de paroi cellulaire et l'absence de chloroplastes (ainsi que l'absence de grande vacuole centrale).
\item Non : $\SI{150}{nm} = \SI{0,15}{\micro m}$, ce qui est inférieur au pouvoir de résolution du MO (\SI{0,2}{\micro m}). Cet organite n'est observable qu'au microscope électronique (résolution \SI{0,2}{nm}).
\item Le schéma doit présenter une membrane plasmique, un cytoplasme et un noyau, avec annotations horizontales reliées par des traits fins, un titre (« Cellule animale vue en MO ») et l'échelle.
\end{enumerate}
\end{corrige}

\begin{aretenir}[L'essentiel de la section]
\begin{itemize}
\item La cellule est l'unité structurale et fonctionnelle du vivant ; les cellules animales sont diverses de forme, adaptées à leur fonction.
\item Le MO (résolution \SI{0,2}{\micro m}) révèle membrane, cytoplasme et noyau ; le ME (résolution \SI{0,2}{nm}) révèle l'ultrastructure des organites.
\item Toute cellule animale comprend une membrane plasmique, un cytoplasme (hyaloplasme + organites) et un noyau : c'est une cellule \textbf{eucaryote sans paroi ni chloroplastes}.
\item Cette organisation, en particulier la membrane plasmique et l'absence de paroi, conditionne les échanges cellulaires que nous étudierons dans la suite de la leçon — et qui expliquent les observations d'Aminatou et la prescription du sérum physiologique.
\end{itemize}
\end{aretenir}

\coursec{Les principaux organites cellulaires et leurs rôles}

Dans la section précédente, l'observation d'une cellule animale en microscopie optique puis en microscopie électronique nous a révélé une organisation générale en trois compartiments : la \cle{membrane plasmique}, le \cle{cytoplasme} et le \cle{noyau}. Mais la microscopie électronique montre bien plus : le cytoplasme n'est pas un gel homogène, il abrite de petites structures spécialisées, comparables aux organes d'un organisme. Ce sont les \cle{organites}. Comprendre la cellule, c'est comprendre chacun d'eux : sa structure, repérable sur un cliché, et son rôle, démontré par l'expérience. C'est l'objet de cette section, véritable carte d'identité fonctionnelle de la cellule animale.

\begin{definition}[Organite]
Un \cle{organite} est une structure cellulaire spécialisée, située dans le cytoplasme, qui remplit une fonction précise indispensable à la vie de la cellule. Certains organites sont entourés d'une membrane (mitochondrie, appareil de Golgi, lysosome\ldots), d'autres non (ribosomes, centrioles).
\end{definition}

\begin{popfigure}
\begin{center}
\begin{tikzpicture}[scale=1.0]
% Cytoplasme
\fill[popcream] (0,0) ellipse (5.2 and 3.6);
\draw[popline, thick] (0,0) ellipse (5.2 and 3.6);
% Noyau
\fill[popblueL] (-1.2,0.3) circle (1.5);
\draw[popblue, thick, double distance=2pt] (-1.2,0.3) circle (1.5);
\fill[popblue] (-1.2,0.3) circle (0.45);
\draw[popblue!60, thick] plot[smooth cycle, tension=0.9] coordinates {(-2.2,0.9) (-1.0,1.4) (-0.3,0.7) (-0.5,-0.4) (-1.6,-0.9) (-2.4,-0.2)};
% Mitochondries
\fill[poporangeL] (2.6,1.6) ellipse (0.9 and 0.45);
\draw[poporange, thick] (2.6,1.6) ellipse (0.9 and 0.45);
\draw[poporange, thick] (1.9,1.6) .. controls (2.2,1.85) and (2.5,1.35) .. (2.8,1.6) .. controls (3.0,1.8) and (3.2,1.5) .. (3.4,1.65);
\fill[poporangeL] (-3.4,-1.9) ellipse (0.8 and 0.4);
\draw[poporange, thick] (-3.4,-1.9) ellipse (0.8 and 0.4);
\draw[poporange, thick] (-4.0,-1.9) .. controls (-3.7,-1.65) and (-3.4,-2.15) .. (-3.1,-1.9) .. controls (-2.9,-1.7) and (-2.8,-2.0) .. (-2.7,-1.85);
% REG
\draw[popgreen, thick] (0.6,-1.6) .. controls (1.2,-1.2) and (1.8,-2.0) .. (2.4,-1.6) .. controls (3.0,-1.2) and (3.6,-2.0) .. (4.2,-1.7);
\draw[popgreen, thick] (0.6,-1.9) .. controls (1.2,-1.5) and (1.8,-2.3) .. (2.4,-1.9) .. controls (3.0,-1.5) and (3.6,-2.3) .. (4.2,-2.0);
\foreach \p in {(1.0,-1.42),(1.7,-1.85),(2.6,-1.5),(3.3,-1.8),(3.9,-1.75)} \fill[popdark] \p circle (0.07);
% Golgi
\foreach \i in {0,1,2,3} \draw[poppurple, thick] (-3.6+0.15*\i, 1.9-0.28*\i) arc (160:20:0.9-0.08*\i);
\fill[poppurpleL] (-2.3,2.3) circle (0.16);
\fill[poppurpleL] (-2.5,1.5) circle (0.13);
% Lysosome
\fill[poppinkL] (3.6,0.2) circle (0.35);
\draw[poppink, thick] (3.6,0.2) circle (0.35);
% Centrioles
\draw[popgold, thick] (0.4,2.4) rectangle (0.75,2.9);
\draw[popgold, thick] (0.35,2.55) rectangle (0.85,2.75);
% REL
\draw[popgreen!70!popblue, thick] (-4.3,0.6) .. controls (-3.9,1.0) and (-3.5,0.3) .. (-3.1,0.7);
\draw[popgreen!70!popblue, thick] (-4.3,0.35) .. controls (-3.9,0.75) and (-3.5,0.05) .. (-3.1,0.45);
% Légendes fléchées
\node[font=\small, anchor=west] at (5.4,2.6) {1. Membrane plasmique};
\draw[->, popmuted] (5.4,2.5) -- (4.6,2.9);
\node[font=\small, anchor=west] at (5.4,1.6) {2. Mitochondrie};
\draw[->, popmuted] (5.4,1.5) -- (3.5,1.65);
\node[font=\small, anchor=west] at (5.4,0.3) {3. Lysosome};
\draw[->, popmuted] (5.4,0.2) -- (3.95,0.2);
\node[font=\small, anchor=west] at (5.4,-1.2) {4. REG + ribosomes};
\draw[->, popmuted] (5.4,-1.3) -- (4.2,-1.8);
\node[font=\small, anchor=east] at (-5.4,2.6) {5. Appareil de Golgi};
\draw[->, popmuted] (-5.4,2.5) -- (-3.4,2.1);
\node[font=\small, anchor=east] at (-5.4,1.2) {6. REL};
\draw[->, popmuted] (-5.4,1.1) -- (-4.2,0.55);
\node[font=\small, anchor=east] at (-5.4,-0.2) {7. Noyau (chromatine)};
\draw[->, popmuted] (-5.4,-0.3) -- (-2.7,0.1);
\node[font=\small, anchor=east] at (-5.4,-1.4) {8. Nucléole};
\draw[->, popmuted] (-5.4,-1.5) -- (-1.7,0.15);
\node[font=\small, anchor=east] at (-5.4,-2.6) {9. Cytoplasme (cytosol)};
\draw[->, popmuted] (-5.4,-2.7) -- (-2.2,-2.6);
\node[font=\small] at (1.3,3.2) {10. Centrioles};
\draw[->, popmuted] (1.2,3.05) -- (0.7,2.85);
\end{tikzpicture}
\end{center}
\legende{Schéma d'interprétation d'une cellule animale observée en microscopie électronique. Chaque organite possède une signature structurale : double membrane crêtée (mitochondrie), saccules empilés (Golgi), réseau membraneux parsemé de grains (REG).}
\end{popfigure}

\courssub{La membrane plasmique : une frontière sélective}

\textbf{Observation et question.} Une cellule placée dans un milieu coloré ne laisse pas entrer n'importe quel colorant : certains traversent, d'autres restent dehors. La membrane plasmique n'est donc pas une simple enveloppe : elle \emph{choisit} ce qui passe. Quelle structure explique ce choix ?

\textbf{Structure : le modèle de la mosaïque fluide.} En microscopie électronique, la membrane plasmique apparaît comme une double couche sombre d'épaisseur voisine de \SI{7,5}{\nano\meter}. Le modèle admis est celui de la \cle{mosaïque fluide} (Singer et Nicolson, 1972) : une \cle{bicouche phospholipidique} dans laquelle flottent des \cle{protéines}. Chaque phospholipide possède une tête hydrophile (attirée par l'eau) tournée vers les milieux aqueux (extérieur et cytosol) et deux queues hydrophobes (qui fuient l'eau) repliées vers l'intérieur de la membrane. Les protéines, enfoncées à des profondeurs variables comme les pièces d'une mosaïque, assurent le passage sélectif des substances : ce sont des canaux, des transporteurs, des récepteurs.

\begin{popfigure}
\begin{center}
\begin{tikzpicture}[scale=1.0]
% Bicouche phospholipidique : têtes + queues
\foreach \x in {0,0.7,1.4,2.1,2.8,3.5,4.9,5.6,6.3,7.0,7.7,8.4} {
  \fill[popblue] (\x,1.0) circle (0.18);
  \draw[popblue, thick] (\x-0.05,0.85) -- (\x-0.08,0.25);
  \draw[popblue, thick] (\x+0.05,0.85) -- (\x+0.08,0.25);
  \fill[popblue] (\x,-0.6) circle (0.18);
  \draw[popblue, thick] (\x-0.05,-0.45) -- (\x-0.08,0.15);
  \draw[popblue, thick] (\x+0.05,-0.45) -- (\x+0.08,0.15);
}
% Protéine transmembranaire (canal)
\fill[poporangeL, rounded corners=3pt] (3.85,1.25) rectangle (4.75,-0.85);
\draw[poporange, thick, rounded corners=3pt] (3.85,1.25) rectangle (4.75,-0.85);
\fill[popcream] (4.15,1.25) rectangle (4.45,-0.85);
\draw[poporange, thick] (4.15,1.25) -- (4.15,-0.85);
\draw[poporange, thick] (4.45,1.25) -- (4.45,-0.85);
% Protéine de surface
\fill[poppurpleL] (1.6,1.45) ellipse (0.55 and 0.3);
\draw[poppurple, thick] (1.6,1.45) ellipse (0.55 and 0.3);
% Glucide (chaîne)
\draw[poppink, thick] (1.6,1.75) -- (1.6,2.15);
\fill[poppink] (1.6,2.25) circle (0.1);
\fill[poppink] (1.75,2.05) circle (0.1);
% Annotations
\node[font=\small, anchor=west] at (8.9,1.6) {Têtes hydrophiles};
\draw[->, popmuted] (8.9,1.5) -- (7.9,1.05);
\node[font=\small, anchor=west] at (8.9,0.2) {Queues hydrophobes};
\draw[->, popmuted] (8.9,0.1) -- (7.9,0.2);
\node[font=\small] at (4.3,1.9) {Protéine canal};
\node[font=\small, anchor=east] at (0.9,2.3) {Glucide (récepteur)};
\draw[->, popmuted] (0.95,2.2) -- (1.5,2.2);
\node[font=\small] at (4.3,-1.5) {Milieu extracellulaire (haut) / Cytosol (bas)};
\end{tikzpicture}
\end{center}
\legende{Modèle de la mosaïque fluide de la membrane plasmique : bicouche de phospholipides (têtes hydrophiles vers l'extérieur, queues hydrophobes vers l'intérieur) parsemée de protéines (canaux, transporteurs, récepteurs) et de glucides de reconnaissance.}
\end{popfigure}

\begin{propriete}[Perméabilité sélective — savoir admis]
La membrane plasmique est \cle{sélectivement perméable} (ou semi-perméable) : grâce à sa structure en mosaïque fluide, elle laisse passer librement certaines substances (eau, petites molécules non polaires), contrôle le passage d'autres (ions, glucose, acides aminés) via ses protéines de transport, et bloque les grosses molécules (protéines, polysaccharides). Cette propriété, admise sans démonstration au programme, est la clé de tous les échanges cellulaires étudiés dans les sections suivantes.
\end{propriete}

\courssub{Le noyau : le centre de commandement}

\textbf{Observation.} Sur les clichés de microscopie électronique, le noyau est la structure la plus volumineuse de la cellule animale (diamètre voisin de \SI{5}{\micro\meter}). Il est délimité par une \cle{membrane nucléaire} double, percée de pores, et contient deux éléments : la \cle{chromatine}, réseau de filaments d'\cle{ADN} associés à des protéines (histones), et le \cle{nucléole}, zone dense où sont fabriqués les ribosomes.

\textbf{Établissement expérimental du rôle.} L'expérience classique de greffe nucléaire chez l'algue unicellulaire \emph{Acetabularia} (Hämmerling, 1930-1950) est éclairante : cette algue possède un chapeau dont la forme dépend de l'espèce. Si l'on transplante le noyau d'une espèce A dans une algue énucléée de l'espèce B, le chapeau qui repousse est celui de l'espèce A. \emph{Interprétation :} c'est le noyau — donc l'ADN qu'il contient — qui porte l'information déterminant les caractères de la cellule.

\begin{aretenir}[Rôles du noyau]
Le noyau \textbf{stocke l'information génétique} (l'ADN de la chromatine) et \textbf{commande les activités cellulaires} : il dirige la synthèse des protéines en envoyant des messages (ARN messagers) vers le cytoplasme à travers les pores nucléaires. Une cellule énucléée (comme le globule rouge humain mature) ne peut ni se diviser ni synthétiser de nouvelles protéines ; sa durée de vie est limitée (environ 120 jours pour l'érythrocyte).
\end{aretenir}

\courssub{Les mitochondries : les centrales énergétiques}

\textbf{Structure.} La mitochondrie est un organite ovoïde (longueur voisine de \SI{1}{\micro\meter}) entouré d'une \cle{double membrane} : une membrane externe lisse et une membrane interne repliée en \cle{crêtes mitochondriales}, qui augmentent considérablement la surface d'échange. L'intérieur contient un gel, la \cle{matrice}, riche en enzymes, ADN mitochondrial et ribosomes.

\textbf{Rôle établi par l'expérience.} Par centrifugation fractionnée (différentielle), on isole les mitochondries d'un broyat cellulaire. Placées dans un milieu contenant du pyruvate (issu de la glycolyse) et de l'oxygène, elles consomment le \ce{O2} et produisent de l'\cle{ATP} (adénosine triphosphate), la molécule d'énergie directement utilisable par la cellule. Sans oxygène, cette production s'arrête. La mitochondrie est donc le siège de la \cle{respiration cellulaire} aérobie, dont le bilan global est :
\[ \ce{C6H12O6 + 6O2 -> 6CO2 + 6H2O} \quad \text{(avec production d'environ 30 à 32 ATP par glucose)} \]

\begin{exemplebox}[Des mitochondries là où l'énergie est demandée]
Les cellules très consommatrices d'énergie sont riches en mitochondries : une cellule musculaire cardiaque en contient plusieurs milliers (elle se contracte toute la vie sans repos), et la pièce intermédiaire du spermatozoïde en est gainée pour fournir l'ATP nécessaire au battement du flagelle — sans elle, pas de fécondation possible.
\end{exemplebox}

\begin{attention}[Erreur fréquente au Bac]
N'attribue \textbf{jamais} la photosynthèse aux mitochondries ! Les mitochondries réalisent la \cle{respiration cellulaire} (consommation de \ce{O2}, rejet de \ce{CO2}, production d'ATP). La photosynthèse est le rôle des chloroplastes, organites des cellules végétales, absents chez l'animal. Confondre ces deux organites est une faute grave, souvent éliminatoire, dans une copie de Terminale.
\end{attention}

\courssub{REG, ribosomes et REL : les usines de synthèse}

\textbf{Le réticulum endoplasmique granuleux (REG).} C'est un réseau de membranes formant des canaux aplatis (cisternae), dont la face cytoplasmique est parsemée de petits grains : les \cle{ribosomes}, d'où son aspect « granuleux ». Les ribosomes (particules ribonucléoprotéiques, non membranaires) sont le siège de la \cle{synthèse des protéines} (traduction) : ils y assemblent les acides aminés selon les instructions venues du noyau (ARNm). Les protéines fabriquées par les ribosomes fixés sur le REG sont destinées à être exportées hors de la cellule (protéines de sécrétion), adressées aux lysosomes ou insérées dans les membranes.

\textbf{Le réticulum endoplasmique lisse (REL).} De structure semblable (cisternae tubulaires) mais dépourvu de ribosomes (d'où son aspect « lisse »), il est le siège de la \cle{synthèse des lipides} (phospholipides des membranes, hormones stéroïdes). Il est particulièrement développé dans les cellules qui sécrètent des hormones lipidiques (cellules interstitielles des testicules, cellules de la corticosurrénale, cellules du corps jaune de l'ovaire) et dans les cellules hépatiques (hépatocytes), où il participe aussi à la détoxication (métabolisme des médicaments, de l'alcool, du bilirubine).

\courssub{L'appareil de Golgi : le centre de tri et d'expédition}

L'appareil de Golgi apparaît comme un empilement de \cle{saccules} (ou cisternae) membraneux aplatis, généralement au nombre de 3 à 8, entourés de petites \cle{vésicules} qui s'en détachent par bourgeonnement. Son rôle, établi par des expériences de marquage radioactif (pulse-chase, voir exercice résolu ci-dessous), est triple : \textbf{maturation} (modifications post-traductionnelles : glycosylation, sulfatation, clivage protéolytique) des protéines venues du REG, \textbf{tri et conditionnement} en vésicules sécrétrices ou lysosomales, et \textbf{expédition} vers leurs destinations finales (membrane plasmique, espace extracellulaire, lysosomes).

\begin{popfigure}
\begin{center}
\begin{tikzpicture}[scale=1.0, >=Stealth]
% Noyau
\fill[popblueL] (0,0) circle (0.9);
\draw[popblue, thick] (0,0) circle (0.9);
\node[font=\small] at (0,0) {Noyau};
% REG
\draw[popgreen, thick] (1.6,0.5) .. controls (2.1,0.9) and (2.6,0.2) .. (3.1,0.6);
\draw[popgreen, thick] (1.6,0.1) .. controls (2.1,0.5) and (2.6,-0.2) .. (3.1,0.2);
\foreach \p in {(1.9,0.62),(2.4,0.35),(2.9,0.5)} \fill[popdark] \p circle (0.06);
\node[font=\small, popgreen!60!popdark] at (2.4,1.15) {REG};
% Golgi
\foreach \i in {0,1,2} \draw[poppurple, thick] (4.6+0.1*\i, 0.75-0.3*\i) arc (150:30:0.75-0.07*\i);
\node[font=\small, poppurple] at (5.6,1.15) {Golgi};
% Vésicule
\fill[poppurpleL] (6.9,0.4) circle (0.22);
\draw[poppurple, thick] (6.9,0.4) circle (0.22);
\node[font=\small] at (6.9,-0.25) {Vésicule};
% Membrane
\draw[popline, very thick] (8.3,-1.2) -- (8.3,1.6);
\node[font=\small, rotate=90] at (8.05,0.2) {Membrane};
% Exocytose
\fill[poporangeL] (8.75,0.4) circle (0.1);
\fill[poporangeL] (9.0,0.55) circle (0.1);
\fill[poporangeL] (9.05,0.3) circle (0.1);
\node[font=\small, anchor=west] at (9.2,0.4) {Sécrétion};
% Flèches
\draw[->, thick, poporange] (0.95,0) -- (1.5,0.3);
\draw[->, thick, poporange] (3.2,0.45) -- (4.3,0.55);
\draw[->, thick, poporange] (5.9,0.5) -- (6.6,0.42);
\draw[->, thick, poporange] (7.15,0.4) -- (8.25,0.4);
% Étapes
\node[font=\footnotesize, popmuted] at (0.9,-0.7) {1. Information (ADN)};
\node[font=\footnotesize, popmuted] at (2.4,-0.7) {2. Synthèse};
\node[font=\footnotesize, popmuted] at (5.2,-0.7) {3. Maturation};
\node[font=\footnotesize, popmuted] at (7.6,-0.7) {4. Exocytose};
\end{tikzpicture}
\end{center}
\legende{Trajet d'une protéine de sécrétion : synthétisée par les ribosomes du REG sur instruction de l'ADN nucléaire (transcription $\rightarrow$ ARNm $\rightarrow$ traduction), elle transite vers l'appareil de Golgi (face \emph{cis} $\rightarrow$ face \emph{trans}) pour maturation et tri, est empaquetée dans une vésicule qui fusionne avec la membrane plasmique et libère la protéine hors de la cellule (exocytose).}
\end{popfigure}

\courssub{Les lysosomes et les centrioles}

\textbf{Les lysosomes} sont de petites vésicules sphériques (diamètre voisin de \SI{0,5}{\micro\meter}) délimitées par une membrane unique et remplies d'\cle{enzymes hydrolytiques} (protéases, lipases, nucléases, glycosidases) actives à pH acide (vers 5). Ils assurent la \cle{digestion intracellulaire} : destruction des particules ingérées par phagocytose (chez les macrophages et neutrophiles qui éliminent les microbes, par exemple lors d'une infection) et recyclage des organites usagés (autophagie). Leurs enzymes sont synthétisées par le REG, maturées et adressées par l'appareil de Golgi (reconnaissance du marqueur mannose-6-phosphate) — bel exemple de coopération entre organites.

\textbf{Les centrioles} sont deux petits cylindres creux formés de neuf triplets de microtubules, disposés perpendiculairement l'un à l'autre près du noyau, au sein du \cle{centrosome} (centre organisateur des microtubules). Ils jouent un rôle essentiel dans la \cle{division cellulaire} (mitose et méiose) : ils organisent le fuseau mitotique (fuseau acentriolaire chez les végétaux) qui répartit les chromosomes entre les cellules filles. Nous y reviendrons en détail dans la leçon consacrée à la reproduction cellulaire.

\courssub{Bilan : tableau récapitulatif et méthode d'annotation}

\begin{center}
\begin{tabularx}{\linewidth}{|l|X|X|}
\hline
\rowcolor{popblueL} \textbf{Organite} & \textbf{Structure (signature en ME)} & \textbf{Rôle principal} \\
\hline
Membrane plasmique & Bicouche phospholipidique + protéines intégrales/périphériques (mosaïque fluide) & Barrière d'échanges sélectivement perméable ; reconnaissance cellulaire \\
\hline
Noyau & Double membrane nucléaire à pores ; chromatine (ADN + histones) ; nucléole & Stockage de l'information génétique ; transcription (synthèse ARNm) ; assemblage sous-unités ribosomiques \\
\hline
Mitochondrie & Double membrane ; membrane interne en crêtes ; matrice (ADNmt, ribosomes, enzymes) & Respiration cellulaire aérobie ; production massive d'ATP (phosphorylation oxydative) \\
\hline
REG + ribosomes & Cisternae aplaties parsemées de ribosomes (côté cytosol) & Synthèse des protéines de sécrétion, membranaires, lysosomales (traduction) \\
\hline
REL & Réseau tubulaire lisse (sans ribosomes) & Synthèse des lipides (stéroïdes, phospholipides) ; détoxication (foie) ; stockage \ce{Ca^2+} (muscle) \\
\hline
Appareil de Golgi & Empilement de saccules (face \emph{cis} / \emph{trans}) + vésicules de bourgeonnement & Maturation (glycosylation\ldots), tri, conditionnement, expédition des protéines et lipides \\
\hline
Lysosome & Vésicule sphérique à membrane unique, contenu dense (enzymes acides) & Digestion intracellulaire (hétérophagie, autophagie) ; recyclage macromolécules \\
\hline
Centrioles / Centrosome & Deux cylindres perpendiculaires (9 triplets microtubules) + matériel péricentriolaire & Organisation du fuseau mitotique ; formation cils/flagelles (basal bodies) \\
\hline
\end{tabularx}
\end{center}

\begin{methode}[Annoter un schéma d'organite en trois questions]
Face à un cliché ou un schéma de microscopie électronique à légender, pose-toi toujours les questions dans cet ordre :
\begin{enumerate}
\item \textbf{Y a-t-il une membrane ?} Si non : ribosome (grain isolé \SI{\approx 20}{\nano\meter}, souvent en grappes = polysomes) ou centriole (cylindre de triplets microtubulaires). Si oui, continue.
\item \textbf{La membrane est-elle double ou simple ?} Double membrane : noyau (volumineux, pores visibles, chromatine périphérique) ou mitochondrie (ovoïde, crêtes internes). Membrane simple : continue.
\item \textbf{Quelle est la morphologie du système membraneux ?} Saccules empilés parallèles : Golgi. Cisternae aplatis en réseau : REG (si grains) ou REL (si lisse). Tubules anastomosés : REL. Petite sphère dense homogène : lysosome.
\end{enumerate}
Conclus en nommant l'organite \emph{et} en justifiant par la structure observée : le correcteur attend cette justification structurale.
\end{methode}

\begin{exoresolu}[Le trajet d'une protéine de sécrétion : l'expérience de « pulse-chase »]
\textbf{Énoncé.} Des cellules pancréatiques (qui sécrètent l'insuline, une protéine) sont incubées \SI{3}{\minute} avec des acides aminés radioactifs (pulse), puis replacées dans un milieu non radioactif en excès (chase). On localise la radioactivité par autoradiographie à différents temps : à $t = \SI{3}{\minute}$, elle est dans le REG ; à $t = \SI{20}{\minute}$, dans l'appareil de Golgi ; à $t = \SI{60}{\minute}$, dans des vésicules proches de la membrane plasmique ; à $t = \SI{90}{\minute}$, à l'extérieur de la cellule. Interprète ces résultats et conclus.
\tcblower
\textbf{Solution.} \emph{Analyse :} la radioactivité suit un trajet ordonné dans le temps : REG $\rightarrow$ Golgi $\rightarrow$ vésicules sécrétrices $\rightarrow$ milieu extracellulaire. \emph{Interprétation :} les acides aminés radioactifs sont d'abord incorporés dans les chaînes polypeptidiques en cours de synthèse au niveau des ribosomes fixés sur le REG (étape de traduction, \SI{3}{\minute}) ; les protéines néo-synthétisées (préproinsuline $\rightarrow$ proinsuline) transitent ensuite dans les vésicules de transport vers l'appareil de Golgi où elles arrivent vers \SI{20}{\minute} (maturation : clivage en insuline, conditionnement dans des granules de sécrétion) ; les granules migrent vers la périphérie (\SI{60}{\minute}) et fusionnent avec la membrane plasmique par exocytose, libérant l'insuline marquée à l'extérieur (\SI{90}{\minute}). \emph{Conclusion :} la sécrétion d'une protéine résulte de la coopération ordonnée de plusieurs organites membranaires : REG (synthèse et translocation dans la lumière), Golgi (maturation, tri, stockage), vésicules (transport vectoriel) et membrane plasmique (fusion et expulsion). C'est la voie de sécrétion constitutive ou régulée, valable pour toute protéine exportée (insuline, enzymes digestives, anticorps, mucus).
\end{exoresolu}

\begin{savaistu}[Quand les lysosomes se dérèglent : maladies lysosomales et biotechnologies]
Certaines maladies génétiques graves, dites \emph{maladies lysosomales} (ex. : maladie de Gaucher, maladie de Pompe, mucoviscidose — bien que cette dernière soit due à un défaut de canal \ce{Cl-} CFTR, pas une enzyme lysosomale, corrigeons : maladie de Tay-Sachs, mucopolysaccharidoses), proviennent de l'absence ou du déficit d'une enzyme hydrolytique dans les lysosomes : les substrats non dégradés s'accumulent, distendent l'organite et détruisent la cellule (stockage lysosomal). À l'inverse, la médecine utilise la voie de sécrétion REG--Golgi : l'insuline humaine injectée aux patients diabétiques, suivis notamment dans les hôpitaux de Yaoundé et de Douala, est aujourd'hui produite industriellement par génie génétique. On insère le gène de l'insuline humaine dans des bactéries (\emph{E. coli}) ou des levures (\emph{S. cerevisiae}) qui la sécrètent, ou dans des cellules animales (cellules CHO) pour les protéines complexes nécessitant des glycosylations spécifiques. C'est une application directe de la connaissance du trafic intracellulaire.
\end{savaistu}

En résumé, la cellule animale fonctionne comme une cité organisée : le noyau en est le centre de commandement (bibliothèque génétique), les mitochondries les centrales énergétiques, le REG et le REL les usines de synthèse, le Golgi le centre de tri, de maturation et d'expédition, les lysosomes le service de recyclage et de destruction, et la membrane plasmique la frontière douanière sélective. C'est précisément le fonctionnement de cette frontière — comment l'eau, les substances dissoutes et les particules la franchissent — que nous allons maintenant étudier dans les sections suivantes.

\coursec{Les échanges d'eau : osmose et dialyse}

Dans la section précédente, nous avons exploré l'organisation interne de la cellule animale et le rôle de chacun de ses organites. Tous ces organites baignent dans le hyaloplasme, un milieu essentiellement aqueux, et la cellule elle-même est entourée d'un liquide extracellulaire, lui aussi aqueux. Or la membrane plasmique, si remarquable soit-elle, n'est pas une barrière étanche : elle laisse passer certaines substances et en bloque d'autres. Comment l'eau, constituant majeur de la cellule (environ 70 \% de sa masse), circule-t-elle entre l'intérieur et l'extérieur ? Que se passe-t-il lorsque les concentrations ne sont pas les mêmes de part et d'autre de la membrane ? C'est à ces questions, d'une importance capitale pour la survie de la cellule, que cette section est consacrée. Tu verras que deux phénomènes physiques simples, l'\cle{osmose} et la \cle{dialyse}, expliquent à eux seuls pourquoi une perfusion mal dosée peut tuer un patient à l'hôpital, et pourquoi le sérum physiologique utilisé dans tous les centres de santé du Cameroun contient exactement 0,9 \% de NaCl.

\courssub{Mise en évidence expérimentale : l'expérience de Dutrochet}

\textbf{Observation de départ.} Un raisin sec plongé dans un verre d'eau pure gonfle progressivement et redevient presque un raisin frais. Inversement, des fruits placés dans du sirop très concentré rendent de l'eau et se ratatinent (c'est le principe des confitures). Ces observations quotidiennes suggèrent que l'eau se déplace spontanément d'un milieu vers un autre. Mais dans quel sens, et selon quelle règle ? C'est le physicien et physiologiste français René Dutrochet qui, dès 1826, a conçu une expérience décisive pour répondre à cette question.

\begin{experience}[Expérience de Dutrochet (1826)]
\textbf{Matériel :} une vessie de porc vidée (membrane naturelle semi-perméable), un tube de verre, de l'eau sucrée (solution de saccharose), un cristallisoir rempli d'eau pure.

\textbf{Protocole :}
\begin{enumerate}
\item On remplit la vessie de porc d'eau sucrée et on la relie à un long tube de verre vertical.
\item On plonge l'ensemble dans un cristallisoir contenant de l'eau pure.
\item On marque le niveau initial du liquide dans le tube, puis on observe l'évolution du niveau au cours du temps.
\end{enumerate}

\textbf{Observations :} le niveau du liquide monte progressivement dans le tube, puis se stabilise à une certaine hauteur. Le volume de liquide à l'intérieur de la vessie a donc augmenté.

\textbf{Interprétation :} la vessie de porc est une membrane \cle{semi-perméable} : elle laisse passer les molécules d'eau (le solvant) mais pas les molécules de saccharose (le soluté, trop grosses). Seule l'eau a pu traverser la membrane, et elle l'a fait de l'extérieur (eau pure) vers l'intérieur (eau sucrée), c'est-à-dire \textbf{du milieu le moins concentré en soluté vers le milieu le plus concentré en soluté}. L'arrêt de la montée s'explique par l'équilibre entre le flux d'eau entrant et la pression exercée par la colonne d'eau montée dans le tube, qui s'oppose à ce flux.

\textbf{Conclusion :} à travers une membrane semi-perméable, l'eau diffuse spontanément vers le milieu le plus concentré en solutés. Ce phénomène est l'\cle{osmose}.
\end{experience}

\begin{popfigure}
\begin{center}
\begin{tikzpicture}[scale=0.9, every node/.style={font=\small}]
% ===== Etat initial =====
\begin{scope}
% cristallisoir
\draw[thick, popdark] (-2.2,0) -- (-2.2,3) -- (2.2,3) -- (2.2,0);
\fill[popblue!15] (-2.2,0) rectangle (2.2,2.6);
\node[popblue, anchor=west] at (-2.1,2.35) {\footnotesize eau pure};
% vessie
\draw[thick, poporange, fill=popgold!25] (0,1.1) ellipse (0.95 and 0.7);
\node[popdark] at (0,1.1) {\footnotesize eau sucrée};
% tube
\draw[thick, popdark] (-0.12,1.7) -- (-0.12,4.4);
\draw[thick, popdark] (0.12,1.7) -- (0.12,4.4);
\fill[popgold!25] (-0.12,1.7) rectangle (0.12,2.3);
\draw[popink, thick] (-0.12,2.3) -- (0.12,2.3);
\node[right, popink] at (0.15,2.3) {\footnotesize niveau initial};
% fleches eau entrante
\foreach \a in {150, 180, 210, 330, 0, 30}{
  \draw[->, popblue, thick] ({\a}:1.55) ++(0,1.1) -- ({\a}:1.05) ++(0,1.1);
}
\node[popdark, font=\bfseries] at (0,-0.5) {État initial};
\end{scope}
% ===== Etat final =====
\begin{scope}[xshift=7.5cm]
\draw[thick, popdark] (-2.2,0) -- (-2.2,3) -- (2.2,3) -- (2.2,0);
\fill[popblue!15] (-2.2,0) rectangle (2.2,2.6);
\draw[thick, poporange, fill=popgold!25] (0,1.15) ellipse (1.25 and 0.85);
\node[popdark] at (0,1.15) {\footnotesize liquide dilué};
\draw[thick, popdark] (-0.12,1.9) -- (-0.12,4.4);
\draw[thick, popdark] (0.12,1.9) -- (0.12,4.4);
\fill[popgold!25] (-0.12,1.9) rectangle (0.12,3.6);
\draw[popink, thick] (-0.12,3.6) -- (0.12,3.6);
\node[right, popink] at (0.15,3.6) {\footnotesize niveau final};
\draw[<->, poppurple, thick] (1.6,2.3) -- (1.6,3.6) node[midway, right, font=\footnotesize] {$h$};
\node[popdark, font=\bfseries] at (0,-0.5) {État final};
\end{scope}
% fleche evolution
\draw[->, very thick, popgreen] (2.7,1.5) -- (4.8,1.5) node[midway, above, font=\footnotesize] {temps};
\end{tikzpicture}
\end{center}
\legende{Expérience de Dutrochet. À gauche : état initial — la vessie de porc (membrane semi-perméable) remplie d'eau sucrée est plongée dans l'eau pure ; les flèches bleues indiquent le flux d'eau entrant. À droite : état final — le liquide est monté d'une hauteur $h$ dans le tube car l'eau a traversé la membrane vers le milieu le plus concentré en saccharose ; la montée s'arrête lorsque la pression de la colonne d'eau compense le flux osmotique.}
\end{popfigure}

\begin{attention}[Le sens du flux d'eau : une erreur à ne jamais commettre]
Beaucoup d'élèves disent que « l'eau va vers l'eau » ou que « l'eau va diluer le milieu concentré ». C'est une formulation dangereuse qui fait perdre des points au Baccalauréat. La formulation rigoureuse est : \textbf{l'eau diffuse du milieu le moins concentré en solutés (hypotonique) vers le milieu le plus concentré en solutés (hypertonique)}. C'est le soluté qui, en quelque sorte, « attire » l'eau : plus un milieu contient de particules dissoutes, plus sa concentration en eau libre est faible, et plus l'eau tend à y entrer. Retiens aussi que ce sont bien les molécules d'\textbf{eau} qui se déplacent dans l'expérience de Dutrochet, pas les molécules de sucre : la membrane les bloque.
\end{attention}

\courssub{Mesure de la pression osmotique : l'osmomètre de Pfeiffer}

L'expérience de Dutrochet est qualitative : elle montre le phénomène mais ne le mesure pas. En 1877, le botaniste allemand Wilhelm Pfeiffer perfectionne le dispositif en fabriquant une membrane semi-perméable artificielle (une paroi d'argile poreuse imprégnée de ferrocyanure de cuivre, imperméable au saccharose mais perméable à l'eau) et en y adaptant un manomètre : c'est l'\cle{osmomètre} de Pfeiffer.

\begin{experience}[Expérience de Pfeiffer (1877)]
\textbf{Protocole :} on remplit la cellule osmométrique (membrane artificielle semi-perméable) avec des solutions de saccharose de concentrations croissantes (par exemple 58,5 ; 117 ; 175,5 g/L, soit environ 0,17 ; 0,34 ; 0,51 mol/L, la masse molaire du saccharose étant de 342 g/mol) et on plonge l'osmomètre dans l'eau pure. On mesure, pour chaque concentration, la pression maximale atteinte à l'équilibre, appelée \cle{pression osmotique} et notée $\pi$.

\textbf{Résultats :} la pression osmotique mesurée augmente avec la concentration de la solution. À température constante, les résultats sont proportionnels à la concentration molaire du soluté : quand la concentration double, la pression osmotique double.
\end{experience}

\begin{popfigure}
\begin{center}
\begin{tikzpicture}[scale=0.85, every node/.style={font=\small}]
% osmometre
\draw[thick, popdark] (-1.8,0) -- (-1.8,2.6) -- (1.8,2.6) -- (1.8,0);
\fill[popblue!15] (-1.8,0) rectangle (1.8,2.2);
\node[popblue] at (-1.15,1.9) {\footnotesize eau pure};
\draw[thick, poporange, fill=popgold!25] (0,0.9) circle (0.6);
\node[popdark, font=\footnotesize] at (0,0.9) {solution};
\draw[thick, popdark] (-0.1,1.5) -- (-0.1,3.4);
\draw[thick, popdark] (0.1,1.5) -- (0.1,3.4);
% manometre
\draw[thick, popdark, fill=popcream] (0,3.9) circle (0.5);
\draw[->, thick, popink] (0,3.9) -- (0.35,4.15);
\node[popdark, font=\footnotesize] at (0,4.6) {manomètre};
\foreach \a in {200, 230, 310, 340}{
  \draw[->, popblue, thick] ({\a}:1.15) ++(0,0.9) -- ({\a}:0.72) ++(0,0.9);
}
\node[popdark, font=\bfseries] at (0,-0.5) {Osmomètre de Pfeiffer};
\end{tikzpicture}
\hfill
\begin{tikzpicture}
\begin{axis}[
  width=0.52\linewidth, height=5.2cm,
  xlabel={Concentration en saccharose (mol/L)},
  ylabel={$\pi$ (atm)},
  xmin=0, xmax=0.6, ymin=0, ymax=14,
  xtick={0,0.17,0.34,0.51}, ytick={0,4,8,12},
  grid=both, grid style={popline!40},
  axis line style={popdark},
  label style={font=\footnotesize}, tick label style={font=\footnotesize},
]
\addplot[popcurve, domain=0:0.55, samples=50] {23.6*x};
\addplot[only marks, mark=*, mark size=1.8pt, color=poporange] coordinates {(0.17,4.0) (0.34,8.0) (0.51,12.0)};
\end{axis}
\end{tikzpicture}
\end{center}
\legende{À gauche : schéma de l'osmomètre de Pfeiffer — la solution à étudier est enfermée dans la cellule à membrane semi-perméable plongée dans l'eau pure ; l'entrée d'eau (flèches bleues) élève la pression mesurée par le manomètre. À droite : évolution de la pression osmotique $\pi$ en fonction de la concentration molaire du saccharose à température constante : la relation est une proportionnalité (droite passant par l'origine).}
\end{popfigure}

\begin{propriete}[Pression osmotique]
La \cle{pression osmotique} $\pi$ d'une solution est la pression qu'il faut exercer sur cette solution pour empêcher l'entrée d'eau à travers une membrane semi-perméable la séparant du solvant pur. À température constante, elle est proportionnelle à la concentration molaire $C$ du soluté :
\[ \pi = R \times T \times C \]
où $R$ est la constante des gaz parfaits ($R = \num{0,082}$ L·atm·mol$^{-1}$·K$^{-1}$) et $T$ la température absolue en kelvins. Plus une solution est concentrée en solutés, plus sa pression osmotique est élevée, et plus elle « attire » l'eau. (La formule n'est pas exigible ; la relation de proportionnalité entre $\pi$ et $C$ à température constante, elle, doit être connue.)
\end{propriete}

\courssub{Définitions rigoureuses et vocabulaire}

\begin{definition}[Osmose]
L'\cle{osmose} est le passage du solvant (l'eau) à travers une membrane semi-perméable, du milieu le moins concentré en solutés (milieu \cle{hypotonique}) vers le milieu le plus concentré en solutés (milieu \cle{hypertonique}). C'est un phénomène passif : il ne consomme aucune énergie cellulaire.
\end{definition}

\begin{definition}[Dialyse]
La \cle{dialyse} est la diffusion des petites molécules dissoutes (ions, glucose, urée...) à travers une membrane perméable, de leur milieu de forte concentration vers leur milieu de faible concentration. Contrairement à l'osmose, ce sont ici les \textbf{solutés} qui se déplacent, et non le solvant. La membrane de dialyse laisse passer l'eau et les petites molécules, mais retient les grosses molécules (protéines) et les cellules.
\end{definition}

\begin{aretenir}[Le vocabulaire indispensable]
Soient deux milieux séparés par une membrane :
\begin{itemize}
\item le milieu le \textbf{moins} concentré en solutés est dit \cle{hypotonique} ;
\item le milieu le \textbf{plus} concentré en solutés est dit \cle{hypertonique} ;
\item deux milieux de \textbf{même} concentration sont dits \cle{isotoniques} l'un par rapport à l'autre.
\end{itemize}
Attention à la nuance entre les membranes : une membrane \cle{semi-perméable} ne laisse passer que l'eau (cas de l'osmose) ; une membrane \cle{perméable} (ou dialysante) laisse passer l'eau \textbf{et} les petites molécules dissoutes (cas de la dialyse). La membrane plasmique de la cellule animale joue les deux rôles selon les substances considérées : on dit qu'elle est \emph{sélectivement} perméable.
\end{aretenir}

\begin{methode}[Interpréter une expérience d'osmose en quatre étapes]
Face à un document expérimental (schéma, tableau de résultats, photographies), procède toujours ainsi :
\begin{enumerate}
\item \textbf{Identifier la membrane} qui sépare les deux milieux et préciser sa nature (semi-perméable : eau seule ; perméable : eau + petits solutés).
\item \textbf{Comparer les concentrations} des deux milieux : lequel est hypotonique ? lequel est hypertonique ?
\item \textbf{Déduire le sens du flux d'eau} : l'eau va toujours du milieu hypotonique vers le milieu hypertonique (et les petits solutés, si la membrane est perméable, diffusent en sens inverse, du milieu concentré vers le milieu dilué).
\item \textbf{Conclure sur l'état de la cellule ou du système} : gonflement, éclatement, rétraction, ou absence de variation si les milieux sont isotoniques.
\end{enumerate}
\end{methode}

\courssub{Application à la cellule animale : le comportement des hématies}

La cellule animale ne possède pas de paroi rigide (contrairement à la cellule végétale) : sa seule enveloppe est la membrane plasmique, fragile et déformable. Elle est donc particulièrement sensible aux variations de concentration du milieu extérieur. Les hématies (globules rouges), faciles à prélever et à observer, constituent le matériel classique pour étudier l'osmose cellulaire.

\begin{experience}[Hématies placées dans des solutions de NaCl de concentrations différentes]
\textbf{Protocole :} on dépose une goutte de sang dans trois tubes contenant respectivement une solution de NaCl à 0,3 \% (hypotonique), à 0,9 \% (isotonique) et à 3 \% (hypertonique). On observe au microscope optique après quelques minutes.

\textbf{Observations :}
\begin{itemize}
\item \textbf{NaCl à 0,3 \% (milieu hypotonique) :} l'eau entre massivement dans les hématies ; elles gonflent (\cle{turgescence}), deviennent sphériques, puis leur membrane éclate : c'est l'\cle{hémolyse}. Le contenu cellulaire (hémoglobine) se répand dans le milieu, qui devient rouge et translucide.
\item \textbf{NaCl à 0,9 \% (milieu isotonique) :} les flux d'eau entrant et sortant se compensent ; les hématies conservent leur forme normale de disque biconcave.
\item \textbf{NaCl à 3 \% (milieu hypertonique) :} l'eau sort des hématies ; elles se ratatinent et se déforment, leur membrane se plisse : c'est la \cle{crénation}.
\end{itemize}

\textbf{Conclusion :} le cytoplasme des hématies est isotonique par rapport à une solution de NaCl à 0,9 \%. Tout milieu moins concentré provoque l'entrée d'eau et l'éclatement ; tout milieu plus concentré provoque la sortie d'eau et la rétraction.
\end{experience}

\begin{popfigure}
\begin{center}
\begin{tikzpicture}[scale=1, every node/.style={font=\small}]
% ===== Hypotonique =====
\begin{scope}
\draw[thick, popdark, rounded corners] (-1.6,-1.9) rectangle (1.6,2.1);
\fill[popblue!10, rounded corners] (-1.6,-1.9) rectangle (1.6,2.1);
\draw[thick, poppink, fill=poppink!45] (0,0.2) circle (0.75);
\node[popdark, font=\footnotesize] at (0,0.2) {gonflée};
\foreach \a in {60, 120, 180, 240, 300, 0}{
  \draw[->, popblue, thick] ({\a}:1.35) ++(0,0.2) -- ({\a}:0.85) ++(0,0.2);
}
\node[popdark, font=\bfseries\footnotesize] at (0,1.75) {Hypotonique (0,3 \%)};
\node[popink, font=\footnotesize] at (0,-1.5) {turgescence puis hémolyse};
\end{scope}
% ===== Isotonique =====
\begin{scope}[xshift=5.2cm]
\draw[thick, popdark, rounded corners] (-1.6,-1.9) rectangle (1.6,2.1);
\fill[popgreen!8, rounded corners] (-1.6,-1.9) rectangle (1.6,2.1);
\draw[thick, poppink, fill=poppink!45] (0,0.2) ellipse (0.8 and 0.45);
\draw[thick, poppink, fill=popcream] (0,0.2) ellipse (0.35 and 0.18);
\foreach \a in {90, 270}{
  \draw[<->, popblue, thick] ({\a}:1.15) ++(0,0.2) -- ({\a}:0.62) ++(0,0.2);
}
\node[popdark, font=\bfseries\footnotesize] at (0,1.75) {Isotonique (0,9 \%)};
\node[popink, font=\footnotesize] at (0,-1.5) {forme normale, flux équilibrés};
\end{scope}
% ===== Hypertonique =====
\begin{scope}[xshift=10.4cm]
\draw[thick, popdark, rounded corners] (-1.6,-1.9) rectangle (1.6,2.1);
\fill[popgold!15, rounded corners] (-1.6,-1.9) rectangle (1.6,2.1);
\draw[thick, poppink, fill=poppink!45, decorate, decoration={zigzag, segment length=4, amplitude=1.2}] (0,0.2) circle (0.5);
\foreach \a in {60, 120, 180, 240, 300, 0}{
  \draw[<-, popblue, thick] ({\a}:1.3) ++(0,0.2) -- ({\a}:0.75) ++(0,0.2);
}
\node[popdark, font=\bfseries\footnotesize] at (0,1.75) {Hypertonique (3 \%)};
\node[popink, font=\footnotesize] at (0,-1.5) {crénation};
\end{scope}
\end{tikzpicture}
\end{center}
\legende{Comportement d'une hématie selon la concentration du milieu extérieur. Les flèches bleues indiquent le sens du flux d'eau. À gauche (milieu hypotonique) : entrée d'eau, gonflement puis éclatement (hémolyse). Au centre (milieu isotonique, NaCl à 0,9 \%) : flux équilibrés, forme normale en disque biconcave. À droite (milieu hypertonique) : sortie d'eau, rétraction et plissement de la membrane (crénation).}
\end{popfigure}

\begin{exemplebox}[Pourquoi le sérum physiologique contient-il 0,9 \% de NaCl ?]
Dans les hôpitaux et centres de santé du Cameroun, toute perfusion intraveineuse utilise du \cle{sérum physiologique}, c'est-à-dire une solution de NaCl à 0,9 \% (9 g/L). Cette solution est \textbf{isotonique par rapport au plasma sanguin} : injectée dans la veine, elle ne provoque ni entrée ni sortie nette d'eau au niveau des hématies, qui conservent leur intégrité. Si l'on perfusait un patient avec de l'eau pure (milieu très hypotonique), l'eau envahirait les hématies qui éclateraient par hémolyse, avec des conséquences potentiellement mortelles. À l'inverse, une solution trop concentrée déshydraterait les cellules sanguines. Le dosage à 0,9 \% n'est donc pas un hasard : c'est une application directe des lois de l'osmose.
\end{exemplebox}

\begin{exoresolu}[Interpréter une expérience d'osmose]
\textbf{Énoncé.} On place des hématies humaines dans trois solutions A, B et C. Après quelques minutes, on observe : dans A, les hématies ont éclaté et le milieu est devenu rouge ; dans B, les hématies ont une forme normale ; dans C, les hématies sont ratatinées et plissées. (1) Identifier, pour chaque solution, sa tonicité par rapport au cytoplasme des hématies. (2) Expliquer les observations dans les tubes A et C. (3) Laquelle de ces solutions pourrait être du sérum physiologique ?
\tcblower
\textbf{Solution détaillée.}

(1) \textbf{Tonicité des solutions.} Dans A, les hématies ont éclaté : l'eau est entrée massivement, donc la solution A est \textbf{hypotonique} (moins concentrée en solutés que le cytoplasme). Dans B, les hématies sont normales : les flux d'eau sont équilibrés, donc B est \textbf{isotonique}. Dans C, les hématies sont ratatinées : l'eau est sortie, donc C est \textbf{hypertonique}.

(2) \textbf{Explications.} Dans le tube A, la membrane plasmique des hématies, semi-perméable, sépare le cytoplasme (concentré en solutés) du milieu extérieur hypotonique : l'eau diffuse de l'extérieur vers l'intérieur, l'hématie gonfle (turgescence) puis sa membrane, sans paroi protectrice, éclate : c'est l'\textbf{hémolyse} ; l'hémoglobine libérée colore le milieu en rouge. Dans le tube C, le milieu extérieur hypertonique attire l'eau : l'eau sort de l'hématie, qui perd du volume et dont la membrane se plisse : c'est la \textbf{crénation}.

(3) Le sérum physiologique (NaCl à 0,9 \%) est isotonique par rapport au sang : c'est donc la \textbf{solution B}.
\end{exoresolu}

\begin{attention}[Pièges de rédaction au Baccalauréat]
\begin{itemize}
\item Ne confonds jamais \textbf{osmose} (mouvement de l'\emph{eau}) et \textbf{dialyse} (mouvement des \emph{petits solutés}). Les deux phénomènes peuvent se produire simultanément, en sens inverse, à travers une membrane perméable.
\item « Hypotonique » et « hypertonique » sont des termes \textbf{relatifs} : un même milieu peut être hypotonique par rapport à une solution et hypertonique par rapport à une autre. Précise toujours « par rapport à... ».
\item Le terme \textbf{plasmolyse} est réservé à la cellule \emph{végétale} : en milieu hypertonique, sa membrane plasmique se décolle de la paroi rigide, mais la cellule n'éclate pas. Chez la cellule \emph{animale} (sans paroi), on parle de \textbf{crénation} en milieu hypertonique et d'\textbf{hémolyse} (pour les hématies) en milieu hypotonique. N'emploie jamais « plasmolyse » pour une hématie.
\item Une hématie ne « meurt » pas seulement parce qu'elle gonfle : c'est l'\textbf{éclatement de la membrane} (hémolyse) qui est létal.
\end{itemize}
\end{attention}

\begin{savaistu}[La dialyse sauve des vies : le rein artificiel]
Chez les patients atteints d'insuffisance rénale chronique, les reins ne filtrent plus le sang : l'urée et d'autres déchets s'y accumulent dangereusement. L'\textbf{hémodialyse}, pratiquée dans les grands hôpitaux camerounais (Yaoundé, Douala), applique directement le principe de la dialyse : le sang du patient circule d'un côté d'une membrane perméable, et un liquide de dialyse de composition contrôlée circule de l'autre. L'urée et les solutés en excès diffusent du sang (concentré) vers le liquide de dialyse (dilué), tandis que les protéines et les cellules sanguines, trop grosses, sont retenues. Trois séances par semaine permettent ainsi de suppléer la fonction rénale défaillante.
\end{savaistu}

\begin{exercice}[Osmose et dialyse réunies]
On place dans un tube en U séparé par une membrane \emph{perméable} (dialysante) deux solutions : à gauche, une solution de glucose à 10 g/L ; à droite, une solution de glucose à 2 g/L. (1) Dans quel sens l'eau diffuse-t-elle par osmose ? (2) Dans quel sens le glucose diffuse-t-il par dialyse ? (3) Décrire l'état final du système.
\end{exercice}

\begin{corrige}[Exercice : Osmose et dialyse réunies]
(1) La solution de gauche est hypertonique par rapport à celle de droite : l'eau diffuse par osmose de la droite (hypotonique) vers la gauche (hypertonique).

(2) La membrane étant perméable aux petites molécules, le glucose diffuse par dialyse de son milieu concentré (gauche, 10 g/L) vers son milieu dilué (droite, 2 g/L).

(3) Les deux flux se poursuivent jusqu'à l'égalisation des concentrations de part et d'autre de la membrane : à l'équilibre, les deux compartiments contiennent une solution de glucose à 6 g/L (les volumes étant supposés égaux) et les flux nets s'annulent. On remarque que l'osmose et la dialyse se produisent simultanément et en sens inverse l'un de l'autre.
\end{corrige}

\begin{aretenir}[L'essentiel de la section]
\begin{itemize}
\item L'\textbf{osmose} est le passage de l'eau à travers une membrane semi-perméable, du milieu hypotonique vers le milieu hypertonique (expériences de Dutrochet et de Pfeiffer).
\item La \textbf{pression osmotique} d'une solution augmente proportionnellement avec sa concentration en solutés (osmomètre de Pfeiffer).
\item La \textbf{dialyse} est la diffusion des petites molécules dissoutes à travers une membrane perméable, de leur milieu concentré vers leur milieu dilué.
\item Une cellule animale en milieu hypotonique gonfle puis éclate (hémolyse) ; en milieu isotonique (NaCl à 0,9 \% pour le sang), elle garde sa forme ; en milieu hypertonique, elle se ratatine (crénation).
\item Le sérum physiologique à 0,9 \% de NaCl est isotonique par rapport au sang : c'est pourquoi il est utilisé pour les perfusions.
\end{itemize}
\end{aretenir}

Maintenant que tu maîtrises les échanges d'eau, il reste à comprendre comment la cellule gère les substances dissoutes qui ne peuvent pas traverser librement sa membrane : c'est l'objet de la section suivante, consacrée aux transports passif et actif.

\coursec{Les échanges de substances dissoutes : transports passif et actif}

Dans la section précédente, nous avons vu que l'eau traverse la membrane plasmique par osmose, toujours du milieu le moins concentré en solutés vers le milieu le plus concentré, sans que la cellule ne dépense la moindre énergie. Mais une cellule animale n'échange pas que de l'eau : elle doit capter du \cle{glucose}, des \cle{ions} ($Na^+$, $K^+$, $Ca^{2+}$, $Cl^-$), du \cle{dioxygène}, et rejeter du \cle{dioxyde de carbone} et divers déchets. Comment ces substances dissoutes franchissent-elles la bicouche lipidique ? Consomment-elles de l'énergie ? Peuvent-elles être concentrées par la cellule ? C'est à ces questions que répond cette section, en distinguant deux grandes catégories de mécanismes : le \cle{transport passif} et le \cle{transport actif}.

\courssub{Notion préalable : le gradient de concentration}

Avant tout raisonnement, il faut disposer d'un vocabulaire précis. Toute la logique des échanges membranaires repose sur une notion simple mais fondamentale.

\begin{definition}[Gradient de concentration]
On appelle \cle{gradient de concentration} la différence de concentration d'une substance entre deux milieux séparés par une membrane. On dit qu'une substance se déplace \emph{dans le sens du gradient} (ou « le long du gradient ») lorsqu'elle passe du milieu où elle est la plus concentrée vers le milieu où elle est la moins concentrée ; elle se déplace \emph{contre le gradient} dans le cas inverse.
\end{definition}

L'intuition est celle d'une odeur de grillade qui se répand dans une cour de quartier à Yaoundé : les molécules odorantes, animées d'un mouvement incessant et désordonné (l'\emph{agitation thermique}), se dispersent spontanément des zones où elles abondent vers les zones où elles sont rares. Nul n'a besoin de « pousser » ces molécules : leur propre agitation suffit. De même, toute substance dissoute tend spontanément à s'égaliser entre deux milieux. Le gradient de concentration est donc un \emph{moteur naturel} des échanges ; aller contre lui, en revanche, exigera un travail, donc de l'énergie.

\begin{attention}[Pour les ions, on parle de gradient électrochimique]
Un ion est sensible non seulement à la différence de concentration, mais aussi à la différence de charge électrique entre les deux faces de la membrane (la face interne de la membrane plasmique est négative par rapport à la face externe). On parle alors de \cle{gradient électrochimique}. Au niveau de la Terminale, on retiendra simplement que le transport passif d'un ion suit toujours son gradient électrochimique, sans dépense d'énergie.
\end{attention}

\courssub{Le transport passif : la diffusion simple}

\begin{definition}[Diffusion simple]
La \cle{diffusion simple} est le passage de molécules à travers la bicouche lipidique, directement, sans l'aide d'aucune protéine, dans le sens du gradient de concentration et \textbf{sans dépense d'énergie} par la cellule.
\end{definition}

Pourquoi certaines molécules traversent-elles directement la membrane ? La bicouche phospholipidique possède un cœur \emph{hydrophobe} (les queues lipidiques). Les molécules petites et apolaires ou liposolubles y sont « à l'aise » et la traversent librement : c'est le cas du \cle{dioxygène} $O_2$, du \cle{dioxyde de carbone} $CO_2$, des petites molécules liposolubles (vitamines A, D, E, K, hormones stéroïdes). En revanche, les ions et les molécules polaires de taille notable (glucose, acides aminés) sont repoussés par ce cœur hydrophobe : la diffusion simple leur est interdite.

\begin{exemplebox}[La diffusion simple au niveau des poumons]
Dans les alvéoles pulmonaires, la pression partielle en $O_2$ est plus élevée que dans le sang des capillaires : $O_2$ diffuse donc simplement de l'alvéole vers le sang. À l'inverse, le $CO_2$, plus concentré dans le sang, diffuse vers l'alvéole pour être rejeté à l'expiration. Ces échanges, vitaux à chaque seconde, ne coûtent \emph{aucune} énergie à la cellule : ils suivent passivement les gradients.
\end{exemplebox}

Une propriété quantitative importante : la vitesse de diffusion simple est \emph{proportionnelle} au gradient de concentration. Plus la différence de concentration est grande, plus le flux est rapide, et ce sans limite autre que le gradient lui-même. Nous verrons que cette linéarité est une signature expérimentale précieuse.

\courssub{Le transport passif : la diffusion facilitée}

Comment le glucose, molécule polaire, parvient-il à entrer dans un hématie (globule rouge) alors que la bicouche lipidique lui est imperméable ? L'observation montre qu'il y entre rapidement, sans énergie dépensée, mais uniquement grâce à des \emph{protéines} de la membrane. C'est la \cle{diffusion facilitée}.

\begin{definition}[Diffusion facilitée]
La \cle{diffusion facilitée} est le passage d'ions ou de molécules polaires à travers la membrane plasmique par l'intermédiaire de \textbf{protéines de transport} (protéines \cle{canaux} ou protéines \cle{transporteurs}), dans le sens du gradient de concentration (ou du gradient électrochimique pour les ions) et \textbf{sans dépense d'énergie}.
\end{definition}

On distingue deux types de protéines impliquées :
\begin{itemize}
\item les \cle{protéines canaux} : ce sont des pores hydrophiles traversant la membrane, souvent spécifiques d'un ion (canaux $Na^+$, canaux $K^+$, canaux $Cl^-$). Certains sont permanents, d'autres s'ouvrent ou se ferment en réponse à un signal (canaux « voltage-dépendants » des neurones) ;
\item les \cle{transporteurs} (ou perméases) : ils lient la molécule transportée d'un côté de la membrane, changent de conformation, puis la relâchent de l'autre côté. C'est le cas du transporteur du glucose \emph{GLUT} de l'hématie.
\end{itemize}

\begin{attention}[Erreur fréquente : la diffusion facilitée ne consomme pas d'ATP]
Beaucoup d'élèves écrivent que la diffusion facilitée « consomme de l'énergie parce qu'elle utilise des protéines ». C'est faux. La présence d'une protéine ne signifie pas dépense d'énergie : la protéine offre simplement un \emph{passage} à travers la bicouche. Le mouvement reste gouverné par le gradient, du milieu le plus concentré vers le milieu le moins concentré. La diffusion facilitée est donc un \textbf{transport passif}. Seul le transport actif hydrolyse de l'ATP.
\end{attention}

Parce que le nombre de protéines de transport dans une membrane est \emph{limité}, la diffusion facilitée présente une propriété que la diffusion simple n'a pas : la \cle{saturation}. Lorsque tous les transporteurs sont occupés, augmenter encore la concentration externe n'accélère plus le transport : la vitesse atteint un \emph{plateau} (vitesse maximale $V_{max}$). C'est exactement ce que montre l'expérience suivante.

\begin{experience}[Comparaison des cinétiques de transport : protocole et résultats]
\textbf{Protocole.} On mesure la vitesse d'entrée d'une substance dans des cellules en fonction de sa concentration extracellulaire, dans deux situations : pour une molécule liposoluble traversant par diffusion simple, et pour une molécule polaire (type glucose) traversant via des transporteurs.

\textbf{Résultats.} Pour la molécule liposoluble, la vitesse d'absorption est une droite passant par l'origine : elle croît proportionnellement à la concentration externe. Pour la molécule polaire, la vitesse croît d'abord, puis plafonne à une valeur maximale $V_{max}$ qu'aucune augmentation de concentration ne permet de dépasser.

\textbf{Interprétation.} La droite signe un passage libre à travers la bicouche (diffusion simple) ; le plateau signe l'intervention d'un nombre fini de protéines de transport, toutes occupées au-delà d'une certaine concentration.
\end{experience}

\begin{popfigure}
\begin{center}
\begin{tikzpicture}
\begin{axis}[
    width=0.85\linewidth, height=6.5cm,
    xlabel={Concentration extracellulaire (mmol.L$^{-1}$)},
    ylabel={Vitesse d'absorption (unités arbitraires)},
    xmin=0, xmax=10, ymin=0, ymax=11,
    legend style={at={(0.03,0.97)}, anchor=north west, font=\small},
    grid=major, grid style={popline!40},
]
\addplot[popcurve, popblue, domain=0:10, samples=100]{x};
\addlegendentry{Diffusion simple (droite : vitesse proportionnelle au gradient)}
\addplot[popcurve, poporange, domain=0:10, samples=200]{10*x/(2+x)};
\addlegendentry{Transport via protéines (plateau de saturation)}
\draw[dashed, popmuted] (axis cs:0,10) -- (axis cs:10,10) node[pos=0.55, above, font=\small\itshape, popdark]{$V_{max}$ : tous les transporteurs sont saturés};
\end{axis}
\end{tikzpicture}
\end{center}
\legende{Vitesse d'absorption d'une substance en fonction de sa concentration extracellulaire. En diffusion simple (bleu), la vitesse croît proportionnellement au gradient, sans limite. Lorsque le transport fait intervenir des protéines (orange), la vitesse plafonne : le plateau de saturation traduit le nombre limité de transporteurs membranaires.}
\end{popfigure}

Ce critère — droite ou plateau — sera au cœur de la méthode d'analyse de courbes présentée plus bas.

\courssub{Mise en évidence expérimentale du transport actif}

\textbf{Le problème posé.} Certaines cellules accumulent des substances en quantité bien supérieure à celle du milieu extérieur. Ainsi, les cellules racinaires des plantes concentrent les ions minéraux du sol, et les cellules animales maintiennent en permanence une concentration intracellulaire en $K^+$ environ 30 fois supérieure à celle du milieu extracellulaire (environ \SI{140}{mmol.L^{-1}} à l'intérieur contre \SI{5}{mmol.L^{-1}} à l'extérieur). Ni la diffusion simple ni la diffusion facilitée ne peuvent expliquer une telle accumulation \emph{contre le gradient}. Il doit exister un troisième mécanisme.

\begin{experience}[Absorption des ions $K^+$ par des cellules : protocole et résultats]
\textbf{Protocole.} On place des cellules vivantes dans des milieux contenant des concentrations croissantes d'ions $K^+$ et l'on mesure la vitesse d'absorption de $K^+$. On renouvelle l'expérience en ajoutant du \cle{cyanure}, poison qui bloque la chaîne respiratoire mitochondriale et donc la production d'ATP, ou en plaçant les cellules en l'absence de dioxygène.

\textbf{Résultats.} En l'absence de cyanure, la vitesse d'absorption croît avec la concentration externe puis atteint un plateau de saturation. En présence de cyanure (ou sans $O_2$), l'absorption s'arrête pratiquement, alors même que le gradient de concentration est inchangé.
\end{experience}

\begin{popfigure}
\begin{center}
\begin{tikzpicture}
\begin{axis}[
    width=0.85\linewidth, height=6.5cm,
    xlabel={Concentration externe en $K^+$ (mmol.L$^{-1}$)},
    ylabel={Vitesse d'absorption de $K^+$ (unités arbitraires)},
    xmin=0, xmax=10, ymin=0, ymax=11,
    legend style={at={(0.03,0.97)}, anchor=north west, font=\small},
    grid=major, grid style={popline!40},
]
\addplot[popcurve, popgreen, domain=0:10, samples=200]{9*x/(1.5+x)};
\addlegendentry{Témoin (cellules respirant normalement)}
\addplot[popcurve, poppink, domain=0:10, samples=100]{0.4};
\addlegendentry{En présence de cyanure (ou sans $O_2$)}
\end{axis}
\end{tikzpicture}
\end{center}
\legende{Absorption des ions $K^+$ en fonction de la concentration externe. Courbe verte : cellules témoins — absorption importante avec plateau de saturation (intervention de transporteurs). Courbe rose : en présence de cyanure, inhibiteur de la respiration cellulaire, l'absorption est quasi nulle. Le transport de $K^+$ nécessite donc l'énergie issue de la respiration : c'est un transport actif.}
\end{popfigure}

\textbf{Interprétation.} Deux faits sont décisifs :
\begin{enumerate}
\item le \emph{plateau de saturation} prouve l'intervention de protéines transporteurs, en nombre limité ;
\item l'\emph{arrêt de l'absorption en présence de cyanure} (ou en anaérobiose) prouve que le transport exige l'\emph{énergie} produite par la respiration cellulaire, c'est-à-dire l'hydrolyse de l'\cle{ATP}.
\end{enumerate}

\begin{propriete}[Caractères du transport actif — établie expérimentalement]
Le \cle{transport actif} est le passage d'une substance à travers la membrane plasmique qui :
\begin{itemize}
\item peut s'effectuer \textbf{contre le gradient de concentration} (du milieu le moins concentré vers le milieu le plus concentré) ;
\item nécessite des \textbf{protéines transporteurs} spécifiques (d'où un plateau de saturation) ;
\item nécessite de l'\textbf{énergie} fournie par l'hydrolyse de l'\cle{ATP}, elle-même produite par la respiration cellulaire (d'où l'inhibition par le cyanure ou l'absence de $O_2$).
\end{itemize}
\end{propriete}

\courssub{Un exemple type : la pompe Na+/K+}

\begin{savaistu}[La pompe sodium-potassium, savoir admis utile au Bac]
La \cle{pompe $Na^+/K^+$} est une protéine transporteur présente dans la membrane de toutes les cellules animales. À chaque cycle, elle hydrolyse une molécule d'ATP pour \textbf{expulser 3 ions $Na^+$} hors de la cellule et faire \textbf{entrer 2 ions $K^+$}, dans les deux cas contre le gradient. Elle maintient ainsi la forte concentration intracellulaire en $K^+$ et la faible concentration en $Na^+$, conditions indispensables à l'excitabilité des neurones et des fibres musculaires. On estime qu'au repos, elle consomme à elle seule une fraction importante de l'ATP cellulaire.
\end{savaistu}

\begin{exemplebox}[Deux destins pour le glucose : un exemple chiffré]
Le glucose illustre parfaitement la coexistence des deux types de transport. Dans l'\emph{hématie}, la concentration en glucose est plus élevée dans le plasma (environ \SI{1}{g.L^{-1}}, soit \SI{5,5}{mmol.L^{-1}}) que dans le cytoplasme : le glucose entre par \textbf{diffusion facilitée} via le transporteur GLUT, sans énergie. En revanche, dans l'\emph{intestin grêle}, après un repas pauvre ou en fin de digestion, la concentration en glucose de la lumière intestinale peut devenir inférieure à celle du cytoplasme des cellules de l'épithélium : le glucose y est pourtant absorbé, par \textbf{transport actif} (cotransport avec $Na^+$), contre le gradient, au prix d'une dépense d'ATP. Même molécule, deux mécanismes distincts selon le sens du gradient.
\end{exemplebox}

\courssub{Schéma de synthèse : les trois modes de transport à travers la membrane}

\begin{popfigure}
\begin{center}
\begin{tikzpicture}[scale=1, every node/.style={font=\small}]
% Bicouche lipidique : deux rangées de phospholipides
\foreach \x in {0,0.6,...,9.6}{
  \draw[popblue] (\x,0.35) circle (0.14);
  \draw[popblue] (\x,-0.35) circle (0.14);
  \draw[popblue] (\x,0.21) -- (\x,0.05);
  \draw[popblue] (\x,-0.21) -- (\x,-0.05);
}
\node[popdark, anchor=west] at (10.1,0.35) {têtes hydrophiles};
\node[popdark, anchor=west] at (10.1,-0.35) {têtes hydrophiles};
\node[popdark, anchor=west] at (10.1,0) {queues hydrophobes};
\node[above, popdark] at (4.8,0.75) {\textbf{Milieu extracellulaire}};
\node[below, popdark] at (4.8,-0.75) {\textbf{Cytoplasme}};

% 1. Diffusion simple
\draw[->, very thick, popgreen] (1.2,1.6) -- (1.2,-1.6);
\node[popgreen, above] at (1.2,1.6) {$O_2$, $CO_2$};
\node[popgreen, below, align=center] at (1.2,-1.7) {\textbf{Diffusion simple}\\ (à travers la bicouche)};

% 2. Diffusion facilitée : canal
\fill[poppurple!70] (4.3,-0.55) rectangle (4.7,0.55);
\fill[white] (4.44,-0.55) rectangle (4.56,0.55);
\draw[->, very thick, poppurple] (4.5,1.6) -- (4.5,-1.6);
\node[poppurple, above] at (4.5,1.6) {ions, glucose};
\node[poppurple, below, align=center] at (4.5,-1.7) {\textbf{Diffusion facilitée}\\ (canal ou transporteur)};

% 3. Transport actif : pompe
\fill[poporange!80] (7.6,-0.6) rectangle (8.4,0.6);
\node[white, font=\footnotesize\bfseries] at (8,0) {pompe};
\draw[->, very thick, poporange] (8,-1.6) -- (8,1.6);
\node[poporange, above] at (8,1.6) {$K^+$ (contre gradient)};
\node[poporange, below, align=center] at (8,-1.7) {\textbf{Transport actif}\\ (transporteur + ATP)};
\node[popgold, font=\footnotesize\bfseries] at (7.1,-1.2) {ATP $\rightarrow$ ADP + Pi};
\end{tikzpicture}
\end{center}
\legende{Les trois modes de passage des substances dissoutes à travers la membrane plasmique. 1 : diffusion simple, directement à travers la bicouche lipidique, dans le sens du gradient. 2 : diffusion facilitée, via une protéine (canal ou transporteur), dans le sens du gradient, sans énergie. 3 : transport actif, via une pompe qui hydrolyse l'ATP, possible contre le gradient.}
\end{popfigure}

\courssub{Méthode d'analyse et tableau comparatif}

\begin{methode}[Distinguer transport passif et transport actif à partir d'une courbe ou d'un protocole]
Face à une courbe de vitesse d'absorption en fonction de la concentration externe, ou face à un protocole expérimental, procède en étapes :
\begin{enumerate}
\item \textbf{Cherche un plateau de saturation.} S'il existe, le transport fait intervenir des \emph{protéines} (diffusion facilitée ou transport actif) ; si la courbe est une droite, c'est une diffusion simple.
\item \textbf{Compare les concentrations de part et d'autre de la membrane.} Si la substance s'accumule \emph{contre le gradient}, c'est nécessairement un transport actif.
\item \textbf{Cherche l'effet d'un inhibiteur de la respiration} (cyanure, absence de $O_2$, refroidissement bloquant le métabolisme). Si l'absorption s'effondre, le transport dépend de l'ATP : c'est un transport \emph{actif}. Si l'absorption persiste, c'est un transport \emph{passif}.
\item \textbf{Conclus} en citant les deux indices expérimentaux qui fondent ta réponse.
\end{enumerate}
\end{methode}

\begin{center}
\begin{tabularx}{\linewidth}{|l|X|X|X|}
\hline
\rowcolor{popblueL} \textbf{Critère} & \textbf{Diffusion simple} & \textbf{Diffusion facilitée} & \textbf{Transport actif} \\
\hline
Sens du transport & Dans le sens du gradient & Dans le sens du gradient & Possible contre le gradient \\
\hline
Protéine membranaire & Aucune & Oui (canal ou transporteur) & Oui (pompe / transporteur) \\
\hline
Dépense d'énergie (ATP) & Non & Non & Oui (hydrolyse de l'ATP) \\
\hline
Plateau de saturation & Non (droite) & Oui & Oui \\
\hline
Inhibé par le cyanure & Non & Non & Oui \\
\hline
Exemples & $O_2$, $CO_2$, molécules liposolubles & Glucose dans l'hématie (GLUT), ions via canaux & Pompe $Na^+/K^+$, absorption du glucose intestinal, $K^+$ racinaire \\
\hline
\end{tabularx}
\end{center}

\begin{exoresolu}[Analyse de résultats expérimentaux]
\textbf{Énoncé.} On mesure l'absorption d'un ion X par des cellules en culture. Expérience 1 : la vitesse d'absorption croît avec la concentration externe puis se stabilise à un plateau. Expérience 2 : en ajoutant du cyanure au milieu, l'absorption chute de \SI{95}{\%}. Expérience 3 : la concentration intracellulaire en X devient, au bout d'une heure, 40 fois supérieure à la concentration externe. Préciser, en le justifiant, le mode de transport de l'ion X.
\tcblower
\textbf{Solution détaillée.}
\begin{itemize}
\item L'expérience 1 montre un \emph{plateau de saturation} : le transport fait donc intervenir des protéines transporteurs en nombre limité. Ce n'est pas une diffusion simple.
\item L'expérience 2 montre que l'absorption dépend de la respiration cellulaire : le cyanure bloque la production d'ATP, et l'absorption s'arrête. Le transport nécessite donc de l'énergie.
\item L'expérience 3 montre que X s'accumule \emph{contre le gradient de concentration} (40 fois plus concentré à l'intérieur), ce qu'aucun transport passif ne peut réaliser.
\end{itemize}
\textbf{Conclusion :} l'ion X est absorbé par \textbf{transport actif}, mécanisme qui met en jeu des protéines transporteurs et l'énergie de l'ATP pour déplacer la substance contre son gradient de concentration.
\end{exoresolu}

\begin{exercice}[À toi de jouer]
On étudie l'entrée d'une molécule Y dans des cellules animales. On observe que : (a) la vitesse d'entrée est strictement proportionnelle à la concentration externe de Y ; (b) l'adjonction de cyanure ne modifie pas cette vitesse ; (c) Y est une petite molécule liposoluble. Identifier le mode de transport de Y et justifier chaque étape du raisonnement.
\end{exercice}

\begin{corrige}[Exercice : mode de transport de Y]
\begin{itemize}
\item La proportionnalité entre vitesse et concentration (a), sans plateau de saturation, indique qu'aucune protéine de transport n'intervient.
\item L'insensibilité au cyanure (b) confirme l'absence de dépense d'ATP : transport passif.
\item Le caractère liposoluble de Y (c) explique qu'elle traverse directement le cœur hydrophobe de la bicouche.
\end{itemize}
\textbf{Conclusion :} Y entre dans la cellule par \textbf{diffusion simple}, dans le sens du gradient de concentration, sans protéine et sans énergie.
\end{corrige}

\begin{aretenir}[L'essentiel de la section]
\begin{itemize}
\item Le \cle{gradient de concentration} est la différence de concentration d'une substance entre deux milieux séparés par une membrane (pour les ions, on parle de gradient électrochimique).
\item Le \cle{transport passif} s'effectue dans le sens du gradient, sans dépense d'énergie : \emph{diffusion simple} (directement à travers la bicouche : $O_2$, $CO_2$) et \emph{diffusion facilitée} (via canaux ou transporteurs : glucose, ions).
\item Le \cle{transport actif} peut s'effectuer contre le gradient ; il exige des protéines transporteurs et l'énergie de l'ATP (inhibé par le cyanure ou l'absence de $O_2$). Exemple type : la pompe $Na^+/K^+$ (3 $Na^+$ sortent, 2 $K^+$ entrent par ATP hydrolysé).
\item Signatures expérimentales : plateau de saturation = protéines ; inhibition par le cyanure = transport actif.
\end{itemize}
\end{aretenir}

Ces mécanismes suffisent pour les molécules et les ions. Mais comment la cellule peut-elle absorber ou rejeter des \emph{particules} volumineuses — bactéries, gouttelettes, macromolécules — qu'aucun transporteur ne peut faire traverser la membrane ? C'est l'objet de la section suivante : l'endocytose et l'exocytose.

\coursec{Les échanges de particules : endocytose et exocytose}

Nous venons d'étudier comment la cellule échange des \cle{substances dissoutes} : petites molécules et ions traversent la membrane plasmique par diffusion simple, diffusion facilitée ou transport actif. Mais une question demeure : comment un globule blanc peut-il \emph{ingérer une bactérie entière}, dont le diamètre atteint \SI{1}{\micro\meter} ? Comment une cellule pancréatique peut-elle \emph{libérer} de grosses molécules d'insuline (masse molaire \SI{5800}{\g\per\mol}) ? Aucun canal, aucun transporteur ne peut laisser passer de tels objets : la membrane plasmique est une barrière sélective qui n'autorise le franchissement direct que de petites molécules. La cellule dispose alors d'un tout autre stratagème : elle \textbf{déforme sa membrane} pour emballer les particules dans des vésicules. C'est l'objet de cette section.

\courssub{Les limites des transports membranaires}

Les mécanismes étudiés précédemment présentent une limite de taille stricte :

\begin{itemize}
\item la \cle{diffusion simple} ne concerne que de petites molécules liposolubles ou très petites (\ce{O2}, \ce{CO2}, éthanol, urée) ;
\item la \cle{diffusion facilitée} et le \cle{transport actif} font appel à des protéines membranaires (canaux, transporteurs, pompes) dont le site de liaison n'accueille qu'une molécule ou un ion de petite taille (glucose, acides aminés, \ce{Na+}, \ce{K+}\dots) ;
\item les \cle{macromolécules} (protéines, polysaccharides) et les \cle{particules} (bactéries, débris cellulaires, gouttelettes lipidiques) sont beaucoup trop volumineuses : leur diamètre dépasse largement celui de tout canal membranaire.
\end{itemize}

\begin{aretenir}[Idée directrice]
Pour échanger des objets volumineux, la cellule n'ouvre pas de « porte » dans sa membrane : elle utilise la \textbf{plasticité} de celle-ci. La membrane peut s'\emph{invaginer} pour englober une particule extérieure (\cle{endocytose}) ou une vésicule interne peut \emph{fusionner} avec elle pour libérer son contenu (\cle{exocytose}). Dans les deux cas, l'objet transporté ne traverse jamais directement la bicouche lipidique : il reste toujours entouré d'une membrane.
\end{aretenir}

\courssub{L'endocytose : l'entrée des particules dans la cellule}

\begin{definition}[Endocytose]
L'\cle{endocytose} est le mécanisme d'\textbf{entrée de particules ou de macromolécules dans la cellule} : la membrane plasmique s'\textbf{invagine} (se creuse vers l'intérieur) autour de la particule, puis se referme et se détache en formant une \textbf{vésicule intracytoplasmique} qui internalise le matériel capturé.
\end{definition}

On distingue deux formes d'endocytose, selon la nature du matériel ingéré :

\begin{itemize}
\item la \cle{phagocytose} (du grec \emph{phagein}, manger) : la cellule ingère des \textbf{particules solides} (bactéries, débris cellulaires, poussières). Elle émet des expansions cytoplasmiques appelées \cle{pseudopodes} qui entourent la particule ; la vésicule formée, ou \cle{phagosome}, est volumineuse (jusqu'à quelques micromètres). C'est le « repas » de la cellule ;
\item la \cle{pinocytose} (du grec \emph{pinein}, boire) : la cellule ingère des \textbf{gouttelettes de liquide} extracellulaire contenant des substances dissoutes. L'invagination est beaucoup plus fine et forme de petites vésicules. C'est la « boisson » de la cellule.
\end{itemize}

\begin{attention}[Erreur fréquente : confondre pinocytose et phagocytose]
Au Baccalauréat, beaucoup de copies mélangent les deux termes. Retiens le moyen mnémotechnique : \textbf{phagocytose = manger} (particules \emph{solides} : bactérie, débris) ; \textbf{pinocytose = boire} (\emph{liquide} extracellulaire en gouttelettes). Dire qu'un leucocyte « pinocyte une bactérie » est une erreur scientifique : il la \emph{phagocyte}.
\end{attention}

\begin{experience}[Observation de la phagocytose en microscopie électronique]
\textbf{Matériel et protocole.} On place des leucocytes (globules blancs) de sang humain en présence de bactéries sur un milieu de culture à \SI{37}{\celsius}. À intervalles de temps réguliers, on fixe les cellules et on réalise des coupes fines observées au microscope électronique à transmission.

\textbf{Observations.} Sur les clichés successifs, on observe :
\begin{enumerate}
\item le leucocyte émet des expansions cytoplasmiques (pseudopodes) en direction de la bactérie ;
\item les pseudopodes s'allongent et entourent progressivement la bactérie ;
\item les extrémités des pseudopodes fusionnent : la bactérie se retrouve enfermée dans une vésicule intracytoplasmique, le phagosome ;
\item un lysosome fusionne avec le phagosome, formant une \cle{vésicule digestive} (phagolysosome) ; la bactérie est alors dégradée par les enzymes hydrolytiques des lysosomes.
\end{enumerate}

\textbf{Interprétation.} La membrane plasmique du leucocyte n'est jamais franchie directement par la bactérie : c'est la membrane qui se déforme activement pour l'englober. La présence de nombreuses mitochondries dans le cytoplasme du leucocyte suggère une forte production d'ATP, nécessaire aux mouvements du cytosquelette qui animent les pseudopodes.

\textbf{Conclusion.} La phagocytose est une endocytose de particules solides, réalisée par déformation active de la membrane plasmique, suivie d'une digestion intracellulaire assurée par les lysosomes.
\end{experience}

\begin{popfigure}
\begin{tikzpicture}[scale=0.92, every node/.style={transform shape}]
% Étape 1
\begin{scope}[shift={(0,0)}]
\fill[popblueL] (0,0) circle (1.15);
\draw[popblue, thick] (0,0) circle (1.15);
\fill[poppurpleL] (0.25,0.35) circle (0.32);
\draw[poppurple] (0.25,0.35) circle (0.32);
\fill[poporange] (2.1,0.15) ellipse (0.42 and 0.2);
\draw[poporange!70!black] (2.1,0.15) ellipse (0.42 and 0.2);
\node[font=\small\bfseries, popdark] at (0,-1.7) {1. Approche};
\end{scope}
% Étape 2
\begin{scope}[shift={(5.4,0)}]
\fill[popblueL] (0,0) circle (1.15);
\draw[popblue, thick] (0,0) circle (1.15);
\fill[poppurpleL] (-0.3,0.3) circle (0.3);
\draw[poppurple] (-0.3,0.3) circle (0.3);
\draw[popblue, thick, fill=popblueL] (0.9,0.55) .. controls (1.7,0.75) and (2.2,0.6) .. (2.35,0.25);
\draw[popblue, thick, fill=popblueL] (0.95,-0.5) .. controls (1.7,-0.7) and (2.2,-0.5) .. (2.35,-0.1);
\fill[poporange] (1.95,0.08) ellipse (0.4 and 0.19);
\draw[poporange!70!black] (1.95,0.08) ellipse (0.4 and 0.19);
\node[font=\small\bfseries, popdark] at (0.6,-1.7) {2. Pseudopodes};
\end{scope}
% Étape 3
\begin{scope}[shift={(0,-4.6)}]
\fill[popblueL] (0,0) circle (1.15);
\draw[popblue, thick] (0,0) circle (1.15);
\fill[poppurpleL] (-0.35,0.35) circle (0.3);
\draw[poppurple] (-0.35,0.35) circle (0.3);
\fill[popcream] (0.35,-0.25) circle (0.42);
\draw[popblue, thick] (0.35,-0.25) circle (0.42);
\fill[poporange] (0.35,-0.25) ellipse (0.28 and 0.13);
\draw[poporange!70!black] (0.35,-0.25) ellipse (0.28 and 0.13);
\fill[popgreen] (-0.55,-0.55) circle (0.18);
\draw[popgreen!60!black] (-0.55,-0.55) circle (0.18);
\node[font=\small\bfseries, popdark] at (0,-1.7) {3. Phagosome};
\end{scope}
% Étape 4
\begin{scope}[shift={(5.4,-4.6)}]
\fill[popblueL] (0,0) circle (1.15);
\draw[popblue, thick] (0,0) circle (1.15);
\fill[poppurpleL] (-0.4,0.4) circle (0.28);
\draw[poppurple] (-0.4,0.4) circle (0.28);
\fill[popgreenL] (0.3,-0.2) circle (0.5);
\draw[popgreen!60!black, thick] (0.3,-0.2) circle (0.5);
\fill[poporange!50] (0.18,-0.12) circle (0.06);
\fill[poporange!50] (0.42,-0.3) circle (0.05);
\fill[poporange!50] (0.3,-0.05) circle (0.04);
\node[font=\small\bfseries, popdark] at (0,-1.7) {4. Digestion};
\end{scope}
% flèches entre étapes
\draw[-{Stealth[length=3mm]}, popdark, thick] (2.9,-0.6) -- (3.6,-0.6);
\draw[-{Stealth[length=3mm]}, popdark, thick] (6,-1.5) -- (1.2,-3.2);
\draw[-{Stealth[length=3mm]}, popdark, thick] (2.9,-5.2) -- (3.6,-5.2);
\end{tikzpicture}
\legende{Schéma de la phagocytose d'une bactérie (orange) par un leucocyte. 1 : la bactérie est proche de la cellule. 2 : la membrane plasmique (bleu) émet des pseudopodes qui entourent la bactérie. 3 : la fusion des pseudopodes isole la bactérie dans un phagosome (crème) ; un lysosome (vert) s'approche. 4 : le lysosome fusionne avec le phagosome ; les enzymes digestives dégradent la bactérie (débris orange). En violet : le noyau.}
\end{popfigure}

\begin{propriete}[Caractère actif de l'endocytose et de l'exocytose — admise]
L'endocytose et l'exocytose sont des \textbf{processus actifs} : ils \textbf{consomment de l'ATP}. L'énergie fournie par l'hydrolyse de l'ATP sert à remodeler le cytosquelette et les membranes (formation des pseudopodes, invagination, déplacement et fusion des vésicules). En présence d'un inhibiteur de la respiration cellulaire (par exemple le cyanure, qui bloque la chaîne respiratoire mitochondriale et la production d'ATP), la phagocytose des leucocytes s'arrête : c'est la preuve expérimentale de cette dépendance énergétique.
\end{propriete}

\courssub{L'exocytose : la sortie des substances}

\begin{definition}[Exocytose]
L'\cle{exocytose} est le mécanisme d'\textbf{expulsion de substances hors de la cellule} : une \textbf{vésicule cytoplasmique} contenant ces substances migre vers la membrane plasmique, \textbf{fusionne} avec elle, et son contenu est \textbf{libéré dans le milieu extracellulaire}.
\end{definition}

L'exocytose est le mode de \cle{sécrétion} de toutes les grosses molécules exportées par la cellule :

\begin{itemize}
\item les \textbf{enzymes digestives} sécrétées par les cellules du pancréas exocrine ou de l'estomac (amylase, pepsinogène) ;
\item les \textbf{hormones} de nature protéique, comme l'insuline sécrétée par les cellules $\beta$ du pancréas endocrine en réponse à une élévation de la glycémie ;
\item le \textbf{mucus} sécrété par les cellules caliciformes des voies respiratoires et digestives ;
\item les \textbf{anticorps} sécrétés par les plasmocytes au cours de la réaction immunitaire.
\end{itemize}

\begin{experience}[Lien avec les organites : la filière de sécrétion]
\textbf{Document.} Des cellules pancréatiques sont incubées avec des acides aminés radioactifs (marqués au tritium). On suit la radioactivité par autoradiographie sur coupes fines au microscope électronique, à différents temps.

\textbf{Observations.} La radioactivité apparaît successivement : d'abord dans le \textbf{réticulum endoplasmique granuleux} (REG) au bout de quelques minutes, puis dans l'\textbf{appareil de Golgi}, ensuite dans des \textbf{vésicules de sécrétion} proches de la membrane plasmique, et enfin dans le \textbf{milieu extracellulaire}.

\textbf{Interprétation et conclusion.} Les protéines destinées à être sécrétées sont synthétisées par les ribosomes du REG, acheminées jusqu'à l'appareil de Golgi qui les modifie, les trie et les conditionne dans des vésicules ; ces vésicules migrent vers la membrane plasmique et libèrent leur contenu par exocytose. La sécrétion est donc le résultat d'une \textbf{coopération entre organites} : REG $\rightarrow$ Golgi $\rightarrow$ vésicules $\rightarrow$ exocytose.
\end{experience}

\begin{popfigure}
\begin{tikzpicture}[scale=0.95, every node/.style={transform shape}]
% Étape 1 : Golgi + vésicule
\begin{scope}[shift={(0,0)}]
\foreach \y in {0.5,0.15,-0.2,-0.55}{
  \draw[poppurple, thick] (-1.2,\y) .. controls (-0.4,\y+0.18) and (0.4,\y-0.18) .. (1.2,\y);
}
\fill[popgoldL] (1.9,-0.1) circle (0.4);
\draw[popgold!70!black, thick] (1.9,-0.1) circle (0.4);
\fill[poporange] (1.9,-0.1) circle (0.22);
\draw[popdark, thick] (-2.2,-1.5) -- (2.6,-1.5);
\node[font=\small\bfseries, popdark] at (0.2,-2.05) {1. Bourgeonnement de la vésicule};
\end{scope}
% Étape 2 : migration
\begin{scope}[shift={(7.2,0)}]
\draw[poppurple, thick] (-1.6,0.5) .. controls (-0.9,0.68) and (-0.2,0.32) .. (0.5,0.5);
\fill[popgoldL] (0.4,-0.9) circle (0.4);
\draw[popgold!70!black, thick] (0.4,-0.9) circle (0.4);
\fill[poporange] (0.4,-0.9) circle (0.22);
\draw[popdark, thick] (-2.2,-1.5) -- (2.6,-1.5);
\draw[-{Stealth[length=2.5mm]}, popgreen!60!black, thick] (0.4,-0.55) -- (0.4,-1.35);
\node[font=\small\bfseries, popdark] at (0.2,-2.05) {2. Migration vers la membrane};
\end{scope}
% Étape 3 : fusion et libération
\begin{scope}[shift={(14.4,0)}]
\draw[popdark, thick] (-2.2,-1.5) -- (-0.45,-1.5);
\draw[popdark, thick] (1.25,-1.5) -- (2.6,-1.5);
\draw[popgold!70!black, thick] (-0.45,-1.5) .. controls (-0.5,-0.9) and (1.3,-0.9) .. (1.25,-1.5);
\fill[poporange] (0.4,-1.05) circle (0.12);
\fill[poporange] (0.1,-0.55) circle (0.08);
\fill[poporange] (0.7,-0.45) circle (0.08);
\fill[poporange] (0.4,-0.1) circle (0.06);
\draw[-{Stealth[length=2.5mm]}, poporange!70!black, thick] (0.4,-1.3) -- (0.4,0.35);
\node[font=\small\bfseries, popdark] at (0.2,-2.05) {3. Fusion et libération};
\end{scope}
% flèches entre étapes
\draw[-{Stealth[length=3mm]}, popdark, thick] (3.1,-0.8) -- (4.4,-0.8);
\draw[-{Stealth[length=3mm]}, popdark, thick] (10.3,-0.8) -- (11.6,-0.8);
\end{tikzpicture}
\legende{Schéma de l'exocytose. 1 : une vésicule de sécrétion (jaune) chargée de protéines (orange) bourgeonne à partir de l'appareil de Golgi (violet). 2 : la vésicule migre vers la membrane plasmique (trait noir). 3 : la membrane de la vésicule fusionne avec la membrane plasmique ; le contenu est libéré dans le milieu extracellulaire. Ce processus consomme de l'ATP.}
\end{popfigure}

\begin{attention}[Pièges de rédaction]
\begin{itemize}
\item Ne dis jamais qu'une particule « traverse la membrane » lors de l'endocytose : elle est \emph{englobée} par la membrane, qui se referme en vésicule. La bicouche lipidique n'est jamais franchie directement.
\item Ne confonds pas les sens : \textbf{endo}cytose = entrée ; \textbf{exo}cytose = sortie.
\item N'oublie pas de préciser que ces deux mécanismes \textbf{consomment de l'ATP} : c'est ce qui les distingue des transports passifs.
\end{itemize}
\end{attention}

\courssub{Importance des échanges cellulaires dans la vie des animaux}

Les mécanismes d'endocytose et d'exocytose sont indispensables aux fonctions vitales des animaux :

\begin{itemize}
\item \textbf{Nutrition et défense immunitaire :} la phagocytose permet aux macrophages et neutrophiles d'éliminer les agents pathogènes (bactéries, champignons) et les débris cellulaires ; la pinocytose assure la prise de nutriments dissous chez de nombreux types cellulaires. Dans les centres de santé du Cameroun, une hyperleucocytose (élévation du nombre de globules blancs) signe souvent une infection bactérienne en cours, témoin de l'activité phagocytaire intense.
\item \textbf{Communication intercellulaire :} l'exocytose assure la sécrétion des hormones (insuline, glucagon, ADH), des neurotransmetteurs (acétylcholine, dopamine) et des facteurs de croissance, coordonnant ainsi le fonctionnement de l'organisme.
\item \textbf{Digestion extracellulaire :} la sécrétion d'enzymes digestives (amylase, lipase, trypsine) par le pancréas exocrine et l'estomac dépend de l'exocytose.
\item \textbf{Protection des surfaces :} la sécrétion de mucus par exocytose (cellules caliciformes) protège les épithéliums respiratoires et digestifs.
\item \textbf{Renouvellement de la membrane plasmique :} l'endocytose et l'exocytose équilibrent la surface membranaire ; lors de la sécrétion, la membrane vésiculaire s'intègre à la membrane plasmique (exocytose), et l'endocytose permet de la recycler.
\end{itemize}

\begin{exoresolu}[Analyser un cliché de microscopie électronique]
\textbf{Énoncé.} Sur un cliché de microscopie électronique à transmission, on observe un globule blanc dont le cytoplasme émet deux expansions qui entourent une bactérie ; par ailleurs, une vésicule contenant des débris bactériens est visible à proximité d'un lysosome. (1) Nommer le phénomène observé. (2) Décrire en quelques étapes ce qui s'est passé et ce qui va se produire. (3) Expliquer pourquoi ce phénomène est bloqué si l'on ajoute du cyanure au milieu.
\tcblower
\textbf{Solution.} (1) Il s'agit de la \textbf{phagocytose}, une forme d'endocytose de particules solides.

(2) Le leucocyte a émis des \textbf{pseudopodes} qui ont entouré la bactérie ; la fusion des pseudopodes a isolé la bactérie dans une vésicule appelée \textbf{phagosome}. Ensuite, le \textbf{lysosome} va fusionner avec le phagosome pour former une vésicule digestive (phagolysosome) : les enzymes lysosomiales dégraderont la bactérie en débris, qui seront ensuite éliminés ou recyclés.

(3) Le cyanure bloque la chaîne respiratoire mitochondriale (complexe IV) et donc la \textbf{production d'ATP}. Or la phagocytose est un processus actif : la formation des pseudopodes et les mouvements membranaires exigent de l'énergie. Sans ATP, les pseudopodes ne peuvent ni se former ni se déplacer : la phagocytose s'arrête.
\end{exoresolu}

\begin{savaistu}[Des macrophages infatigables]
Un \cle{macrophage} (grosse cellule phagocytaire de nos tissus) peut ingérer et détruire \textbf{plusieurs dizaines de bactéries} au cours de sa vie, parfois plus d'une centaine. Les macrophages alvéolaires des poumons nettoient en permanence l'air inspiré de ses poussières et microbes ; ceux de la rate (macrophages spléniques) recyclent les globules rouges vieillis. La phagocytose est ainsi un \textbf{pilier de la défense de l'organisme} : lors d'une infection, le pus qui se forme dans une plaie est en grande partie constitué de leucocytes morts après avoir phagocyté des bactéries.
\end{savaistu}

\begin{exercice}[Pour s'entraîner]
Une cellule du pancréas sécrète de l'insuline, une hormone protéique. (1) Expliquer pourquoi l'insuline ne peut pas être exportée par un transporteur membranaire. (2) Décrire le trajet complet de l'insuline, de sa synthèse à sa libération dans le sang, en citant les organites impliqués. (3) Prévoir l'effet d'un poison bloquant la production d'ATP sur cette sécrétion, en justifiant.
\end{exercice}

\begin{corrige}[Exercice : sécrétion de l'insuline]
(1) L'insuline est une \textbf{macromolécule} (protéine de masse molaire \SI{5800}{\g\per\mol}), beaucoup trop volumineuse pour le site de liaison d'un transporteur ou le pore d'un canal : les transports membranaires classiques ne concernent que de petites molécules et des ions.

(2) L'insuline est synthétisée par les \textbf{ribosomes du réticulum endoplasmique granuleux} (REG), puis acheminée vers l'\textbf{appareil de Golgi} qui la modifie, la trie et la conditionne dans des \textbf{vésicules de sécrétion}. Ces vésicules migrent vers la membrane plasmique, \textbf{fusionnent} avec elle et libèrent l'insuline dans le milieu extracellulaire par \textbf{exocytose} ; l'hormone gagne ensuite le sang.

(3) Un poison bloquant la production d'ATP \textbf{arrête la sécrétion} : l'exocytose (migration et fusion des vésicules) est un processus actif consommateur d'ATP. Sans énergie, les vésicules s'accumulent dans le cytoplasme et l'insuline n'est plus libérée.
\end{corrige}

\begin{aretenir}[L'essentiel de la section]
\begin{itemize}
\item Les grosses particules et macromolécules ne franchissent jamais directement la membrane : la cellule utilise la \textbf{déformation membranaire}.
\item \textbf{Endocytose} = entrée par invagination de la membrane formant une vésicule ; deux formes : \textbf{phagocytose} (solides, pseudopodes, phagosome puis digestion lysosomiale) et \textbf{pinocytose} (gouttelettes liquides).
\item \textbf{Exocytose} = sortie par fusion d'une vésicule avec la membrane plasmique ; c'est le mode de sécrétion des enzymes, hormones, mucus, anticorps.
\item Filière de sécrétion : REG $\rightarrow$ appareil de Golgi $\rightarrow$ vésicules $\rightarrow$ exocytose ; digestion après phagocytose : phagosome + lysosome $\rightarrow$ phagolysosome.
\item Endocytose et exocytose sont des \textbf{processus actifs consommateurs d'ATP}.
\item Ces échanges assurent nutrition, défense, communication, digestion, protection et renouvellement membranaire chez l'animal.
\end{itemize}
\end{aretenir}

\coursec{Importance des échanges cellulaires dans la vie des animaux}

Dans la section précédente, nous avons vu comment la cellule animale fait entrer ou sortir de grosses particules grâce à l'\cle{endocytose} et à l'\cle{exocytose}. Nous disposons désormais de la panoplie complète des mécanismes d'échanges : osmose et dialyse pour l'eau et les solutés, diffusion simple, diffusion facilitée et transport actif pour les substances dissoutes, endocytose et exocytose pour les particules. Il reste à répondre à la question essentielle : \emph{à quoi servent tous ces mécanismes dans la vie quotidienne d'un animal ?} Tu vas le découvrir : chaque grande fonction de l'organisme — se nourrir, respirer, éliminer, se défendre — n'est, au fond, qu'une application concertée des échanges cellulaires que tu viens d'étudier.

\begin{propriete}[Principe fondamental]
La vie de toute cellule animale dépend d'\cle{échanges permanents et régulés} avec le \cle{milieu intérieur} (sang, lymphe, liquide interstitiel) qui l'entoure. Ces échanges assurent l'approvisionnement en matière et en énergie, l'élimination des déchets et le maintien de la stabilité du milieu intérieur (\cle{homéostasie}). Toute perturbation durable de ces échanges entraîne la souffrance, puis la mort de la cellule.
\end{propriete}

\courssub{La nutrition : l'absorption des nutriments au niveau intestinal}

\textbf{Observation et problème.}\par
Après un repas de foutou et de ndolé, les aliments sont digérés dans l'intestin grêle : les glucides donnent du glucose, les protéines des acides aminés. Or ces nutriments se trouvent \emph{dans le tube digestif}, c'est-à-dire encore à l'extérieur du corps au sens strict. Comment traversent-ils la paroi intestinale pour rejoindre le sang ?

\textbf{Données expérimentales.}\par
Des mesures réalisées sur des anses intestinales de mammifères montrent que l'absorption du glucose se poursuit même lorsque la concentration en glucose est \emph{plus faible} dans la lumière intestinale que dans le sang : le glucose est absorbé \emph{contre son gradient de concentration}. En revanche, cette absorption s'arrête si l'on bloque la respiration cellulaire (par du cyanure, poison des mitochondries) ou si l'on supprime le sodium du milieu.

\textbf{Interprétation.}\par
L'absorption contre le gradient et sa dépendance à l'énergie prouvent qu'il s'agit d'un \cle{transport actif} : les cellules de l'épithélium intestinal (entérocytes) possèdent, sur leur membrane tournée vers le sang, des pompes Na$^+$/K$^+$ ATPases qui chassent le Na$^+$ hors de la cellule, créant un gradient de sodium ; le glucose entre alors dans l'entérocyte, côté lumière intestinale, par un \emph{cotransport} Na$^+$-glucose qui exploite ce gradient. Les acides aminés sont absorbés selon le même principe. Une fois concentrés dans l'entérocyte, glucose et acides aminés sortent vers le sang par \cle{diffusion facilitée} (via des transporteurs spécifiques), c'est-à-dire dans le sens de leur gradient, sans dépense d'énergie. L'eau suit ensuite par \cle{osmose}.

\begin{aretenir}[Absorption intestinale]
L'absorption des nutriments combine plusieurs mécanismes : \textbf{transport actif} (entrée du glucose et des acides aminés dans l'entérocyte, contre le gradient, avec dépense d'ATP), \textbf{diffusion facilitée} (sortie vers le sang, dans le sens du gradient, sans ATP) et \textbf{osmose} (l'eau suit les solutés absorbés).
\end{aretenir}

\courssub{La respiration : les échanges de O$_2$ et CO$_2$}

\textbf{Observation.}\par
Chaque inspiration renouvelle l'air des alvéoles pulmonaires ; chaque cellule du corps consomme du dioxygène et rejette du dioxyde de carbone (respiration cellulaire : \ce{C6H12O6 + 6O2 -> 6CO2 + 6H2O} + énergie). Comment ces gaz traversent-ils les membranes ?

\textbf{Raisonnement.}\par
O$_2$ et CO$_2$ sont de petites molécules non chargées, solubles dans les lipides : elles traversent librement les bicouches phospholipidiques par \cle{diffusion simple}, toujours \emph{dans le sens du gradient de pression partielle}.
\begin{itemize}
\item \textbf{Au niveau des alvéoles pulmonaires} : l'air alvéolaire est riche en O$_2$ (pression partielle élevée) et le sang qui arrive dans les capillaires en est pauvre ; O$_2$ diffuse donc de l'alvéole vers le sang. Inversement, le sang veineux est riche en CO$_2$ qui diffuse du sang vers l'alvéole, puis est expiré.
\item \textbf{Au niveau des tissus} : les cellules en activité consomment O$_2$ (sa pression partielle y est faible) : O$_2$ diffuse du sang vers les cellules. Le CO$_2$ produit diffuse des cellules vers le sang.
\end{itemize}
La finesse des parois alvéolaires et capillaires (une seule couche de cellules aplaties) et l'immense surface d'échange des poumons (environ \SI{70}{m^2} chez l'adulte) rendent cette diffusion extrêmement efficace.

\begin{attention}[Piège fréquent]
Les échanges gazeux respiratoires sont des diffusions \emph{simples} : ils ne nécessitent \textbf{ni transporteur, ni ATP}. Ne parle jamais de « transport actif de l'oxygène » au niveau des alvéoles ! C'est la circulation sanguine (travail du cœur) et la ventilation pulmonaire qui maintiennent les gradients, pas la membrane alvéolaire.
\end{attention}

\courssub{L'excrétion : le rôle du rein et l'application médicale de la dialyse}

\textbf{Problème.}\par
L'activité cellulaire produit des déchets toxiques, notamment l'\cle{urée} issue de la dégradation des protéines. Le rein élimine ces déchets tout en conservant l'eau et les substances utiles. Comment ?

\textbf{Mécanismes.}\par
Au niveau du \cle{néphron} (unité fonctionnelle du rein), deux étapes mettent en jeu les échanges cellulaires :
\begin{enumerate}
\item \textbf{La filtration} : sous l'effet de la pression sanguine, l'eau et les petites molécules (urée, glucose, sels) passent du sang du glomérule vers le tube urinifère, à travers une paroi qui agit comme un filtre ; les grosses protéines et les cellules sanguines restent dans le sang. C'est un passage sélectif à travers une membrane, comparable à la \cle{dialyse}.
\item \textbf{La réabsorption} : le long du tube urinifère, les substances utiles sont récupérées : le glucose est totalement réabsorbé par \cle{transport actif}, les ions Na$^+$ aussi, et l'eau suit par \cle{osmose} (environ 99 \% de l'eau filtrée est réabsorbée : sur environ \SI{180}{L} filtrés par jour, seulement \SI{1,5}{L} d'urine est excrété). L'urée, déchet, est très peu réabsorbée et part dans les urines.
\end{enumerate}

\begin{exemplebox}[La dialyse rénale artificielle]
Lorsque les reins d'un patient cessent de fonctionner (insuffisance rénale chronique), l'urée s'accumule dans le sang : c'est l'urémie, mortelle sans traitement. Le \cle{rein artificiel} exploite le principe de la \cle{dialyse} étudié avec l'expérience de Pfeiffer : le sang du patient circule d'un côté d'une \textbf{membrane perméable} (aux petites molécules, imperméable aux protéines et aux cellules sanguines), tandis qu'un \textbf{liquide de dialyse} — dépourvu d'urée mais contenant les sels et le glucose aux bonnes concentrations — circule de l'autre côté. L'urée diffuse du sang (concentrée) vers le liquide de dialyse (pauvre), selon son gradient de concentration ; les substances utiles ne quittent pas le sang car leur concentration est identique de part et d'autre. Le sang épuré est rendu au patient. Des séances de 4 heures, trois fois par semaine, sont réalisées dans les centres d'hémodialyse des hôpitaux de Yaoundé et de Douala.
\end{exemplebox}

\begin{popfigure}
\begin{center}
\begin{tikzpicture}[font=\small, >=Stealth]
% Membrane
\fill[poppurple!25] (0,-2.2) rectangle (0.35,2.2);
\node[rotate=90] at (0.175,0) {\footnotesize membrane perméable};
% Compartiment sang
\draw[thick, poppink, rounded corners] (-4.2,-2.2) rectangle (0,2.2);
\node[poppink, font=\bfseries] at (-2.1,1.8) {Sang du patient};
\node at (-2.1,1.1) {urée \textbf{+++}};
\node at (-2.1,0.5) {sels, glucose};
\node at (-2.1,-0.1) {protéines, globules};
% Compartiment dialysat
\draw[thick, popblue, rounded corners] (0.35,-2.2) rectangle (4.5,2.2);
\node[popblue, font=\bfseries] at (2.4,1.8) {Liquide de dialyse};
\node at (2.4,1.1) {sans urée};
\node at (2.4,0.5) {sels, glucose (mêmes};
\node at (2.4,0.1) {concentrations que le sang)};
% Flèches
\draw[->, very thick, poporange] (-1.6,-1.0) -- (1.9,-1.0) node[midway, below, popdark]{\footnotesize diffusion de l'urée};
\draw[->, very thick, poppink] (-3.6,2.6) -- (-3.6,2.25);
\node[poppink, left] at (-3.6,2.7) {\footnotesize sang souillé};
\draw[->, very thick, poppink] (-0.8,-2.25) -- (-0.8,-2.6);
\node[poppink, left] at (-0.8,-2.75) {\footnotesize sang épuré};
\draw[->, very thick, popblue] (1.0,-2.6) -- (1.0,-2.25);
\node[popblue, left] at (1.0,-2.75) {\footnotesize dialysat frais};
\draw[->, very thick, popblue] (3.9,2.25) -- (3.9,2.6);
\node[popblue, right] at (3.9,2.7) {\footnotesize dialysat usé};
\end{tikzpicture}
\end{center}
\legende{Principe du rein artificiel : le sang et le liquide de dialyse sont séparés par une membrane perméable. L'urée diffuse du sang vers le dialysat (gradient de concentration) ; les protéines et les globules, trop gros, restent dans le sang.}
\end{popfigure}

\courssub{L'immunité : la phagocytose par les leucocytes}

Lorsqu'une bactérie pénètre dans l'organisme (plaie, infection), certains globules blancs (\cle{leucocytes}, notamment les polynucléaires neutrophiles et les macrophages) la capturent par \cle{phagocytose}, une forme d'\cle{endocytose} : la membrane du leucocyte émet des pseudopodes qui entourent le microbe, puis se soudent pour former une vésicule interne, le \emph{phagosome}. Celui-ci fusionne avec des lysosomes dont les enzymes digèrent le microbe. Les résidus sont ensuite rejetés hors de la cellule par \cle{exocytose}. C'est ainsi que le pus d'une plaie infectée contient des millions de leucocytes ayant phagocyté des microbes. Sans endocytose, aucune défense immunitaire cellulaire ne serait possible.

\courssub{La régulation : le maintien de l'équilibre hydrique et ionique}

Nous avons vu (section sur l'osmose) qu'un globule rouge placé dans un milieu hypotonique gonfle et éclate (hémolyse), et qu'il se ratatine (crénation) en milieu hypertonique. Pour que les cellules vivent, le milieu intérieur doit donc rester \textbf{isotonique} : sa concentration en sels et sa teneur en eau doivent être stables. Cette stabilité, appelée \cle{homéostasie}, est assurée en permanence par :
\begin{itemize}
\item les \textbf{pompes Na$^+$/K$^+$} des membranes cellulaires (transport actif), qui compensent en continu les fuites d'ions par diffusion ;
\item le \textbf{rein}, qui ajuste l'excrétion d'eau (osmose) et de sels selon les besoins, sous contrôle hormonal ;
\item l'\textbf{alimentation et l'hydratation}, qui apportent l'eau et les sels perdus (sueur, urines).
\end{itemize}
Lors d'un effort sous le soleil de Maroua, la sueur fait perdre eau et sels : la soif et la réduction des urines rétablissent l'équilibre. La vie cellulaire n'est possible que dans une gamme étroite de concentration du milieu intérieur.

\begin{savaistu}[Le choléra : quand les échanges d'eau deviennent mortels]
Le \cle{choléra}, maladie qui sévit encore par épidémies en Afrique centrale (notamment dans l'Extrême-Nord du Cameroun), est dû à la bactérie \emph{Vibrio cholerae}. Sa toxine dérègle les transports d'ions dans les cellules intestinales : celles-ci rejettent massivement des ions (Cl$^-$, Na$^+$) dans la lumière intestinale, et l'eau suit par \textbf{osmose} : jusqu'à \SI{15}{L} de liquide peuvent être perdus par jour en diarrhées ! La déshydratation tue en quelques heures. Le traitement repose sur les \textbf{sels de réhydratation orale} (SRO) : une solution d'eau, de sel et de glucose. Le glucose permet l'absorption du sodium par cotransport dans les entérocytes, et l'eau suit par osmose : le patient se réhydrate. Comprendre les échanges cellulaires sauve littéralement des vies.
\end{savaistu}

\courssub{Synthèse générale : les mécanismes d'échanges au service de la vie}

\begin{popfigure}
\begin{center}
\begin{tikzpicture}[font=\small, >=Stealth]
% Cellule centrale
\fill[popcream] (0,0) ellipse (2.3 and 1.7);
\draw[very thick, poppurple] (0,0) ellipse (2.3 and 1.7);
\fill[poppinkL] (0.2,0.1) circle (0.55);
\draw[thick, poppink] (0.2,0.1) circle (0.55);
\node at (0.2,0.1) {\footnotesize noyau};
\node[popdark] at (0,-1.25) {\footnotesize cytoplasme};
\node[font=\bfseries, poppurple] at (0,2.15) {Cellule animale};
% Flèches entrantes
\draw[->, very thick, popgreen] (-5.4,1.3) -- (-2.5,0.9) node[midway, above, sloped, popdark]{\footnotesize O$_2$ : diffusion simple};
\draw[->, very thick, popgreen] (-5.4,0.1) -- (-2.5,0.0) node[midway, above, sloped, popdark]{\footnotesize glucose : transport actif};
\draw[->, very thick, popgreen] (-5.4,-1.1) -- (-2.5,-0.8) node[midway, above, sloped, popdark]{\footnotesize eau : osmose};
% Flèches sortantes
\draw[->, very thick, poporange] (2.5,0.9) -- (5.4,1.3) node[midway, above, sloped, popdark]{\footnotesize CO$_2$ : diffusion simple};
\draw[->, very thick, poporange] (2.5,0.0) -- (5.4,0.1) node[midway, above, sloped, popdark]{\footnotesize urée : diffusion};
\draw[->, very thick, poporange] (2.5,-0.8) -- (5.4,-1.1) node[midway, above, sloped, popdark]{\footnotesize sécrétions : exocytose};
% Endocytose
\draw[->, very thick, popblue] (-1.2,-3.0) -- (-0.7,-1.75) node[midway, left, popdark, xshift=-2pt]{\footnotesize particules : endocytose};
\end{tikzpicture}
\end{center}
\legende{Schéma de synthèse : une cellule animale échange en permanence avec le milieu intérieur. Flèches vertes : entrées (O$_2$, nutriments, eau). Flèches orange : sorties (CO$_2$, déchets, sécrétions). Flèche bleue : capture de particules. Chaque flèche est associée à son mécanisme.}
\end{popfigure}

\begin{center}
\begin{tabularx}{\linewidth}{|l|X|X|X|}
\hline
\rowcolor{popblueL}
\textbf{Élément échangé} & \textbf{Mécanisme} & \textbf{Énergie (ATP)} & \textbf{Exemple physiologique} \\
\hline
Eau & Osmose & Non & Réabsorption rénale ; réhydratation intestinale \\
\hline
Gaz (O$_2$, CO$_2$) & Diffusion simple & Non & Échanges alvéoles--sang et sang--tissus \\
\hline
Glucose, acides aminés & Transport actif / diffusion facilitée & Oui (transport actif) & Absorption intestinale ; réabsorption rénale du glucose \\
\hline
Ions (Na$^+$, K$^+$) & Transport actif (pompe) / diffusion facilitée (canaux) & Oui (pompe) & Équilibre ionique ; influx nerveux \\
\hline
Particules, microbes & Endocytose (phagocytose) & Oui & Défense immunitaire par les leucocytes \\
\hline
Protéines sécrétées, résidus & Exocytose & Oui & Sécrétion d'hormones, d'enzymes digestives \\
\hline
\end{tabularx}
\end{center}

\begin{methode}[Rédiger la conclusion d'un exercice sur les échanges cellulaires]
Pour conclure rigoureusement un exercice d'interprétation, ta phrase doit toujours contenir quatre éléments :
\begin{enumerate}
\item \textbf{Nommer le mécanisme} : osmose, dialyse, diffusion simple, diffusion facilitée, transport actif, endocytose ou exocytose ;
\item \textbf{Préciser le sens du flux} : quelle substance se déplace, d'où vers où (du milieu le plus concentré vers le moins concentré pour les transports passifs ; contre le gradient pour le transport actif) ;
\item \textbf{Identifier la membrane concernée} : membrane plasmique de la cellule, paroi alvéolaire, membrane du rein artificiel, etc. ;
\item \textbf{Indiquer le besoin énergétique} : avec consommation d'ATP (transport actif, endocytose, exocytose) ou sans ATP (osmose, diffusions).
\end{enumerate}
\emph{Exemple de rédaction} : « Le glucose passe de la lumière intestinale dans l'entérocyte par transport actif (contre son gradient, avec ATP), à travers la membrane plasmique, puis rejoint le sang par diffusion facilitée. »
\end{methode}

\begin{exoresolu}[Exercice résolu : analyse d'une situation clinique]
\textbf{Énoncé.} Un patient atteint d'insuffisance rénale présente un taux d'urée sanguin de \SI{2,5}{g/L} (norme : inférieur à \SI{0,45}{g/L}). Il est branché sur un rein artificiel dont le liquide de dialyse ne contient pas d'urée mais contient du glucose à \SI{1}{g/L}, concentration identique à celle du sang. (1) Expliquer le devenir de l'urée au cours de la dialyse. (2) Pourquoi le glucose du sang n'est-il pas éliminé ? (3) Pourquoi les protéines sanguines restent-elles dans le sang ?
\tcblower
\textbf{Solution.} (1) L'urée, petite molécule, traverse la membrane perméable du rein artificiel par \textbf{diffusion} (dialyse) : elle passe du sang, où elle est très concentrée (\SI{2,5}{g/L}), vers le liquide de dialyse, qui n'en contient pas, c'est-à-dire \emph{dans le sens du gradient de concentration}, sans dépense d'ATP. Le sang est ainsi progressivement épuré. (2) Le glucose a la même concentration (\SI{1}{g/L}) des deux côtés de la membrane : il n'existe pas de gradient, donc pas de flux net ; le glucose du sang est conservé. (3) Les protéines sont de grosses molécules dont la taille dépasse le diamètre des pores de la membrane : celle-ci, comme la membrane de l'expérience de Pfeiffer, leur est \textbf{imperméable} ; elles restent donc dans le sang. \emph{Conclusion} : le rein artificiel exploite le principe de la dialyse à travers une membrane sélectivement perméable.
\end{exoresolu}

\begin{attention}[Erreurs fréquentes à éviter]
\begin{itemize}
\item Confondre \textbf{dialyse} (passage de solutés dissous à travers une membrane perméable) et \textbf{osmose} (passage de l'\emph{eau} à travers une membrane semi-perméable). Le rein artificiel est une application de la dialyse, pas de l'osmose.
\item Affirmer que l'absorption intestinale du glucose est « toujours » un transport actif : l'entrée dans l'entérocyte est active, mais sa sortie vers le sang est une diffusion facilitée.
\item Oublier que la diffusion se fait toujours du milieu le plus concentré vers le moins concentré : vérifie le sens du gradient avant de conclure.
\item Croire que la cellule n'échange que lorsqu'elle en a « besoin » : les échanges sont \emph{permanents}, même au repos, car la cellule consomme et produit en continu.
\end{itemize}
\end{attention}

\begin{aretenir}[L'essentiel de la section]
Toutes les grandes fonctions de l'organisme animal reposent sur les échanges cellulaires : la \textbf{nutrition} (transport actif et diffusion facilitée au niveau intestinal), la \textbf{respiration} (diffusion simple des gaz), l'\textbf{excrétion} (filtration, réabsorption active et osmose au niveau rénal ; dialyse artificielle en médecine), l'\textbf{immunité} (phagocytose) et la \textbf{régulation} du milieu intérieur (pompes ioniques, osmose). La cellule vit parce qu'elle échange : arrêter les échanges, c'est arrêter la vie.
\end{aretenir}

\newpage

\coursec{Activité d'intégration}

\coursec{Leçon 1 : Organisation de la cellule et échanges cellulaires — Activité d'intégration}

\courssub{Objectifs de l'activité}
Mobiliser les savoirs et savoir-faire de la leçon sur l'organisation cellulaire et les échanges cellulaires pour résoudre un problème de santé publique concret : le choix du sérum de réhydratation adapté à un enfant sévèrement déshydraté. Cette activité vise à :
\begin{itemize}
    \item Observer, décrire et schématiser le devenir d'hématies placées dans des milieux de tonicité différente.
    \item Interpréter des résultats expérimentaux (tableau, courbe, documents historiques) pour identifier le mécanisme d'osmose et en préciser le sens.
    \item Établir le lien entre la tonicité du milieu extracellulaire, le volume cellulaire et l'intégrité de la membrane plasmique.
    \item Justifier scientifiquement le choix thérapeutique du sérum physiologique (0,9 \% NaCl) et comprendre la physiologie de la réhydratation orale (couplage glucose/Na$^+$).
\end{itemize}

\courssub{Situation : Urgence au Centre de Santé de Bafoussam}
Un enfant de 18 mois, nommé \textbf{Amadou}, est admis en urgence au Centre de Santé Urbain de Bafoussam (Région de l'Ouest). Il présente une diarrhée aqueuse profuse depuis 24 heures, des vomissements, une fontanelle antérieure déprimée, une peau qui se replie lentement (pli cutané \textasciitilde 2 s), une tachycardie à 160 battements/min et une hypotension. Le diagnostic est sans équivoque : \textbf{déshydratation sévère avec hypovolémie} (perte de \textasciitilde 10 \% du poids corporel en eau et électrolytes).

L'infirmier chef de poste, \textbf{M. Kamga}, doit mettre en place une réhydratation intraveineuse immédiate (Plan C de l'OMS) avant le relais par voie orale. Trois poches de perfusion sont disponibles sur le chariot d'urgence, mais leurs étiquettes sont partiellement arrachées. Il ne reste lisible que la composition en chlorure de sodium :
\begin{enumerate}
    \item \textbf{Poche A :} Eau distillée (0 g/L NaCl).
    \item \textbf{Poche B :} Solution de NaCl à 0,9 \% (9 g/L).
    \item \textbf{Poche C :} Solution de NaCl à 5 \% (50 g/L).
\end{enumerate}

M. Kamga hésite. Il sait qu'une mauvaise choix peut être fatal : une hémolyse massive (éclatement des globules rouges) aggraverait l'anémie fréquente chez l'enfant camerounais (paludisme, drépanocytose), tandis qu'une shrinkage cellulaire (création) perturberait la microcirculation et la fonction d'organes vitaux.

Pour l'aider à décider, le biologiste du laboratoire vous remet trois documents issus d'une série d'expériences réalisées sur du sang frais prélevé sur un donneur sain (hématocrite initial = 45~\%, volume corpusculaire moyen VGM = 90 fL).

\medskip
\noindent\textbf{Document 1 : Résultats de l'expérience sur hématies isolées (observation au microscope optique $\times$ 400 et mesure de volume par hématologie automate).}

\begin{center}
\begin{tabularx}{\linewidth}{|l|X|X|X|}\hline
\rowcolor{popblueL}
\textbf{Milieu d'incubation} & \textbf{Aspect microscopique (schéma simplifié)} & \textbf{Volume moyen mesuré (fL)} & \textbf{État de la cellule} \\ \hline
\textbf{Milieu 1 : Eau distillée} & \textit{Hématies « fantômes » : membranes vides, circulaires, très claires, pas de hémoglobine visible.} & \textbf{> 150 fL} (éclatement lors de la mesure) & \textbf{Hémolyse totale} \\ \hline
\textbf{Milieu 2 : NaCl 0,9 \%} & \textit{Disques biconcaves classiques, bien individués, couleur rouge uniforme.} & \textbf{90 fL} (identique au témoin) & \textbf{Intactes (normales)} \\ \hline
\textbf{Milieu 3 : NaCl 5 \%} & \textit{Formes échinocytaires (bords dentelés, « boules à piquants »), volume réduit, couleur rouge foncé concentré.} & \textbf{55 fL} & \textbf{Création (crenation)} \\ \hline
\end{tabularx}
\end{center}

\medskip
\noindent\textbf{Document 2 : Courbe de variation du volume des hématies en fonction de la concentration en NaCl.}
\begin{popfigure}
\begin{tikzpicture}
\begin{axis}[
    width=\linewidth, height=8cm,
    axis lines=middle,
    xlabel={Concentration en NaCl (\% massique)}, ylabel={Volume relatif des hématies (\% du VGM)},
    xmin=0, xmax=6, ymin=0, ymax=180,
    xtick={0,0.9,5}, xticklabels={0, 0,9, 5},
    ytick={0,50,100,150}, yticklabels={0, 50, 100, 150},
    domain=0:6, samples=200,
    clip=false,
    tick style={popdark},
    xlabel style={anchor=north, yshift=-5pt},
    ylabel style={anchor=east, xshift=-5pt},
    grid=major, grid style={popline},
]
% Courbe théorique : pic à 0.9, chute brutale à gauche (lyse), décroissance à droite
\addplot[popcurve, very thick, domain=0:0.85] {100 * exp(-((x-0.9)/0.3)^2) + 20*(0.85-x)/0.85};
\addplot[popcurve, very thick, domain=0.85:0.95] {100 * exp(-((x-0.9)/0.3)^2)};
\addplot[popcurve, very thick, domain=0.95:6] {100 * exp(-(x-0.9)/1.2)};
% Ligne pointillée isotonicité
\draw[popmuted, dashed] (axis cs:0.9,0) -- (axis cs:0.9,100) node[above, pos=0.9, popdark, font=\small] {Isotonicité};
\draw[popmuted, dashed] (axis cs:0,100) -- (axis cs:0.9,100);
% Annotations
\node[popdark, font=\small, align=center] at (axis cs:0.3,160) {\textbf{Domaine hypotonique}\\ Hémolyse};
\node[popdark, font=\small, align=center] at (axis cs:3.5,40) {\textbf{Domaine hypertonique}\\ Création};
\node[popdark, font=\small, align=center] at (axis cs:0.9,115) {\textbf{Domaine isotonicité}\\ Intégrité cellulaire};
\end{axis}
\end{tikzpicture}
\legende{Variation du volume relatif des hématies humaines en fonction de la concentration massique en NaCl du milieu extracellulaire. Le maximum (100 \%) correspond à l'isotonicité (0,9 \% NaCl).}
\end{popfigure}

\medskip
\noindent\textbf{Document 3 : Extrait du mémoire de Henri Dutrochet (1826) sur l'osmose.}
\begin{quote}
\emph{« J'ai adapté à un tube de verre une vessie de porc remplie d'eau pure, que j'ai plongée dans un vase contenant de l'alcool dilué. Bientôt le liquide s'est élevé dans le tube à une hauteur considérable. La vessie de porc, bien que perméable à l'eau, s'est opposée au passage de l'alcool. J'ai reconnu que ce phénomène, que j'appelle \textbf{endosmose}, est dû au passage de l'eau à travers une membrane semi-perméable, du milieu moins concentré vers le milieu plus concentré. L'exosmose est le phénomène inverse. La pression nécessaire pour arrêter ce flux d'eau définit la \textbf{pression osmotique}. »}
\end{quote}

\courssub{Consignes}
En vous appuyant sur les Documents 1, 2, 3 et sur vos connaissances de la leçon, répondez aux questions suivantes en justifiant rigoureusement chaque affirmation.

\begin{enumerate}
    \item \textbf{Analyse microscopique :} À l'aide d'un crayon, réalisez trois schémas biologiques annotés (taille \textasciitilde 3 cm de diamètre) représentant l'aspect d'une hématie typique dans le Milieu 1, le Milieu 2 et le Milieu 3. Respectez les codes graphiques (membrane plasmique, cytoplasme, noyau absent, forme biconcave/échinocyte/lyse).
    \item \textbf{Interprétation quantitative :} 
    \begin{enumerate}
        \item Calculez la variation relative du volume (\% par rapport au VGM témoin) pour les Milieux 1 et 3.
        \item D'après le Document 2, déterminez la concentration en NaCl (\% massique) pour laquelle le volume cellulaire est maximal. Quel terme qualifie ce milieu par rapport au cytoplasme de l'hématie ?
    \end{enumerate}
    \item \textbf{Identification du mécanisme :} 
    \begin{enumerate}
        \item Nommez le phénomène physique mis en évidence par les Documents 1 et 2.
        \item En vous référant au Document 3, définissez ce phénomène et précisez le sens du mouvement net d'eau pour le Milieu 1 et pour le Milieu 3.
        \item Pourquoi la membrane de l'hématie se comporte-t-elle comme une « membrane semi-perméable » au sens de Dutrochet ?
    \end{enumerate}
    \item \textbf{Choix thérapeutique :} 
    \begin{enumerate}
        \item Quelle poche (A, B ou C) M. Kamga doit-il perfuser à Amadou ? Justifiez votre choix en reliant la tonicité du sérum, le volume des hématies et le risque clinique (hémolyse / création / isovolémie).
        \item Expliquez pourquoi la perfusion d'eau distillée (Poche A) chez un patient hypovolémique est une erreur médicale grave (choc hémolytique).
    \end{enumerate}
    \item \textbf{Physiologie de la réhydratation orale (Relais) :} Une fois la perfusion terminée, le médecin prescrit une solution de réhydratation orale (SRO) type OMS contenant : NaCl (3,5 g/L), Glucose (20 g/L), KCl (1,5 g/L), Citrate de trisodium (2,9 g/L).
    \begin{enumerate}
        \item Rappellez le mécanisme de transport du glucose au niveau de l'entérocyte (bordure en brosse). Précisez le type de transport, le sens, l'énergie requise et la protéine impliquée.
        \item Expliquez comment l'absorption active du glucose « tire » l'absorption du sodium et de l'eau (couplage). Pourquoi le glucose est-il indispensable dans la SRO alors que le sodium seul ne suffit pas en cas de diarrhée sécrétoire (choléra, rotavirus) ?
    \end{enumerate}
\end{enumerate}

\courssub{Ressources à mobiliser}
\begin{itemize}
    \item \textbf{Définitions :} Milieu hypotonique, isotonicité, hypertonique ; Osmose (endosmose / exosmose) ; Pression osmotique ; Transport actif secondaire (cotransport / symport) ; Protéine SGLT1 (Sodium-Glucose Linked Transporter 1).
    \item \textbf{Savoir-faire :} Schématisation biologique (membrane, organites, flèches d'échanges) ; Lecture de courbe (extremum, variations) ; Démarche d'investigation (Observation $\rightarrow$ Hypothèse $\rightarrow$ Interprétation $\rightarrow$ Conclusion).
    \item \textbf{Formule utile :} Variation relative (\%) $= \dfrac{V_{\text{final}} - V_{\text{initial}}}{V_{\text{initial}}} \times 100$.
    \item \textbf{Connaissances clés :} La membrane plasmique est perméable à l'eau (aquaporines) et quasi imperméable aux ions Na$^+$, Cl$^-$, aux sucres (sans transporteur). Le gradient de Na$^+$ (fort $\rightarrow$ faible) est maintenu par la Na$^+$/K$^+$-ATPase (transport actif primaire, consommation d'ATP).
\end{itemize}

\courssub{Démarche et corrigé commenté}
\begin{exoresolu}{Activité d'intégration : Le mystère du sérum physiologique}
\textbf{Énoncé :} Voir consignes ci-dessus.

\tcblower
\textbf{Solution détaillée et commentée}

\medskip
\noindent\textbf{1. Analyse microscopique — Schémas biologiques annotés}
\begin{popfigure}
\begin{tikzpicture}[scale=0.9, every node/.style={font=\small}]
% Styles
\tikzset{
    membrane/.style={very thick, popdark},
    cyto/.style={fill=poppink!30, draw=poppink!80!popdark, thick},
    hb/.style={fill=poppink!80!popdark, draw=poppink!80!popdark},
    arrow/.style={->, >=Stealth, thick, popblue},
    label/.style={popdark, align=center}
}
% --- Milieu 1 : Eau distillée (Hémolyse) ---
\begin{scope}[xshift=0cm]
    % Milieu extracellulaire (eau)
    \fill[popblueL!30] (-2.5,-2.5) rectangle (2.5,2.5);
    % Hématie lysée : membrane éclatée (restes)
    \draw[membrane, dashed, poppink!80!popdark] (0,0) circle (1.8cm);
    \draw[membrane, dashed, poppink!80!popdark] (0,0) circle (2.0cm);
    % Hémoglobine diffusée
    \foreach \i in {1,...,30} {
        \pgfmathsetmacro{\ang}{rand*360}
        \pgfmathsetmacro{\rad}{1.5+rand*0.8}
        \fill[hb, opacity=0.4] (\ang:\rad) circle (0.12cm);
    }
    % Légende
    \node[label] at (0, -3.2) {\textbf{Milieu 1 : Eau distillée (Hypotonique)}};
    \node[label] at (0, -3.8) {Hémolyse : Éclatement + Diffusion Hb};
    % Flèche eau entrant
    \draw[arrow] (-3, 0) -- (-2, 0) node[midway, above] {$\text{H}_2\text{O}$ entrant};
\end{scope}

% --- Milieu 2 : NaCl 0,9% (Isotonique) ---
\begin{scope}[xshift=7cm]
    \fill[popblueL!30] (4.5,-2.5) rectangle (9.5,2.5);
    % Hématie normale disque biconcave (vue de face simplifiée)
    \draw[membrane] (7,0) circle (1.5cm);
    \fill[cyto] (7,0) circle (1.45cm);
    % Aspect biconcave suggéré par ombre interne
    \draw[popdark!30] (7,0) ellipse (1.2cm and 0.3cm);
    \draw[popdark!30] (7,0) ellipse (1.2cm and -0.3cm);
    % Légende
    \node[label] at (7, -3.2) {\textbf{Milieu 2 : NaCl 0,9\% (Isotonique)}};
    \node[label] at (7, -3.8) {Intégrité : Volume stable (VGM = 90 fL)};
    % Flèches équilibre
    \draw[arrow, popmuted] (4, 1.5) -- (5.5, 1.5) node[midway, above] {$\text{H}_2\text{O}$};
    \draw[arrow, popmuted] (8.5, -1.5) -- (10, -1.5) node[midway, below] {$\text{H}_2\text{O}$};
\end{scope}

% --- Milieu 3 : NaCl 5% (Hypertonique) ---
\begin{scope}[xshift=14cm]
    \fill[poporangeL!30] (11.5,-2.5) rectangle (16.5,2.5);
    % Hématie crénée (échinocyte)
    \draw[membrane, poporange!80!popdark] plot[smooth cycle, tension=0.7] coordinates {
        (14, 1.6) (13.5, 1.2) (13.8, 0.8) (13.2, 0.4) (13.6, 0) (13.2, -0.4) (13.8, -0.8) (13.5, -1.2) (14, -1.6)
        (14.5, -1.2) (14.2, -0.8) (14.6, -0.4) (14.2, 0) (14.6, 0.4) (14.2, 0.8) (14.5, 1.2)
    };
    \fill[cyto, opacity=0.8] plot[smooth cycle, tension=0.7] coordinates {
        (14, 1.5) (13.55, 1.1) (13.85, 0.7) (13.25, 0.3) (13.65, -0.05) (13.25, -0.35) (13.85, -0.75) (13.55, -1.1) (14, -1.5)
        (14.45, -1.1) (14.15, -0.7) (14.55, -0.3) (14.15, 0.05) (14.55, 0.35) (14.15, 0.7) (14.45, 1.1)
    };
    % cytoplasme concentré (plus foncé)
    \fill[poppink!80!popdark, opacity=0.5] plot[smooth cycle, tension=0.7] coordinates {
        (14, 1.4) (13.6, 1.0) (13.9, 0.6) (13.35, 0.2) (13.75, -0.15) (13.35, -0.25) (13.9, -0.65) (13.6, -1.0) (14, -1.4)
        (14.4, -1.0) (14.1, -0.6) (14.5, -0.2) (14.1, 0.15) (14.5, 0.25) (14.1, 0.6) (14.4, 1.0)
    };
    \node[label] at (14, -3.2) {\textbf{Milieu 3 : NaCl 5\% (Hypertonique)}};
    \node[label] at (14, -3.8) {Création : Volume $\downarrow$ (55 fL), Hb concentrée};
    % Flèche eau sortant
    \draw[arrow, poporange] (14, 2.2) -- (14, 3) node[midway, right] {$\text{H}_2\text{O}$ sortant};
\end{scope}
\end{tikzpicture}
\legende{Schémas biologiques annotés d'hématies humaines dans les trois milieux testés. Échelle non respectée entre milieux pour lisibilité. Membrane plasmique en trait plein (intacte) ou pointillé (rompue). Cytoplasme/hémoglobine en rose/rouge. Flèches bleues/oranges : sens net du flux d'eau (osmose).}
\end{popfigure}

\medskip
\noindent\textbf{Commentaire du correcteur :} Les schémas doivent montrer : (1) Milieu 1 : membrane rompue (pointillés), hémoglobine diffusée dans le milieu (pas de noyau), flèche d'eau \textbf{entrante} massive. (2) Milieu 2 : disque biconcave régulier, double concavité suggérée, flèches d'eau \textbf{équilibrées} (aller/retour). (3) Milieu 3 : membrane repliée en éperons (échinocyte), cytoplasme rétracté et foncé (concentration Hb), flèche d'eau \textbf{sortante}. \textbf{Barème :} 1 pt par schéma correct (forme + annotations + flèche eau).

\medskip
\noindent\textbf{2. Interprétation quantitative}
\begin{enumerate}
    \item \textbf{Calcul des variations relatives :}
    \begin{itemize}
        \item \textbf{Milieu 1 (Eau distillée) :} Le volume mesuré est $> 150$ fL (lyse lors de la mesure). Prenons la valeur seuil $V_{\text{lyse}} \approx 150$ fL (volume critique avant rupture).
        \[ \Delta V_{\%} = \frac{150 - 90}{90} \times 100 = \frac{60}{90} \times 100 \approx \mathbf{+66{,}7\%} \]
        \textit{Note :} En réalité, le volume augmente jusqu'à la rupture (lyse), la mesure automate ne capture que les débris ou surestime par diffraction.
        \item \textbf{Milieu 3 (NaCl 5 \%) :} $V_{\text{final}} = 55$ fL.
        \[ \Delta V_{\%} = \frac{55 - 90}{90} \times 100 = \frac{-35}{90} \times 100 \approx \mathbf{-38{,}9\%} \]
    \end{itemize}
    \item \textbf{Lecture de la courbe (Document 2) :} Le volume relatif est maximal (\textbf{100 \%}) pour une concentration en NaCl de \textbf{0,9 \%}. Ce milieu est qualifié \textbf{d'isotonique} (ou isosmotique) par rapport au cytoplasme de l'hématie. À cette concentration, la pression osmotique du milieu extracellulaire est égale à la pression osmotique intracellulaire : il n'y a pas de flux net d'eau.
\end{enumerate}

\medskip
\noindent\textbf{3. Identification du mécanisme}
\begin{enumerate}
    \item Le phénomène mis en évidence est l'\textbf{osmose} (plus précisément l'endosmose en milieu hypotonique, l'exosmose en milieu hypertonique).
    \item \textbf{Définition (d'après Dutrochet, Doc 3) :} L'osmose est le passage net d'eau à travers une \textbf{membrane semi-perméable} (perméable à l'eau, imperméable aux solutés), du milieu \textbf{moins concentré} (hypotonique, potentiel d'eau élevé) vers le milieu \textbf{plus concentré} (hypertonique, potentiel d'eau bas).
    \begin{itemize}
        \item \textbf{Milieu 1 (Eau distillée) :} Le milieu extracellulaire est hypotonique (0 \% NaCl) par rapport au cytoplasme (\textasciitilde 0,9 \% équivalent NaCl). L'eau entre massivement dans la cellule (\textbf{endosmose}) $\rightarrow$ gonflement $\rightarrow$ lyse.
        \item \textbf{Milieu 3 (NaCl 5 \%) :} Le milieu extracellulaire est hypertonique par rapport au cytoplasme. L'eau sort de la cellule (\textbf{exosmose}) $\rightarrow$ rétraction cytoplasmique $\rightarrow$ création.
    \end{itemize}
    \item La membrane plasmique de l'hématie possède des \textbf{aquaporines} (canaux à eau) qui la rendent très perméable à l'eau. En revanche, elle est \textbf{imperméable aux ions Na$^+$ et Cl$^-$} en l'absence de transporteurs spécifiques (canaux ioniques fermés au repos, pas de transporteur Na$^+$/glucose sur hématie). Elle remplit donc parfaitement le critère de « membrane semi-perméable » de Dutrochet : sélective pour le solvant (eau), barrière pour les solutés (ions).
\end{enumerate}

\medskip
\noindent\textbf{4. Choix thérapeutique}
\begin{enumerate}
    \item \textbf{Choix : Poche B (NaCl 0,9 \%).}
    \textbf{Justification rigoureuse :}
    \begin{itemize}
        \item La Poche B est \textbf{isotonique} par rapport au sang (plasme) et aux hématies (Doc 1 : volume stable 90 fL, forme normale ; Doc 2 : maximum de la courbe à 0,9 \%).
        \item Perfuser cette solution maintient l'\textbf{isovolémie} globulaire : pas de flux net d'eau entrant (risque hémolyse) ni sortant (risque création). Les hématies conservent leur fonction de transport d'O$_2$ (déformabilité, surface d'échange).
        \item \textbf{Poche A (Eau distillée) :} Hypotonique. Provoquerait une hémolyse massive intravasculaire $\rightarrow$ libération d'hémoglobine libre (néphrotoxique, risque d'insuffisance rénale aiguë), hyperkaliémie (libération K$^+$ intracellulaire $\rightarrow$ arythmies cardiaques fatales), coagulation intravasculaire disseminée (CIVD). \textbf{Contre-indication absolue.}
        \item \textbf{Poche C (NaCl 5 \%) :} Hypertonique. Provoquerait une création généralisée des hématies $\rightarrow$ augmentation de la viscosité sanguine, ralentissement de la microcirculation, risque thrombose, douleur à la perfusion, déshydratation cellulaire intracellulaire aggravée (eau sort des cellules vers le sang). \textbf{Inadaptée pour une réhydratation volémique.}
    \end{itemize}
    \item L'erreur médicale de la Poche A s'explique par le \textbf{gradient de pression osmotique} extrême. Le plasma du patient, bien qu'hémoconcentré par la déshydratation (hématocrite $>$ 55 \%), reste hypertonique par rapport à l'eau distillée. L'infusion d'eau pure dilue brutalement le plasma $\rightarrow$ chute de l'osmolarité plasmatique $\rightarrow$ entrée massive d'eau dans les hématies (endosmose) $\rightarrow$ \textbf{choc hémolytique}. La libération d'hémoglobine (Hb) libre happe le NO (vasodilatateur) $\rightarrow$ vasoconstriction rénale + toxicité tubulaire directe $\rightarrow$ \textbf{insuffisance rénale aiguë}. Chez un enfant drépanocytaire (fréquent au Cameroun), le risque est majoré.
\end{enumerate}

\medskip
\noindent\textbf{5. Physiologie de la réhydratation orale (Relais)}
\begin{enumerate}
    \item \textbf{Transport du glucose au niveau de l'entéocyte (bordure en brosse) :}
    \begin{itemize}
        \item \textbf{Type :} Transport actif \textbf{secondaire} (cotransport / \textbf{symport}).
        \item \textbf{Sens :} Du lumen intestinal (milieu extérieur) $\rightarrow$ cytoplasme de l'entéocyte (milieu intérieur). \textbf{Dans le sens du gradient de concentration du Na$^+$} (fort $\rightarrow$ faible) et \textbf{contre le gradient de concentration du glucose} (faible $\rightarrow$ fort).
        \item \textbf{Énergie :} Fournie indirectement par le gradient électrochimique du Na$^+$, lui-même maintenu par la \textbf{Na$^+$/K$^+$-ATPase} (transport actif primaire, consommation d'ATP) localisée sur la membrane basolatérale.
        \item \textbf{Protéine :} \textbf{SGLT1} (Sodium-Glucose Linked Transporter 1). Stœchiométrie : 2 Na$^+$ / 1 Glucose.
    \end{itemize}
    \item \textbf{Couplage Glucose / Sodium / Eau (Mécanisme de la SRO) :}
    \begin{enumerate}
        \item L'absorption active du glucose via SGLT1 \textbf{« tire »} le Na$^+$ depuis le lumen vers la cellule (symport obligatoire).
        \item Cette entrée de Na$^+$ (et du glucose) abaisse le potentiel d'eau (osmolarité $\uparrow$) à l'intérieur de l'entéocyte et de l'espace intercellulaire (jonctions serrées).
        \item L'eau suit passivement le gradient osmotique (\textbf{osmose}) : Lumen $\rightarrow$ Entéocyte $\rightarrow$ Espace intercellulaire $\rightarrow$ Capillaire sanguin (voie paracellulaire et transcellulaire via aquaporines).
        \item Le Na$^+$ est ensuite extrudé vers le sang par la Na$^+$/K$^+$-ATPase basolatérale, maintenant le gradient.
    \end{enumerate}
    \textbf{Pourquoi le glucose est-il indispensable ?}
    En cas de diarrhée sécrétoire (choléra : toxine CTX active adénylate cyclase $\rightarrow$ cAMP $\rightarrow$ ouverture canal CFTR $\rightarrow$ sécrétion massive Cl$^-$ + Na$^+$ + H$_2$O ; Rotavirus : atteinte entérocytes), le transport \textbf{passif} du Na$^+$ (canaux ENaC) est soit inhibé, soit submergé par la sécrétion active. Le transporteur \textbf{SGLT1 n'est pas affecté} par la toxine cholérique (voie de signalisation différente). L'apport de glucose dans la SRO \textbf{réactive l'absorption couplée Na$^+$-Glucose}, créant un gradient osmotique local qui « pompe » l'eau vers le sang, \textbf{contrecarrant la sécrétion pathologique}. Sans glucose, l'absorption de Na$^+$ (et donc d'eau) reste défaillante. C'est la base physiologique du « \textbf{couplage glucose-sodium} » qui a valu le prix Nobel à Crane (1960) et révolutionné la prise en charge du choléra (réduction mortalité de 30\% à $<$ 1\%).
\end{enumerate}
\end{exoresolu}

\courssub{Bilan de l'activité}
Cette activité d'intégration a permis de mettre en évidence, par une démarche d'investigation complète, les points fondamentaux suivants :
\begin{enumerate}
    \item \textbf{L'osmose est le moteur principal des échanges d'eau} entre la cellule et son milieu. Elle est dictée par la différence de pression osmotique (tonicité) de part et d'autre de la membrane plasmique, semi-perméable (perméable à l'eau via aquaporines, imperméable aux solutés sans transporteurs).
    \item \textbf{L'isotonicité (0,9 \% NaCl pour le sang humain) est la condition sine qua non de l'intégrité cellulaire} (isovolémie). Tout écart vers l'hypotonicité (hémolyse) ou l'hypertonicité (création) est pathologique et potentiellement mortel en clinique.
    \item \textbf{L'expérience de Dutrochet (1826)} a posé les bases conceptuelles (membrane semi-perméable, flux d'eau, pression osmotique) qui expliquent encore aujourd'hui la physiologie des perfusions et la pathologie de la déshydratation.
    \item \textbf{La réhydratation orale (SRO/OMS) repose sur un transport actif secondaire (SGLT1)} qui couple l'absorption du glucose à celle du sodium, créant le gradient osmotique nécessaire à l'absorption passive de l'eau. C'est une application directe et salvatrice de la connaissance des transports membranaires (actif primaire Na$^+$/K$^+$-ATPase $\rightarrow$ gradient Na$^+$ $\rightarrow$ symport SGLT1 $\rightarrow$ osmose).
    \item \textbf{Compétences mobilisées :} Observation microscopique (différenciation hémolyse/création/normale), schématisation biologique rigoureuse, analyse quantitative (calculs \%, lecture courbe), modélisation (membrane semi-perméable), argumentation clinique (choix thérapeutique justifié), transfert de connaissances (perfusion IV $\rightarrow$ réhydratation orale).
\end{enumerate}
\begin{aretenir}[Message clé pour le Bac]
\begin{itemize}
    \item \textbf{Sérum physiologique = NaCl 0,9 \% (isotonique).} C'est le \textbf{seul} liquide de remplissage vasculaire de référence en urgence (hémorragie, déshydratation sévère, choc).
    \item \textbf{Jamais d'eau distillée en IV :} Risque d'hémolyse massive $\rightarrow$ choc hémolytique + IRA.
    \item \textbf{SRO = Eau + Sel (NaCl) + Sucre (Glucose).} Le glucose n'est pas un « apport énergétique » primaire ici, mais le \textbf{co-substrat obligatoire} du transporteur SGLT1 pour « pomper » le Na$^+$ et l'eau contre la sécrétion diarrhéique.
    \item \textbf{Schéma type hématie :} Disque biconcave (isotonique) / Sphère éclatée (hypotonique) / Échinocyte (hypertonique). Flèche eau : entrant / sortant / bidirectionnelle.
\end{itemize}
\end{aretenir}

\begin{attention}[Pièges fréquents à l'examen]
\begin{itemize}
    \item Confondre \textbf{osmose} (eau, membrane semi-perméable) et \textbf{diffusion} (solutés, gradient de concentration, pas besoin de membrane semi-perméable).
    \item Dire « l'eau va vers le plus concentré » sans préciser : « en \textbf{solutés non pérméants} » (tonicité effective). L'urée, par ex., traverse la membrane $\rightarrow$ ne crée pas d'osmose durable.
    \item Oublier que le transport actif du glucose (SGLT1) est \textbf{secondaire} : l'énergie vient du gradient Na$^+$ (Na$^+$/K$^+$-ATPase), pas de l'ATP directement sur SGLT1.
    \item Schémas bâclés : pas de membrane, pas de flèche eau, forme biconcave non suggérée, noyau dessiné (hématie anucléée !).
    \item Unités oubliées sur les axes de courbe ou dans les calculs (fL, \%, g/L).
\end{itemize}
\end{attention}

\newpage

\coursec{Méthodes et exercices résolus}

\coursec{Organisation de la cellule et échanges cellulaires}

\courssub{Savoir-faire 1 : Observer et schématiser la cellule animale en microscopie optique et électronique}

\begin{methode}[Schématiser une cellule animale à partir d'une observation]
\begin{enumerate}
\item \textbf{Identifier l'échelle d'observation.} En microscopie optique (MO), le pouvoir de résolution est d'environ \SI{0,2}{\micro\metre} : on distingue trois structures principales : la \cle{membrane cytoplasmique}, le \cle{cytoplasme} et le \cle{noyau}. En microscopie électronique (ME), le pouvoir de résolution atteint \SI{0,2}{nm} (soit 1000 fois meilleur) : on visualise les \cle{organites} (mitochondries, réticulum endoplasmique, appareil de Golgi, lysosomes, ribosomes, centrioles).
\item \textbf{Repérer les structures visibles} sur le document (photo, micrographie) à l'aide des contrastes et des flèches indicatrices : le noyau est la structure la plus volumineuse et la plus colorée (hématéine, bleu de méthylène) ; les mitochondries apparaissent comme des bâtonnets.
\item \textbf{Dessiner au crayon à papier}, en traits fins et continus, sans ombres ni hachures : le schéma biologique est un schéma d'interprétation, non un dessin artistique.
\item \textbf{Annoter avec des lignes d'attache horizontales}, tracées à la règle, aboutissant exactement sur la structure désignée, sans jamais se croiser. Le titre du schéma est souligné et précise : l'objet observé, le type de microscopie, le grossissement.
\item \textbf{Calculer le grossissement ou la taille réelle} avec la relation : $G = \dfrac{\text{taille mesurée sur le document}}{\text{taille réelle}}$, en convertissant toutes les longueurs dans la même unité (souvent le micromètre \si{\micro\metre}).
\end{enumerate}
\end{methode}

\begin{popfigure}
\begin{tikzpicture}[scale=0.8, every node/.style={font=\small}]
% Cellule MO
\draw[popdark, thick] (0,0) ellipse (3cm and 2cm); % Membrane
\draw[popblue, thick] (0,0.2) ellipse (1.2cm and 1cm); % Noyau
\draw[popblue, fill=popblue!20] (0,0.2) ellipse (0.4cm and 0.3cm); % Nucléole
\foreach \x/\y in {-2.2/-1.2, -1.5/1.5, 1.8/-0.5, 2.2/1.2, -0.5/-1.5, 0.8/1.6} {
    \fill[popmuted] (\x,\y) circle (0.08cm); % Granularité cytoplasme
}
% Annotations
\draw[->, popdark] (3.5,1.5) -- (2.5,1.2) node[right, align=left] {Membrane\\cytoplasmique};
\draw[->, popdark] (3.5,0.5) -- (1.8,0.5) node[right] {Cytoplasme};
\draw[->, popdark] (3.5,-0.5) -- (1.2,0.2) node[right, align=left] {Noyau\\(coloré)};
\draw[->, popdark] (3.5,-1.5) -- (0.4,0.0) node[right] {Nucléole};
\end{tikzpicture}
\legende{Schéma d'interprétation d'une cellule animale (épithélium buccal) observée en microscopie optique (coloration bleu de méthylène), grossissement $\times 1000$.}
\end{popfigure}

\begin{exoresolu}[Schématisation d'une cellule animale observée en MO]
\textbf{Énoncé.} Un élève du lycée de Nkongsamba observe au microscope optique (grossissement $G = \times 1000$) des cellules de l'épithélium de sa joue, colorées au bleu de méthylène. Sur sa préparation, il mesure le diamètre d'une cellule : 4 cm sur la photo prise avec son téléphone.
\begin{enumerate}
\item Cite les structures que l'élève peut observer à ce grossissement.
\item Calcule le diamètre réel de la cellule.
\item Décris les règles à respecter pour réaliser le schéma d'observation.
\end{enumerate}
\tcblower
\textbf{Solution détaillée.}
\begin{enumerate}
\item \textbf{Structures observables en MO.} Le pouvoir de résolution de la microscopie optique est limité à environ \SI{0,2}{\micro\metre} : seules trois structures sont nettement visibles dans une cellule animale :
\begin{itemize}
\item la \cle{membrane cytoplasmique} (limite externe de la cellule) ;
\item le \cle{cytoplasme} (milieu granuleux occupant l'intérieur) ;
\item le \cle{noyau}, fortement coloré par le bleu de méthylène car il contient l'ADN (chromatine).
\end{itemize}
Les organites (mitochondries, réticulum, Golgi, lysosomes) ont des dimensions inférieures au pouvoir de résolution : ils ne sont pas visibles en MO.

\item \textbf{Calcul du diamètre réel.} On applique la relation du grossissement :
\[ G = \frac{\text{taille sur le document}}{\text{taille réelle}} \quad\Longrightarrow\quad \text{taille réelle} = \frac{\text{taille sur le document}}{G}. \]
Conversion dans la même unité : 4 cm $= \SI{40000}{\micro\metre}$. D'où :
\[ d_{\text{réel}} = \frac{\SI{40000}{\micro\metre}}{1000} = \SI{40}{\micro\metre}. \]
Le diamètre réel de la cellule est de \SI{40}{\micro\metre}, valeur cohérente avec la taille moyenne d'une cellule animale (\SIrange{10}{50}{\micro\metre}).

\item \textbf{Règles du schéma d'observation.} Dessin au crayon à papier, traits continus et fins, proportions respectées ; titre souligné précisant la nature de la cellule, la coloration et le grossissement ; annotations par lignes d'attache horizontales tracées à la règle, ne se croisant pas, désignant : membrane cytoplasmique, cytoplasme, noyau, nucléole.
\end{enumerate}
\textbf{Conclusion.} En microscopie optique, la cellule animale montre trois structures principales (membrane, cytoplasme, noyau) ; la cellule observée a un diamètre réel de \SI{40}{\micro\metre}.
\end{exoresolu}

\courssub{Savoir-faire 2 : Schématiser et annoter les organites cellulaires}

\begin{methode}[Identifier et annoter les organites sur une micrographie électronique]
\begin{enumerate}
\item \textbf{Repérer les indices morphologiques} caractéristiques de chaque organite :
\begin{itemize}
\item \cle{Mitochondrie} : bâtonnet à double membrane, la membrane interne formant des \emph{crêtes mitochondriales} (site de la respiration cellulaire) ;
\item \cle{Réticulum endoplasmique granuleux (REG)} : réseau de saccules (cisternae) aplatis parsemés de grains sombres (ribosomes attachés) ; le \cle{réticulum endoplasmique lisse (REL)} en est dépourvu (aspect tubulaire lisse) ;
\item \cle{Appareil de Golgi} : empilement de 3 à 8 saccules aplatis (cisternae) avec des vésicules de bourgeonnement aux extrémités (face \emph{cis} et face \emph{trans}) ;
\item \cle{Lysosomes} : petites vésicules sphériques à contenu dense et homogène (enzymes hydrolytiques) ;
\item \cle{Ribosomes} : petits grains (20-30 nm) libres dans le cytoplasme ou fixés sur le REG (site de synthèse des protéines) ;
\item \cle{Centrioles} : deux cylindres courts et perpendiculaires (diplosome), composés de triplets de microtubules, près du noyau (centrosome) ;
\item \cle{Noyau} : grande sphère entourée de l'\cle{enveloppe nucléaire} (double membrane percée de \emph{pores nucléaires}), contenant la \cle{chromatine} (filaments d'ADN) et le \cle{nucléole} (site de synthèse de l'ARN ribosomal).
\end{itemize}
\item \textbf{Associer chaque organite à son rôle fonctionnel} : mitochondrie $\to$ respiration cellulaire (production d'ATP) ; REG $\to$ synthèse des protéines destinées à la sécrétion ou aux membranes ; REL $\to$ synthèse des lipides, détoxication ; Golgi $\to$ maturation (glycosylation), tri et conditionnement des protéines dans des vésicules de sécrétion ; lysosomes $\to$ digestion intracellulaire (autophagie, hétérophagie) ; ribosomes $\to$ traduction de l'ARNm en protéines.
\item \textbf{Annoter} selon les règles du schéma biologique (lignes d'attache horizontales, titre souligné avec échelle/grossissement).
\item \textbf{Vérifier la cohérence structure-fonction} : une cellule très sécrétrice (ex. : cellule pancréatique $\beta$, plasmocyte) présente un REG et un appareil de Golgi très développés ; une cellule très consommatrice d'énergie (ex. : cardiomyocyte, fibre musculaire striée, tubule proximal du rein) est riche en mitochondries.
\end{enumerate}
\end{methode}

\begin{popfigure}
\begin{tikzpicture}[scale=0.7, every node/.style={font=\small}]
% Membrane cytoplasmique
\draw[popdark, thick] (-4,-3) rectangle (4,3);
% Noyau
\draw[popblue, thick] (-1,0.5) ellipse (1.5cm and 1.2cm); % Enveloppe nucléaire externe
\draw[popblue, thick] (-1,0.5) ellipse (1.35cm and 1.05cm); % Interne
% Pores nucléaires
\foreach \angle in {30, 90, 150, 210, 270, 330} {
    \draw[popblue] (-1,0.5) ++(\angle:1.45) -- ++(\angle:0.2);
}
% Nucléole
\draw[popblue, fill=popblue!30] (-1.3,0.8) circle (0.35cm);
% Chromatine
\foreach \i in {1,...,15} {
    \draw[popblue!70] (-2.2+0.25*rand, -0.2+0.2*rand) -- ++(0.15+0.1*rand, 0.1+0.05*rand);
}
% Mitochondries
\foreach \pos in {(-3.2,-2.2), (2.5,-2.5), (-3.5,1.5), (3.0,1.8)} {
    \draw[poporange, thick, rotate=45] \pos ellipse (0.4cm and 0.15cm);
    \draw[poporange, thick, rotate=45] \pos ellipse (0.3cm and 0.08cm); % Crêtes
}
% REG (accollé au noyau)
\draw[poppurple, thick] (-2.8,1.8) .. controls (-3.5,1.5) and (-3.5,-1.5) .. (-2.8,-1.8);
\draw[poppurple, thick] (-2.5,1.8) .. controls (-3.2,1.5) and (-3.2,-1.5) .. (-2.5,-1.8);
% Ribosomes sur REG
\foreach \y in {1.5, 1.0, 0.5, 0, -0.5, -1.0, -1.5} {
    \fill[poppurple] (-3.0, \y) circle (0.04cm);
    \fill[poppurple] (-3.1, \y) circle (0.04cm);
}
% REL
\draw[poppurple, thick, dashed] (2.0, -1.0) .. controls (2.5, -0.5) .. (2.0, 0.0);
\draw[poppurple, thick, dashed] (2.2, -1.0) .. controls (2.7, -0.5) .. (2.2, 0.0);
% Appareil de Golgi
\foreach \y/\xoff in {1.2/0, 0.8/0.1, 0.4/0, 0.0/-0.1, -0.4/0} {
    \draw[popgreen, thick] (1.5+\xoff, \y) -- (2.5+\xoff, \y);
    \draw[popgreen, thick] (1.5+\xoff, \y+0.08) -- (2.5+\xoff, \y+0.08);
}
% Vésicules Golgi
\foreach \pos in {(2.7,1.4), (2.8,0.6), (2.6,-0.2)} {
    \draw[popgreen, fill=popgreen!20] \pos circle (0.12cm);
}
% Lysosomes
\foreach \pos in {(-2.5,-2.5), (0.5,-2.8), (3.5,-1.0)} {
    \draw[poppink, fill=poppink!30] \pos circle (0.15cm);
}
% Ribosomes libres
\foreach \pos in {(-3.0, -0.5), (-0.5, -2.0), (0.0, 2.0), (-3.8, 2.5)} {
    \fill[popmuted] \pos circle (0.03cm);
}
% Centrioles
\draw[popgold, thick] (0.5, 1.0) rectangle (0.8, 1.5);
\draw[popgold, thick] (0.3, 1.2) rectangle (1.0, 1.25); % Perpendiculaire (simplifié)
% Annotations (flèches simplifiées pour l'exemple)
\node[popdark, align=left, anchor=west] at (4.5, 2.5) {\textbf{Légende :}};
\node[popdark, align=left, anchor=west] at (4.5, 2.1) {1. Membrane cytoplasmique};
\node[popblue, align=left, anchor=west] at (4.5, 1.7) {2. Noyau (enveloppe, pores, nucléole, chromatine)};
\node[poporange, align=left, anchor=west] at (4.5, 1.3) {3. Mitochondrie (crêtes)};
\node[poppurple, align=left, anchor=west] at (4.5, 0.9) {4. Réticulum endoplasmique granuleux (REG) + Ribosomes};
\node[poppurple, align=left, anchor=west] at (4.5, 0.5) {5. Réticulum endoplasmique lisse (REL)};
\node[popgreen, align=left, anchor=west] at (4.5, 0.1) {6. Appareil de Golgi + Vésicules};
\node[poppink, align=left, anchor=west] at (4.5, -0.3) {7. Lysosome};
\node[popgold, align=left, anchor=west] at (4.5, -0.7) {8. Centrioles (diplosome)};
\end{tikzpicture}
\legende{Schéma d'interprétation d'une cellule animale observée en microscopie électronique (coupe ultramince). Les organites sont représentés avec leur ultrastructure caractéristique.}
\end{popfigure}

\begin{exoresolu}[Analyse d'une micrographie de cellule pancréatique]
\textbf{Énoncé.} Une micrographie électronique d'une cellule du pancréas (organe qui sécrète l'insuline, une protéine) montre : un noyau volumineux, un réticulum endoplasmique granuleux très développé, un appareil de Golgi important, de nombreuses vésicules et des mitochondries.
\begin{enumerate}
\item Schématise cette cellule en annotant quatre organites.
\item Explique l'abondance du REG et de l'appareil de Golgi dans cette cellule.
\end{enumerate}
\tcblower
\textbf{Solution détaillée.}
\begin{enumerate}
\item \textbf{Schéma.} Le schéma représente une cellule de contour irrégulier limitée par la membrane cytoplasmique, contenant : le noyau (grande sphère à double membrane poreuse), le REG (réseau de saccules granuleux accolé à l'enveloppe nucléaire), l'appareil de Golgi (empilement de saccules lisses avec vésicules de sécrétion bourgeonnant face \emph{trans}), les mitochondries (bâtonnets à crêtes internes). Annotations par lignes d'attache horizontales ; titre : \emph{« Schéma d'interprétation d'une cellule pancréatique $\beta$ observée en ME, $\times 10\,000$ »}.
\item \textbf{Interprétation.} L'insuline est une \cle{protéine} destinée à être \cle{sécrétée} par exocytose (voie de sécrétion régulée). Or :
\begin{itemize}
\item les protéines sécrétées sont synthétisées par les \emph{ribosomes} fixés sur le \emph{réticulum endoplasmique granuleux} : d'où un REG très développé (forte densité de ribosomes) ;
\item les protéines nouvellement formées transitent dans la lumière du REG vers l'\emph{appareil de Golgi} (face \emph{cis}), où elles subissent des modifications post-traductionnelles (glycosylation), sont triées et conditionnées dans des \emph{vésicules de sécrétion} (face \emph{trans}) ;
\item ces vésicules migrent vers la membrane plasmique pour libérer leur contenu par \cle{exocytose} ; ces processus (synthèse, trafic vésiculaire, exocytose) consomment de l'énergie fournie sous forme d'ATP par les \emph{mitochondries}.
\end{itemize}
\end{enumerate}
\textbf{Conclusion.} L'abondance du REG, de l'appareil de Golgi et des mitochondries dans la cellule pancréatique $\beta$ est en relation directe avec sa fonction spécialisée : la synthèse, la maturation et la sécrétion d'une protéine hormonale, l'insuline. C'est une illustration majeure de l'adéquation structure-fonction.
\end{exoresolu}

\courssub{Savoir-faire 3 : Analyser, interpréter et conclure à partir des expériences sur les échanges d'eau}

\begin{definition}[Milieux toniques et osmose]
On compare la concentration en solutés d'une solution (milieu extérieur) à celle du cytoplasme (milieu intérieur) :
\begin{itemize}
\item \cle{Milieu isotonique} : même concentration en solutés. Échanges d'eau équilibrés (volume cellulaire constant).
\item \cle{Milieu hypotonique} : concentration en solutés plus faible. L'eau entre dans la cellule par \cle{osmose} (risque d'éclatement = \cle{hémolyse} pour l'hématie).
\item \cle{Milieu hypertonique} : concentration en solutés plus forte. L'eau sort de la cellule par osmose (déshydratation = \cle{crénelation} pour l'hématie).
\end{itemize}
L'\cle{osmose} est le déplacement net de l'eau (solvant) à travers une membrane \emph{semi-perméable} (perméable à l'eau, imperméable aux solutés), du milieu hypotonique vers le milieu hypertonique.
\end{definition}

\begin{definition}[Dialyse]
La \cle{dialyse} est le déplacement net des \emph{solutés} (petites molécules dissoutes : ions, glucose, urée...) à travers une membrane \emph{dialysante} (perméable à l'eau et aux petites molécules, imperméable aux macromolécules), du milieu le plus concentré vers le milieu le moins concentré (sens du gradient de concentration). Les macromolécules (protéines, amidon) sont retenues.
\end{definition}

\begin{methode}[Interpréter une expérience d'osmose ou de dialyse]
\begin{enumerate}
\item \textbf{Dégager le problème} posé par l'observation (variation de volume, de masse, de niveau liquide, coloration) et formuler une \cle{hypothèse} explicative.
\item \textbf{Identifier les conditions expérimentales} : nature de la membrane (semi-perméable vs dialysante), nature et concentration des solutions de part et d'autre (hypotonique, hypertonique, isotonique par rapport au contenu cellulaire ou à la solution testée).
\item \textbf{Comparer les résultats} des différents lots (témoin / expérimental) : variation de masse, de volume, aspect microscopique des cellules (turgescence, plasmolyse, éclatement, crénelation).
\item \textbf{Interpréter} en précisant le sens du mouvement net d'eau (osmose) ou de solutés (dialyse) et la nature de la membrane.
\item \textbf{Conclure} en répondant au problème posé : la conclusion généralise le mécanisme mis en évidence (loi de l'osmose, sélectivité de la membrane).
\end{enumerate}
\textbf{Réflexes :} Toujours préciser le sens du mouvement net d'eau (du milieu hypotonique vers le milieu hypertonique) ; relier le phénomène au comportement des hématies (éclatement en milieu hypotonique = hémolyse ; flétrissement en milieu hypertonique = crénelation) ; pour la dialyse, tester la présence des solutés des deux côtés (ex. : iode pour l'amidon, bandelette pour le glucose).
\end{methode}

\begin{popfigure}
\begin{tikzpicture}[scale=1.0, every node/.style={font=\small}]
% Tube A : Isotonique
\begin{scope}[xshift=-6cm]
\draw[popdark, thick] (-1,0) rectangle (1,4);
\draw[popblue!30, fill=popblue!30] (-1,0) rectangle (1,4);
% Hématies normales (disques biconcaves)
\foreach \y in {0.5, 1.2, 1.9, 2.6, 3.3} {
    \draw[popdark, thick] (-0.5,\y) ellipse (0.4cm and 0.1cm);
    \draw[popdark, thick] (0.2,\y) ellipse (0.4cm and 0.1cm);
}
\node[popdark] at (0,4.3) {Tube A : \SI{9}{g/L} NaCl};
\node[popdark] at (0,4.6) {\textbf{Isotonique}};
\node[popdark, font=\tiny, align=center] at (0,-0.5){Hématies normales\\(disques biconcaves)};
\end{scope}
% Tube B : Hypotonique
\begin{scope}[xshift=0cm]
\draw[popdark, thick] (-1,0) rectangle (1,4);
\draw[popblue!30, fill=popblue!30] (-1,0) rectangle (1,4);
% Hématies gonflées / éclatées
\foreach \y in {0.5, 1.5, 2.5, 3.5} {
    \draw[poporange, thick, fill=poporange!20] (0,\y) circle (0.35cm); % Gonflées
}
\draw[poporange, thick, fill=poporange!20] (0,3.5) circle (0.35cm);
\draw[poporange, dashed] (-0.5,3.5) -- (-0.8,3.8); % Éclatement
\draw[poporange, dashed] (0.5,3.5) -- (0.8,3.8);
\node[popdark] at (0,4.3) {Tube B : \SI{3}{g/L} NaCl};
\node[popdark] at (0,4.6) {\textbf{Hypotonique}};
\node[popdark, font=\tiny, align=center] at (0,-0.5){Hématies gonflées,\\certaines éclatées\\(\textbf{Hémolyse})};
\draw[->, popblue, thick] (1.5,2) -- (2.5,2) node[right] {Entrée d'eau (Osmose)};
\end{scope}
% Tube C : Hypertonique
\begin{scope}[xshift=6cm]
\draw[popdark, thick] (-1,0) rectangle (1,4);
\draw[popblue!30, fill=popblue!30] (-1,0) rectangle (1,4);
% Hématies crénelées
\foreach \y in {0.5, 1.3, 2.1, 2.9, 3.7} {
    \draw[poppurple, thick] (0,\y) plot[smooth cycle, tension=0.7] coordinates {(-0.35,\y) (-0.15,\y+0.15) (0.1,\y) (0.15,\y-0.15) (0.35,\y) (0.15,\y+0.15) (-0.1,\y) (-0.15,\y-0.15)};
}
\node[popdark] at (0,4.3) {Tube C : \SI{15}{g/L} NaCl};
\node[popdark] at (0,4.6) {\textbf{Hypertonique}};
\node[popdark, font=\tiny, align=center] at (0,-0.5){Hématies ratatinées\\(\textbf{Crénelation})};
\draw[->, popblue, thick] (-2.5,2) -- (-3.5,2) node[left] {Sortie d'eau (Osmose)};
\end{scope}
\end{tikzpicture}
\legende{Comportement des hématies humaines dans des milieux de tonicité différente. A : Milieu isotonique (sérum physiologique), équilibre. B : Milieu hypotonique, entrée d'eau $\to$ hémolyse. C : Milieu hypertonique, sortie d'eau $\to$ crénelation.}
\end{popfigure}

\begin{exoresolu}[Hématies placées dans des solutions de concentrations différentes]
\textbf{Énoncé.} On place des hématies humaines dans trois tubes contenant des solutions de NaCl de concentrations différentes. Après 30 minutes, on observe au microscope : tube A (NaCl à \SI{9}{\gram\per\litre}) : hématies normales ; tube B (NaCl à \SI{3}{\gram\per\litre}) : hématies gonflées, certaines éclatées ; tube C (NaCl à \SI{15}{\gram\per\litre}) : hématies ratatinées.
\begin{enumerate}
\item Formule le problème posé par cette expérience et une hypothèse.
\item Interprète les résultats de chaque tube.
\item Conclus.
\end{enumerate}
\tcblower
\textbf{Solution détaillée.}
\begin{enumerate}
\item \textbf{Problème.} Comment expliquer les modifications d'aspect et de volume des hématies placées dans des solutions de salinité différente ? \textbf{Hypothèse.} Les hématies échangent de l'eau avec le milieu extérieur à travers leur membrane plasmique, selon le gradient de concentration en solutés (tonicité du milieu).

\item \textbf{Interprétation.}
\begin{itemize}
\item \emph{Tube A (\SI{9}{\gram\per\litre} NaCl)} : la solution est \emph{isotonique} par rapport au cytoplasme de l'hématie (c'est la concentration du sérum physiologique utilisé en perfusion). Les échanges d'eau sont équilibrés dans les deux sens (mouvements égaux entrée/sortie) : le volume et la forme discoïde biconcave des hématies sont maintenus.
\item \emph{Tube B (\SI{3}{\gram\per\litre} NaCl)} : la solution est \emph{hypotonique} (moins concentrée en solutés que le cytoplasme). L'eau entre massivement dans les hématies par \cle{osmose} (milieu extérieur hypotonique $\to$ milieu intérieur hypertonique). Ne possédant pas de paroi rigule (contrairement aux cellules végétales), les hématies gonflent jusqu'à la rupture de leur membrane : c'est l'\emph{hémolyse} (libération d'hémoglobine dans le milieu).
\item \emph{Tube C (\SI{15}{\gram\per\litre} NaCl)} : la solution est \emph{hypertonique}. L'eau sort des hématies par osmose (milieu intérieur hypotonique $\to$ milieu extérieur hypertonique) ; elles se déshydratent, leur membrane se replie en formant des excroissances : c'est la \emph{crénelation}.
\end{itemize}
\item \textbf{Conclusion.} La membrane plasmique de l'hématie se comporte comme une membrane \emph{semi-perméable} : à travers elle, l'eau se déplace par osmose du milieu le moins concentré (hypotonique) vers le milieu le plus concentré (hypertonique) en solutés, ce qui modifie le volume cellulaire. C'est pourquoi les perfusions hospitalières utilisent obligatoirement du sérum physiologique isotonique (\SI{9}{\gram\per\litre} NaCl, ou \SI{0,9}{\%} massique) pour éviter l'hémolyse des globules rouges du patient.
\end{enumerate}
\end{exoresolu}

\begin{experience}[Expérience de dialyse : séparation mélange amidon/glucose]
\textbf{Matériel.} Poche de dialyse (cellophane, membrane dialysante), amidon, glucose, eau distillée, solution d'iode (révélateur amidon), bandelettes réactives glucose (ou réactif de Fehling), deux béchers.
\textbf{Protocole.}
\begin{enumerate}
\item Préparer un mélange \SI{10}{mL} eau + \SI{1}{g} amidon + \SI{1}{g} glucose dans la poche de dialyse. Fermer hermétiquement.
\item Plonger la poche dans un bécher contenant \SI{200}{mL} d'eau distillée (milieu hypotonique pour les solutés).
\item Laisser reposer \SI{30}{min} sous agitation douce.
\item Prélèver \SI{2}{mL} de liquide \textbf{à l'intérieur} de la poche et \textbf{à l'extérieur} (bécher).
\item Tester l'amidon (goutte d'iode : coloration bleu-noir si présence) et le glucose (bandelette/Fehling : coloration rouge brique si présence).
\end{enumerate}
\textbf{Observations.}
\begin{itemize}
\item \emph{Intérieur de la poche :} Test iode $+$ (amidon présent), Test glucose $+$ (glucose présent, mais concentration diminuée).
\item \emph{Extérieur (bécher) :} Test iode $-$ (pas d'amidon), Test glucose $+$ (glucose présent).
\end{itemize}
\textbf{Interprétation.} La membrane de cellophane est \emph{dialysante} : elle laisse passer l'eau et les petites molécules (glucose, $M \approx \SI{180}{g/mol}$) mais retient les macromolécules (amidon, $M > \SI{10000}{g/mol}$). Le glucose a diffusé de la poche (concentré) vers le bécher (dilué) : c'est la \cle{dialyse}. L'amidon n'a pas traversé la membrane.
\textbf{Conclusion.} La dialyse permet de séparer les petites molécules des macromolécules en utilisant une membrane sélective. C'est le principe de l'hémodialyse (rein artificiel) pour épurer le sang des patients insuffisants rénaux (élimination urée, créatinine, excès d'ions ; rétention des protéines).
\end{experience}

\courssub{Savoir-faire 4 : Interpréter les résultats de l'expérience de Dutrochet}

\begin{methode}[Exploiter l'expérience de Dutrochet (osmomètre à membrane biologique)]
\begin{enumerate}
\item \textbf{Décrire le dispositif} : un osmomètre constitué d'une poche à membrane semi-perméable (historiquement : vessie de porc ; aujourd'hui : cellophane ou tube de dialyse) remplie d'une solution concentrée (ex. : eau sucrée \SI{1}{mol/L}), reliée à un tube vertical gradué (manomètre à colonne d'eau), le tout plongé dans un récipient d'eau pure (eau distillée).
\item \textbf{Observer} : le niveau du liquide monte dans le tube gradué ; la masse de l'ensemble augmente. La montée ralentit puis s'arrête.
\item \textbf{Interpréter} : la membrane laisse passer l'eau (solvant) mais retient le saccharose (soluté). L'eau pure (milieu \emph{hypotonique}, concentration nulle) entre dans la poche (milieu \emph{hypertonique}) : c'est l'\cle{osmose}. L'entrée d'eau augmente le volume et la pression hydrostatique dans la colonne. L'équilibre est atteint quand la \cle{pression hydrostatique} de la colonne d'eau ($P = \rho g h$) compense la \cle{pression osmotique} ($\pi$) de la solution.
\item \textbf{Conclure} : l'osmose est le passage de l'eau à travers une membrane semi-perméable, du milieu moins concentré vers le milieu plus concentré. La pression osmotique est la pression qu'il faut appliquer pour arrêter l'osmose.
\item \textbf{Vérifier les contrôles} : 
\begin{itemize}
\item Témoin 1 : Eau pure des deux côtés $\to$ aucune montée (pas de gradient de concentration).
\item Témoin 2 : Membrane perméable au saccharose (ex. : toile d'étamine) $\to$ le niveau monte puis redescend (diffusion du saccharose vers l'extérieur jusqu'à équilibre des concentrations).
\end{itemize}
\end{enumerate}
\end{methode}

\begin{popfigure}
\begin{tikzpicture}[scale=0.9, every node/.style={font=\small}]
% Cristallisoir eau pure
\draw[popdark, thick] (-4,-1) rectangle (-1,5);
\draw[popblue!20, fill=popblue!20] (-4,-1) rectangle (-1,5);
\node[popdark] at (-2.5,5.3) {Eau distillée (Hypotonique)};
% Osmomètre
\draw[popdark, thick] (-0.5,0) -- (-0.5,4) -- (0.5,4) -- (0.5,0); % Tube
\draw[popdark, thick] (-1.5,0) -- (-0.5,0) -- (-0.5,-1) -- (-1.5,-1) -- cycle; % Poche (simplifiée)
\draw[poporange, fill=poporange!30] (-1.5,0) rectangle (-0.5,-1); % Solution sucrée
\draw[popblue, fill=popblue!50] (-0.5,0) rectangle (0.5,3.2); % Montée eau dans tube
\node[popdark] at (0,3.5) {Niveau monté : $h$};
\node[popdark] at (0,-1.5) {Eau sucrée (Hypertonique)};
\node[popdark] at (0,4.3) {Tube gradué};
% Flèches osmose
\draw[->, popblue, very thick] (-2.5,-0.5) -- (-1.5,-0.5) node[midway, above] {Osose : $H_2O$};
% Légende
\node[popdark, align=left, anchor=west] at (1.5, 3) {Membrane semi-perméable};
\node[popdark, align=left, anchor=west] at (1.5, 2.5) {(Vessie porc / Cellophane)};
\node[popdark, align=left, anchor=west] at (1.5, 2.0) {Retient saccharose};
\node[popdark, align=left, anchor=west] at (1.5, 1.5) {Laisse passer $H_2O$};
\end{tikzpicture}
\legende{Osomètre de Dutrochet. L'eau pure (hypotonique) entre dans la poche contenant la solution sucrée (hypertonique) par osmose, faisant monter le niveau dans le tube. La hauteur $h$ à l'équilibre est proportionnelle à la pression osmotique $\pi = \rho g h$.}
\end{popfigure}

\begin{exoresolu}[L'osmomètre de Dutrochet]
\textbf{Énoncé.} Dutrochet réalise le montage suivant : une poche de vessie de porc (membrane semi-perméable) remplie d'eau sucrée est reliée à un tube vertical gradué, puis plongée dans un cristallisoir d'eau distillée. Au bout d'une heure, le niveau du liquide dans le tube a monté de 12 cm. Un dispositif témoin, contenant de l'eau distillée dans la poche, ne montre aucune variation de niveau.
\begin{enumerate}
\item Explique la montée du liquide dans le tube.
\item Précise le rôle du dispositif témoin.
\item Déduis-en une définition de l'osmose.
\end{enumerate}
\tcblower
\textbf{Solution détaillée.}
\begin{enumerate}
\item \textbf{Explication de la montée.} La vessie de porc est une membrane \emph{semi-perméable} : elle laisse passer les molécules d'eau (solvant) mais retient les molécules de saccharose (soluté). L'eau distillée du cristallisoir constitue le milieu \emph{hypotonique} (concentration en solutés nulle) et l'eau sucrée le milieu \emph{hypertonique}. L'eau traverse donc la membrane de l'extérieur vers l'intérieur de la poche : le volume de liquide augmente et le niveau monte de \SI{12}{cm} dans le tube. La montée cesse lorsque la pression hydrostatique exercée par la colonne d'eau ($P = \rho g h$) compense exactement la \cle{pression osmotique} ($\pi$) de la solution sucrée.
\item \textbf{Rôle du témoin.} Le dispositif témoin (eau distillée des deux côtés de la membrane) ne montre aucune montée : il prouve que le phénomène observé n'est pas dû au montage lui-même (poussée d'Archimède, capillarité, dilatation thermique) mais bien à la \emph{différence de concentration} en solutés de part et d'autre de la membrane semi-perméable. Il valide l'interprétation osmotique.
\item \textbf{Définition.} L'\cle{osmose} est le déplacement net du solvant (l'eau) à travers une membrane semi-perméable, du milieu le moins concentré (hypotonique) vers le milieu le plus concentré (hypertonique) en substances dissoutes (solutés).
\end{enumerate}
\textbf{Conclusion.} L'expérience de Dutrochet (1826) démontre quantitativement l'existence des échanges osmotiques d'eau à travers les membranes semi-perméables, mécanisme fondamental des échanges d'eau des cellules animales (réabsorption tubulaire rénale, résorption intestinale, équilibre hydrique).
\end{exoresolu}

\courssub{Savoir-faire 5 : Interpréter les résultats de l'expérience de Pfeffer}

\begin{methode}[Exploiter l'expérience de Pfeffer (mesure quantitative de la pression osmotique)]
\begin{enumerate}
\item \textbf{Décrire le dispositif} : Pfeffer (1877) utilise une membrane artificielle semi-perméable (précipité de \emph{ferrocyanure de cuivre} $Cu_2[Fe(CN)_6]$ déposé dans les parois d'un vase poreux en terre cuite) qui résiste à de fortes pressions. La poche contient la solution à étudier (ex. : saccharose) et plonge dans l'eau pure. Le dispositif est relié à un \emph{manomètre} (ou un tube à mercure) mesurant directement la pression.
\item \textbf{Observer et exploiter les données} : le manomètre indique une pression qui croît avec la concentration de la solution. La \cle{pression osmotique} $\pi$ (en atm ou Pa) est proportionnelle à la concentration molaire $C$ (en \si{mol\per\litre}) et à la température absolue $T$ (en kelvins) : c'est la \cle{loi de Van't Hoff} (1886) :
\[ \pi = R \times C \times T \]
avec $R = \SI{0,082}{L.atm.mol^{-1}.K^{-1}}$ (constante des gaz parfaits), $C$ en \si{mol\per\litre}, $T(\text{K}) = T(°\text{C}) + 273$.
\item \textbf{Interpréter} : la pression osmotique mesure la "force d'attraction" de l'eau par la solution ; plus une solution est concentrée en particules solutées, plus sa pression osmotique est élevée, plus elle attire l'eau.
\item \textbf{Conclure} en reliant au phénomène d'osmose : c'est la différence de pression osmotique ($\Delta \pi$) entre deux compartiments séparés par une membrane semi-perméable qui est la cause motrice du déplacement net d'eau.
\end{enumerate}
\textbf{Réflexe calcul :} Convertir systématiquement la température en kelvins ($T = t + 273$) et la concentration en \si{mol\per\litre} ($C = \frac{m}{M \times V}$) avant d'appliquer la formule $\pi = RCT$.
\end{methode}

\begin{popfigure}
\begin{tikzpicture}[scale=0.8, every node/.style={font=\small}]
% Vase poreux (terre cuite)
\draw[popmuted, thick, fill=popmuted!20] (-2,0) rectangle (2,4);
\draw[popdark, thick] (-2,0) -- (-2,4) (2,0) -- (2,4);
% Membrane ferrocyanure de cuivre (parois)
\draw[poppurple, thick, fill=poppurple!10] (-2,0) rectangle (-1.8,4);
\draw[poppurple, thick, fill=poppurple!10] (1.8,0) rectangle (2,4);
% Solution saccharose
\draw[poporange, fill=poporange!30] (-1.8,0) rectangle (1.8,3);
\node[popdark] at (0,3.3) {Solution saccharose $C$};
% Manomètre (tube U Hg)
\draw[popdark, thick] (2.5,0) -- (2.5,4) -- (4,4) -- (4,0);
\draw[popdark, fill=popdark!50] (2.5,0) rectangle (4,1.5); % Mercure bas
\draw[popdark, fill=popdark!50] (2.5,2.5) rectangle (4,4); % Mercure haut (pression)
\draw[<->, popdark, thick] (4.2,1.5) -- (4.2,2.5) node[midway, right] {$\Delta h$ (Pression $\pi$)};
\node[popdark] at (3.2,4.3) {Manomètre};
% Eau pure extérieure
\draw[popblue!20, fill=popblue!20] (-4,-1) rectangle (4,-0.5);
\draw[popblue!20, fill=popblue!20] (-4,4.5) rectangle (4,5);
\node[popdark] at (0,-0.7) {Eau pure (Hypotonique)};
\node[popdark] at (0,4.7) {Eau pure (Hypotonique)};
% Flèche osmose
\draw[->, popblue, very thick] (0,-0.5) -- (0,0) node[midway, right] {Osose $H_2O$};
\end{tikzpicture}
\legende{Appareil de Pfeffer pour la mesure de la pression osmotique. La membrane de ferrocyanure de cuivre (violette) résiste à la pression. Le manomètre indique la pression $\pi$ nécessaire pour empêcher l'entrée d'eau.}
\end{popfigure}

\begin{exoresolu}[Calcul de la pression osmotique d'une solution de saccharose]
\textbf{Énoncé.} Dans l'osmomètre de Pfeffer maintenu à \SI{27}{\celsius}, on introduit une solution contenant \SI{34,2}{g} de saccharose ($M = \SI{342}{g/mol}$) par litre. Le vase poreux plonge dans l'eau pure.
\begin{enumerate}
\item Calcule la concentration molaire de la solution.
\item Calcule la pression osmotique mesurée par le manomètre.
\item Que devient cette pression si l'on double la concentration ?
\end{enumerate}
\tcblower
\textbf{Solution détaillée.}
\begin{enumerate}
\item \textbf{Concentration molaire.}
\[ C = \frac{n}{V} = \frac{m}{M \times V} = \frac{34,2}{342 \times 1} = \SI{0,10}{mol\per\litre}. \]
\item \textbf{Pression osmotique.} Conversion de la température en kelvins : $T = 27 + 273 = \SI{300}{K}$. Application de la loi de Van't Hoff :
\[ \pi = R \times C \times T = 0,082 \times 0,10 \times 300 = \SI{2,46}{atm}. \]
La pression osmotique de la solution est de \SI{2,46}{atm} (soit environ \SI{2,49}{bar} ou \SI{249}{kPa}).
\item \textbf{Doublement de la concentration.} La pression osmotique étant proportionnelle à la concentration molaire ($\pi = R C T$), si $C$ double ($C' = \SI{0,20}{mol/L}$), $\pi$ double :
\[ \pi' = 2 \times 2,46 = \SI{4,92}{atm}. \]
\end{enumerate}
\textbf{Conclusion.} La pression osmotique croît proportionnellement à la concentration molaire en particules solutées (loi de Van't Hoff) : plus une solution est concentrée, plus elle attire l'eau à travers la membrane semi-perméable. L'expérience de Pfeffer a permis la première mesure quantitative rigoureuse de cette pression, validant l'analogie entre solutions diluées et gaz parfaits.
\end{exoresolu}

\courssub{Savoir-faire 6 : Interpréter les échanges de substances dissoutes et de particules}

\begin{definition}[Transport passif vs Transport actif]
\begin{itemize}
\item \cle{Transport passif} : déplacement spontané dans le sens du gradient de concentration (du concentré vers le dilué), \textbf{sans dépense d'énergie métabolique (ATP)}.
\begin{itemize}
\item \cle{Diffusion simple} : passage direct à travers la bicouche lipidique (petites molécules apolaires : \ce{O2}, \ce{CO2}, \ce{N2}, stéroïdes) ou par canaux aquaporines (\ce{H2O}). Aucune protéine transporteur spécifique saturable.
\item \cle{Diffusion facilitée} : passage via une protéine membranaire spécifique (canal ionique ou transporteur/perméase), saturable, stéréospécifique, \textbf{sans ATP}. Ex. : entrée du glucose dans l'hématie (GLUT1), entrée des \ce{K+} dans la cellule au repos.
\end{itemize}
\item \cle{Transport actif} : déplacement \textbf{contre le gradient de concentration} (du dilué vers le concentré), \textbf{nécessite une protéine transporteur (pompe) et de l'énergie (ATP direct ou indirect)}.
\begin{itemize}
\item \cle{Transport actif primaire} : hydrolyse directe d'ATP par la pompe (ex. : pompe \ce{Na+}/\ce{K+}-ATPase, pompe \ce{Ca^{2+}}-ATPase, pompe \ce{H+}-ATPase).
\item \cle{Transport actif secondaire (cotransport)} : utilisation de l'énergie du gradient électrochimique d'un ion (souvent \ce{Na+}) créé par une pompe primaire. \cle{Symport} (même sens, ex. : \ce{Na+}/\ce{glucose} SGLT1 intestin/rene) ou \cle{Antiport} (sens inverse, ex. : \ce{Na+}/\ce{H+}, \ce{Na+}/\ce{Ca^{2+}}).
\end{itemize}
\item \cle{Endocytose} / \cle{Exocytose} : transport de \emph{particules}, macromolécules ou gros volumes par formation/fusion de vésicules membranaires, \textbf{nécessite de l'ATP} (dynamine, actinomyosine, protéines SNARE). \cle{Phagocytose} ("manger", particules solides), \cle{Pinocytose} ("boire", fluide), \cle{Endocytose médiée par récepteurs} (spécifique, ex. : LDL, ferritine, hormones peptidiques).
\end{itemize}
\end{definition}

\begin{methode}[Distinguer et justifier les modes de transport membranaire]
\begin{enumerate}
\item \textbf{Identifier le mode de transport} à partir des indices de l'énoncé (courbe vitesse/concentration, effet inhibiteurs, sens du gradient) :
\begin{itemize}
\item \textbf{Courbe vitesse = f([S]) hyperbolique (saturation)} $\to$ Protéine transporteur/canal (Diffusion facilitée ou Transport actif).
\item \textbf{Courbe linéaire non saturante} $\to$ Diffusion simple (traversée lipidique).
\item \textbf{Transport contre le gradient} (concentration intracellulaire $>$ extracellulaire pour un nutriment) $\to$ Transport actif.
\item \textbf{Inhibiteur respiratoire (cyanure, DNP, ou absence de glucose/oxygène)} bloque le transport $\to$ Dépendance à l'ATP $\to$ Transport actif (ou endo/exocytose).
\item \textbf{Inhibiteur spécifique} (ex. : ouabaïne pour \ce{Na+}/\ce{K+}-ATPase, phlorizine pour SGLT1) $\to$ Identification de la pompe/transporteur.
\end{itemize}
\item \textbf{Argumenter rigoureusement} en citant les trois critères discriminants :
\begin{enumerate}
\item Sens par rapport au gradient de concentration (avec = passif / contre = actif).
\item Nécessité d'une protéine membranaire spécifique (saturation, inhibiteur compétitif).
\item Nécessité d'énergie métabolique (ATP) : testé par inhibiteurs de la phosphorylation oxydative (cyanure) ou glycolytique (iodoacétate).
\end{enumerate}
\end{enumerate}
\end{methode}

\begin{popfigure}
\begin{tikzpicture}[scale=0.85, every node/.style={font=\small}]
% Membrane
\draw[popdark, thick, fill=popcream] (-5,-0.5) rectangle (5,0.5);
\node[popdark] at (0,0.9) {\textbf{Membrane plasmique (Bicouche phospholipidique)}};
% Extracellulaire (haut)
\node[popblue] at (-4.5, 1.5) {Milieu extracellulaire};
% Intracellulaire (bas)
\node[poporange] at (-4.5, -1.5) {Cytoplasme};

% 1. Diffusion simple O2/CO2
\draw[->, popblue, thick] (-4.5, 1.2) -- (-4.5, 0.6);
\draw[->, poporange, thick] (-4.5, -0.6) -- (-4.5, -1.2);
\node[popdark, align=center] at (-4.5, 0) {\textbf{Diffusion simple}\\(\ce{O2}, \ce{CO2}, lipophiles)\\\textit{Sans protéine, sans ATP}};

% 2. Diffusion facilitée (Canal ionique)
\draw[popgreen, thick] (-1.5, 0.5) -- (-1.5, -0.5); % Canal
\draw[popgreen, fill=popgreen!30] (-1.5, 0) circle (0.15); % Pore
\draw[->, popblue] (-1.5, 1.2) -- (-1.5, 0.6);
\draw[->, popblue] (-1.5, -0.6) -- (-1.5, -1.2);
\node[popdark, align=center] at (-1.5, 0) {\textbf{Diffusion facilitée}\\(Canal ionique \ce{K+}, \ce{Na+}, \ce{Cl-})\\\textit{Protéine, saturable, sans ATP}};

% 3. Diffusion facilitée (Transporteur GLUT)
\draw[poppurple, thick, fill=poppurple!20] (1.5, 0.3) rectangle (2.5, -0.3); % Transporteur
\node at (2, 0) {GLUT};
\draw[->, popblue] (2, 1.2) -- (2, 0.4);
\draw[->, poporange] (2, -0.4) -- (2, -1.2);
\node[popdark, align=center] at (2, 0) {\textbf{Diffusion facilitée}\\(Transporteur GLUT1 glucose)\\\textit{Protéine, saturable, sans ATP}};

% 4. Transport actif primaire (Na/K ATPase)
\draw[poppink, thick, fill=poppink!20] (-4.5, -1.8) rectangle (-2.5, -2.8);
\node at (-3.5, -2.3) {\ce{Na+/K+}-ATPase};
\draw[->, popblue] (-3.5, -1.5) -- (-3.5, -1.8) node[midway, left] {3 \ce{Na+} (sortie)};
\draw[->, poporange] (-3.5, -2.8) -- (-3.5, -3.1) node[midway, left] {2 \ce{K+} (entrée)};
\node[rotate=90] at (-4.8, -2.3) {ATP $\to$ ADP + Pi};
\node[popdark, align=center] at (-3.5, -3.5) {\textbf{Transport actif primaire}\\\textit{Pompe, contre gradient, ATP direct}};

% 5. Transport actif secondaire (Symport Na/Glucose SGLT1)
\draw[popgold, thick, fill=popgold!20] (1.5, -1.8) rectangle (3.5, -2.8);
\node at (2.5, -2.3) {SGLT1};
\draw[->, popblue] (2.5, -1.5) -- (2.5, -1.8) node[midway, right] {\ce{Na+} (gradient)};
\draw[->, popblue] (2.5, -1.5) -- (2.5, -1.8) node[midway, left] {Glucose};
\node[popdark, align=center] at (2.5, -3.5) {\textbf{Transport actif secondaire}\\\textit{Symport, contre gradient glucose,}\\\textit{énergie gradient \ce{Na+}}};

% 6. Endocytose / Exocytose
\draw[->, popdark, thick, decorate, decoration={snake, amplitude=1pt}] (4.5, 0.5) .. controls (5.5, 0.5) and (5.5, -0.5) .. (4.5, -0.5);
\draw[popdark, fill=popdark!10] (4.8, 0.0) circle (0.25); % Vésicule
\node[popdark, align=center] at (5.5, 0) {\textbf{Endocytose / Exocytose}\\\textit{Vésicules, particules/macromolécules,}\\\textit{ATP (dynamine, SNARE)}};
\end{tikzpicture}
\legende{Résumé des principaux modes de transport transmembranaire. Flèches bleues = sens extracellulaire $\to$ intracellulaire (entrée) ; Flèches oranges = sortie. L'épaisseur des flèches ne représente pas la vitesse.}
\end{popfigure}

\begin{exoresolu}[Absorption du glucose par les cellules intestinales]
\textbf{Énoncé.} On mesure la vitesse d'absorption du glucose par des cellules de l'épithélium intestinal (bord en brosse) en fonction de la concentration extracellulaire en glucose, en l'absence puis en présence de cyanure (bloqueur de la chaîne respiratoire mitochondriale, arrêt de la production d'ATP). Résultats : sans cyanure, l'absorption se poursuit même lorsque la concentration intracellulaire dépasse la concentration extracellulaire ; avec cyanure, l'absorption s'arrête dès que les concentrations intra- et extracellulaires s'équilibrent.
\begin{enumerate}
\item Quel mode de transport ces résultats mettent-ils en évidence ? Justifie.
\item Pourquoi l'entrée du glucose dans les hématies, elle, n'est pas affectée par le cyanure ?
\end{enumerate}
\tcblower
\textbf{Solution détaillée.}
\begin{enumerate}
\item \textbf{Mode de transport mis en évidence : le transport actif (secondaire, symport \ce{Na+}/\ce{glucose} SGLT1).} Deux arguments expérimentaux majeurs :
\begin{itemize}
\item Le glucose continue d'être absorbé \emph{contre son gradient de concentration} (la concentration intracellulaire devient supérieure à la concentration extracellulaire) : seul un transport actif permet un tel déplacement "en montée".
\item Le cyanure, qui bloque la production d'ATP par phosphorylation oxydative, supprime totalement cette capacité d'absorption contre le gradient : le transport nécessite donc de l'\emph{énergie métabolique}. L'arrêt de l'absorption net dès l'équilibre des concentrations (avec cyanure) montre que seule la composante passive (diffusion facilitée résiduelle ou fuite) persiste sans ATP.
\end{itemize}
\textit{Précision mécanistique :} Il s'agit d'un transport actif \emph{secondaire}. L'énergie directe est fournie par le gradient électrochimique du \ce{Na+} (haut dehors, bas dedans), lui-même maintenu par la pompe \ce{Na+}/\ce{K+}-ATPase (transport actif primaire) qui consomme l'ATP. Le cyanure arrête la pompe $\to$ gradient \ce{Na+} s'effondre $\to$ symport SGLT1 s'arrête.

\item \textbf{Cas de l'hématie (érythrocyte).} Dans l'hématie, le glucose pénètre par \emph{diffusion facilitée} via le transporteur \cle{GLUT1} (uniporteur). Ce transport s'effectue \emph{dans le sens du gradient de concentration} : le glucose entrant est immédiatement phosphorylé en glucose-6-phosphate par l'hexokinase (glycolyse), ce qui maintient la concentration intracellulaire en glucose libre très faible. Ce mécanisme ne fait intervenir \textbf{aucune dépense d'ATP direct pour le transport lui-même} (l'ATP est consommé ensuite par l'hexokinase, mais le transporteur GLUT1 ne l'utilise pas). Le cyanure, en bloquant la respiration, n'affecte pas le transporteur GLUT1 (pas de saturation, pas de sens imposé contre gradient).
\end{enumerate}
\textbf{Conclusion.} Une même molécule (glucose) peut traverser la membrane plasmique par des mécanismes radicalement différents selon le type cellulaire et sa fonction physiologique : transport actif secondaire (absorption nutritionnelle intestinale/réabsorption rénale) vs diffusion facilitée (approvisionnement énergétique des cellules hématopoïétiques, cérébrales, musculaires). Ces échanges sont essentiels : ils assurent la nutrition cellulaire, l'élimination des déchets, le maintien de la composition ionique du milieu intérieur (homéostasie) et la signalisation cellulaire.
\end{exoresolu}

\courssub{Synthèse : Importance des échanges cellulaires dans la vie des animaux}

\begin{aretenir}[Rôles physiologiques majeurs des échanges cellulaires]
\begin{enumerate}
\item \textbf{Nutrition et approvisionnement énergétique :} Entrée du glucose (transport actif intestinal/rénal, diffusion facilitée hématies/cerveau), acides aminés, acides gras ; entrée \ce{O2} (diffusion simple) pour la respiration mitochondriale.
\item \textbf{Élimination des déchets métaboliques :} Sortie \ce{CO2} (diffusion simple), urée, acide urique, cratinine (diffusion facilitée/dialyse rénale) ; sécrétion active \ce{H+} (estomac, rein).
\item \textbf{Maintien de l'homéostasie ionique et osmotique (Milieu intérieur) :} Pompe \ce{Na+}/\ce{K+}-ATPase (potentiel de membrane, volume cellulaire, excitabilité) ; réabsorption \ce{Na+}, \ce{Cl-}, \ce{HCO3-}, \ce{Ca^{2+}}, \ce{Mg^{2+}}, eau (rein, intestin) ; régulation du pH (échangeurs \ce{Na+}/\ce{H+}, \ce{Cl-}/\ce{HCO3-}).
\item \textbf{Signalisation et communication intercellulaire :} \cle{Exocytose} des neurotransmetteurs (synapse), hormones (endocrine), facteurs de croissance, cytokines ; \cle{Endocytose médiée par récepteurs} pour l'internalisation des signaux (LDL, insuline, facteurs de croissance, antigènes) et la régulation de la sensibilité cellulaire (down-regulation).
\item \textbf{Défense immunitaire :} \cle{Phagocytose} (macrophages, neutrophiles) des bactéries/débris ; \cle{Exocytose} des granules cytotoxiques (perforine, granzymes) par les lymphocytes T CD8+ et NK ; présentation d'antigènes (endocytose $\to$ traitement $\to$ exocytose MHC-classe II).
\item \textbf{Spécialisation cellulaire (Différenciation) :} L'expression différentielle des protéines de transport (canaux, pompes, transporteurs, récepteurs) définit le phénotype fonctionnel de la cellule (ex. : cellule principale du tubule collecteur rénal : aquaporines-2 régulées par ADH ; cellule bêta pancréatique : GLUT2, glucokinase, canaux \ce{K+}$_{ATP}$, vésicules d'insuline).
\end{enumerate}
\end{aretenir}

\begin{attention}[Les 5 erreurs qui coûtent des points au Bac]
\begin{enumerate}
\item \textbf{Confondre le sens du mouvement net d'eau.} L'eau se déplace du milieu \emph{hypotonique} (peu concentré en solutés, potentiel d'eau élevé) vers le milieu \emph{hypertonique} (très concentré, potentiel d'eau bas), jamais l'inverse. \textbf{Astuce mnémotechnique :} « L'eau va vers le sel » (ou vers le plus concentré). Dire « l'eau sort de la solution concentrée » ou « l'eau suit le gradient de concentration de l'eau » sans préciser le sens est une erreur éliminatoire.
\item \textbf{Confondre osmose et dialyse.} L'\cle{osmose} concerne le déplacement de l'\emph{eau (solvant)} à travers une membrane \emph{semi-perméable} (imperméable aux solutés) ; la \cle{dialyse} concerne le déplacement des \emph{solutés (petites molécules dissoutes)} à travers une membrane \emph{dialysante} (perméable aux petites molécules, imperméable aux macromolécules), du milieu concentré vers le milieu dilué. Ne jamais employer l'un pour l'autre.
\item \textbf{Oublier les conversions d'unités dans les calculs.} Pour le grossissement ($G = \frac{\text{image}}{\text{réel}}$), convertir les deux longueurs dans la \textbf{même unité} (1 cm $= \SI{10000}{\micro\metre}$ ; 1 mm $= \SI{1000}{\micro\metre}$). Pour la pression osmotique ($\pi = R C T$), convertir \textbf{impérativement} la température en kelvins ($T = t(°\text{C}) + 273$) et la concentration en \si{mol\per\litre}. Utiliser $t = 27$ au lieu de $T = 300$ K donne un résultat faux d'un facteur $\approx 11$.
\item \textbf{Affirmer que la diffusion facilitée consomme de l'ATP.} La diffusion (simple ou facilitée) est toujours \emph{passive} : elle suit le gradient de concentration chimique/électrochimique et ne consomme \textbf{aucune énergie métabolique} (ATP). La saturation de la courbe vitesse = f([S]) prouve l'intervention d'une protéine, pas l'usage d'ATP. Seuls le transport actif (primaire ou secondaire), l'endocytose et l'exocytose consomment de l'ATP.
\item \textbf{Bâcler le schéma biologique et la conclusion.} Un schéma sans titre souligné, sans indication du grossissement/type de microscopie, avec des lignes d'attache qui se croisent, qui ne sont pas horizontales, ou avec des hachures/ombres perd des points faciles. De même, une conclusion qui se contente de recopier les résultats bruts (« le niveau a monté de 12 cm », « les hématies ont éclaté ») au lieu de dégager le mécanisme général (« l'eau passe par osmose du milieu hypotonique vers le milieu hypertonique à travers une membrane semi-perméable ») ne répond pas au problème posé et ne vaut pas les points de « Conclure ».
\end{enumerate}
\end{attention}

\begin{savaistu}[Histoire des sciences \& Applications camerounaises]
\textbf{Histoire.} Henri \textsc{Dutrochet} (1776-1847), médecin et botaniste français, découvre l'osmose en 1826 avec son osmomètre à vessie de porc. Il invente le terme « osmose » (du grec \emph{osmos}, poussée). Wilhelm \textsc{Pfeffer} (1845-1920) met au point en 1877 l'osmomètre à membrane de ferrocyanure de cuivre résistant à la pression, permettant la mesure quantitative de la pression osmotique et la validation de la loi de Van't Hoff ($\pi = CRT$), étendant les lois des gaz parfaits aux solutions diluées.

\textbf{Applications au Cameroun.}
\begin{itemize}
\item \textbf{Santé publique :} La préparation rigoureuse du \textbf{SRO (Sérum de Réhydratation Orale)} pour le traitement de la diarrhée aiguë (choléra, rotavirus) repose sur la compréhension du transport actif couplé \ce{Na+}/\ce{glucose} (SGLT1) au niveau de l'intestin grêle : le glucose favorise l'absorption du \ce{Na+} et de l'eau par osmose, sauvant des milliers d'enfants chaque année. Les perfusions intraveineuses (sérum physiologique \SI{0,9}{\%} NaCl, Ringer lactate) doivent être isotoniques pour éviter l'hémolyse.
\item \textbf{Insuffisance rénale :} Les centres d'hémodialyse (Yaoundé, Douala, Garoua) utilisent le principe de la \textbf{dialyse} à travers une membrane semi-perméable artificielle pour épurer le sang des déchets azotés (urée, créatinine) et corriger les désordres hydro-électrolytiques, en attendant une éventuelle greffe.
\item \textbf{Agriculture :} La tolérance à la salinité des cultures (riz dans la vallée du Logone, coton au Nord) dépend de la capacité des cellules racinaires à réguler leur pression osmotique (accumulation de prolines, sucres, ions \ce{K+}) pour maintenir l'entrée d'eau (osmose) malgré un sol hypertonique.
\item \textbf{Environnement :} L'éruption limnique du \textbf{Lac Nyos} (1986) a libéré un nuage de \ce{CO2} dense. La diffusion simple du \ce{CO2} à travers les membranes alvéolaires (gradient de pression partielle) a causé l'asphyxie massive. La compréhension des échanges gazeux (diffusion simple \ce{O2}/\ce{CO2}) est cruciale pour la réanimation.
\end{itemize}
\end{savaistu}

\newpage

\coursec{Exercices}

\courssub{Je vérifie mes connaissances}

\begin{exercice}[QCM : la cellule et ses organites]
Pour chaque question, une seule réponse est exacte. Entoure la lettre correspondante.

\textbf{1.} L'organite exclusivement visible au microscope électronique est :
\begin{itemize}
\item[a)] le noyau
\item[b)] la mitochondrie
\item[c)] le réticulum endoplasmique
\item[d)] la membrane plasmique
\end{itemize}

\textbf{2.} Les ribosomes sont le siège :
\begin{itemize}
\item[a)] de la respiration cellulaire
\item[b)] de la synthèse des protéines
\item[c)] de la digestion intracellulaire
\item[d)] de la photosynthèse
\end{itemize}

\textbf{3.} L'appareil de Golgi intervient dans :
\begin{itemize}
\item[a)] la production d'ATP
\item[b)] la réplication de l'ADN
\item[c)] la maturation et la sécrétion des protéines
\item[d)] la dégradation des déchets cellulaires
\end{itemize}

\textbf{4.} Le transport actif se caractérise par :
\begin{itemize}
\item[a)] un déplacement dans le sens du gradient de concentration
\item[b)] une consommation d'ATP et l'intervention de protéines transporteuses
\item[c)] l'absence totale de protéines membranaires
\item[d)] un déplacement uniquement des molécules d'eau
\end{itemize}
\end{exercice}

\begin{exercice}[Vrai ou faux justifié]
Réponds par \textbf{vrai} ou \textbf{faux}, puis justifie ta réponse en une ou deux phrases.
\begin{enumerate}
\item Une hématie placée dans l'eau distillée se rétracte.
\item La diffusion simple nécessite une dépense d'énergie par la cellule.
\item Les lysosomes contiennent des enzymes digestives.
\item L'osmose est un déplacement de solvant à travers une membrane semi-perméable, du milieu le moins concentré vers le milieu le plus concentré en soluté.
\item L'exocytose permet à une cellule d'expulser des substances vers le milieu extracellulaire.
\item La pompe Na$^+$/K$^+$ fonctionne sans ATP.
\end{enumerate}
\end{exercice}

\begin{exercice}[Définitions]
Définis brièvement et précisément chacun des termes suivants :
\begin{enumerate}
\item Osmose ;
\item Dialyse ;
\item Diffusion facilitée ;
\item Endocytose ;
\item Solution isotonique.
\end{enumerate}
\end{exercice}

\begin{exercice}[Calculs directs : pression osmotique]
On rappelle la loi de Van't Hoff : $\pi = R \times T \times C$, avec $R = \num{0,082}$~L$\cdot$atm$\cdot$K$^{-1}\cdot$mol$^{-1}$, $T$ la température absolue en kelvins (K) et $C$ la concentration molaire du soluté en mol$\cdot$L$^{-1}$.
\begin{enumerate}
\item Calcule la pression osmotique d'une solution de saccharose de concentration $C = \num{0,30}$~mol$\cdot$L$^{-1}$ à $T = 300$~K.
\item Calcule la pression osmotique d'une solution de glucose de concentration $C = \num{0,15}$~mol$\cdot$L$^{-1}$ à la même température.
\item Laquelle de ces deux solutions est hypertonique par rapport à l'autre ? Justifie.
\end{enumerate}
\end{exercice}

\courssub{Je m'entraîne}

\begin{exercice}[Schéma d'une cellule animale]
Un élève de Terminale C du lycée de Ngaoundéré observe au microscope électronique une cellule de pancréas.
\begin{enumerate}
\item Réalise un schéma annoté de cette cellule animale telle qu'elle apparaît en microscopie électronique, en y représentant au moins huit structures : membrane plasmique, cytoplasme, noyau (avec nucléole et membrane nucléaire), mitochondries, réticulum endoplasmique granuleux, réticulum endoplasmique lisse, appareil de Golgi, ribosomes.
\item Indique sur ton schéma, par une flèche, le trajet d'une protéine de sécrétion depuis sa synthèse jusqu'à sa sortie de la cellule.
\item Nomme le mécanisme par lequel cette protéine quitte la cellule.
\end{enumerate}
\end{exercice}

\begin{exercice}[Expérience de Dutrochet]
Un osmomètre de Dutrochet est constitué d'une cloche dont l'ouverture est fermée par une membrane de vessie de porc (membrane semi-perméable) et surmontée d'un tube capillaire. La cloche contient une solution de saccharose à \num{58,5}~g$\cdot$L$^{-1}$ ; elle est plongée dans un cristallisoir rempli d'eau distillée, à la température de 25~\textdegree C.
\begin{enumerate}
\item Schématise le dispositif expérimental au début de l'expérience.
\item Prédis et justifie l'évolution du niveau du liquide dans le tube capillaire.
\item Explique, en termes de mouvements moléculaires, le phénomène observé.
\item Pourquoi la montée du liquide s'arrête-t-elle au bout d'un certain temps ?
\item On recommence l'expérience avec une solution de saccharose deux fois plus concentrée. Compare qualitativement les hauteurs finales atteintes dans les deux cas.
\end{enumerate}
\end{exercice}

\begin{exercice}[Hématies et tonicité des milieux]
Au laboratoire du Centre Hospitalier d'Essos (Yaoundé), un technicien place des hématies humaines dans trois solutions de chlorure de sodium (NaCl) :
\begin{itemize}
\item solution A : NaCl à \num{0,3}~\% (masse/volume) ;
\item solution B : NaCl à \num{0,9}~\% ;
\item solution C : NaCl à \num{3}~\%.
\end{itemize}
Après 15 minutes, il observe les cellules au microscope optique.
\begin{enumerate}
\item Prédis l'aspect des hématies dans chacune des trois solutions et schématise-les.
\item Qualifie chaque solution (hypotonique, isotonique ou hypertonique) par rapport au cytoplasme des hématies.
\item Explique les modifications observées dans les solutions A et C en précisant le sens des mouvements d'eau.
\item Pourquoi le sérum physiologique injecté aux malades est-il une solution de NaCl à \num{0,9}~\% ?
\end{enumerate}
\end{exercice}

\begin{exercice}[Tableau comparatif des échanges cellulaires]
Recopie et complète le tableau suivant, qui compare les différents modes d'échanges transmembranaires :
\begin{center}
\begin{tabularx}{\linewidth}{|l|X|X|X|X|}
\hline
\rowcolor{popblueL} \textbf{Mode de transport} & \textbf{Substance transportée} & \textbf{Sens du déplacement} & \textbf{Énergie (ATP)} & \textbf{Exemple physiologique} \\
\hline
Osmose & & & & \\
\hline
Diffusion simple & & & & \\
\hline
Diffusion facilitée & & & & \\
\hline
Transport actif & & & & \\
\hline
Endocytose & & & & \\
\hline
Exocytose & & & & \\
\hline
\end{tabularx}
\end{center}
\end{exercice}

\courssub{Je me prépare au Bac}

\begin{exercice}[Type Bac : l'expérience de Pfeiffer et les facteurs de l'osmose]
Au cours d'une séance de travaux pratiques, un groupe d'élèves de Terminale TI réalise l'expérience de Pfeiffer : un osmomètre à membrane de ferrocyanure de cuivre (imperméable au saccharose, perméable à l'eau) contient des solutions de saccharose de concentrations croissantes ; il est plongé dans l'eau distillée à 20~\textdegree C. Le manomètre associé permet de mesurer la pression osmotique $\pi$ de chaque solution. Les résultats sont consignés dans le tableau ci-dessous :

\begin{center}
\begin{tabularx}{\linewidth}{|l|*{6}{X|}}
\hline
\rowcolor{popblueL} \textbf{Concentration $C$ (g$\cdot$L$^{-1}$)} & 0 & 17,1 & 34,2 & 51,3 & 68,4 & 85,5 \\
\hline
\textbf{Pression osmotique $\pi$ (atm)} & 0 & 1,2 & 2,4 & 3,6 & 4,8 & 6,0 \\
\hline
\end{tabularx}
\end{center}

Masse molaire du saccharose : $M = 342$~g$\cdot$mol$^{-1}$.

\begin{enumerate}
\item Rappelle les conditions nécessaires à la réalisation de l'osmose.
\item Trace sur papier millimétré la courbe $\pi = f(C)$. Échelles : 1~cm pour 10~g$\cdot$L$^{-1}$ en abscisse ; 1~cm pour 0,5~atm en ordonnée.
\item Déduis de l'allure de la courbe la relation mathématique entre la pression osmotique et la concentration.
\item Convertis les concentrations massiques en concentrations molaires, puis calcule le rapport $\dfrac{\pi}{C_{\text{mol}}}$ pour chaque solution. Que constates-tu ? En déduire la valeur du produit $R \times T$ à 20~\textdegree C (293~K) et comparer à la valeur théorique $R \times T = \num{0,082} \times 293$.
\item Une solution X de saccharose de concentration inconnue développe une pression osmotique de 3,0~atm. Détermine graphiquement puis par le calcul sa concentration massique.
\item Conclus sur les facteurs dont dépend la pression osmotique d'une solution.
\end{enumerate}
\end{exercice}

\begin{exercice}[Type Bac : absorption du potassium par des cellules racinaires]
Dans une station de recherche agronomique de Dschang, on étudie l'absorption des ions K$^+$ par des cellules de racines de maïs placées dans une solution nutritive contenant \num{0,5}~mmol$\cdot$L$^{-1}$ de K$^+$. On mesure la quantité de K$^+$ absorbée (en $\mu$mol par gramme de matière fraîche) en fonction du temps, dans trois conditions expérimentales :

\begin{itemize}
\item \textbf{Lot 1 (témoin)} : racines à 25~\textdegree C, en présence d'air ;
\item \textbf{Lot 2} : racines à 25~\textdegree C, en présence de cyanure de potassium (KCN), inhibiteur de la respiration cellulaire ;
\item \textbf{Lot 3} : racines à 4~\textdegree C, en présence d'air.
\end{itemize}

\begin{center}
\begin{tabularx}{\linewidth}{|l|*{6}{X|}}
\hline
\rowcolor{popblueL} \textbf{Temps (min)} & 0 & 20 & 40 & 60 & 90 & 120 \\
\hline
\textbf{Lot 1 : K$^+$ absorbé ($\mu$mol$\cdot$g$^{-1}$)} & 0 & 4,0 & 7,5 & 10,0 & 12,5 & 14,0 \\
\hline
\textbf{Lot 2 : K$^+$ absorbé ($\mu$mol$\cdot$g$^{-1}$)} & 0 & 1,2 & 1,6 & 1,7 & 1,7 & 1,7 \\
\hline
\textbf{Lot 3 : K$^+$ absorbé ($\mu$mol$\cdot$g$^{-1}$)} & 0 & 1,0 & 1,8 & 2,3 & 2,8 & 3,0 \\
\hline
\end{tabularx}
\end{center}

On précise que la concentration intracellulaire en K$^+$ des cellules racinaires est de \num{80}~mmol$\cdot$L$^{-1}$.

\begin{enumerate}
\item Trace sur un même graphique les trois courbes d'absorption du K$^+$ en fonction du temps. Échelles : 1~cm pour 20~min ; 1~cm pour 2~$\mu$mol$\cdot$g$^{-1}$.
\item Compare les courbes des lots 1 et 2. Quel est le rôle du cyanure ? Que révèle son effet sur la nature de l'absorption du K$^+$ ?
\item Compare les courbes des lots 1 et 3. Quel facteur est mis en évidence et comment agit-il ?
\item En tenant compte des concentrations intra- et extracellulaires en K$^+$, montre que l'absorption du K$^+$ s'effectue contre le gradient de concentration.
\item Identifie le type de transport mis en jeu et dégage les trois conditions nécessaires à ce transport.
\item La faible absorption résiduelle observée dans le lot 2 correspond à un autre type de transport : lequel ? Justifie.
\end{enumerate}
\end{exercice}

\begin{exercice}[Type Bac : phagocytose et défense de l'organisme]
Au cours d'une infection à \emph{Plasmodium} (paludisme) suivie à l'Hôpital Central de Yaoundé, on observe au microscope électronique à transmission un macrophage (leucocyte) en interaction avec une bactérie. Un cliché montre successivement : (a) la bactérie accolée à la surface de la cellule ; (b) l'émission par le macrophage de prolongements cytoplasmiques qui entourent la bactérie ; (c) la bactérie enfermée dans une vésicule intracytoplasmique ; (d) la fusion de cette vésicule avec d'autres vésicules riches en enzymes, puis la disparition progressive de la structure bactérienne.

\begin{enumerate}
\item Nomme le phénomène général illustré par ce cliché et précise la catégorie d'endocytose à laquelle il appartient.
\item Décris, en quatre étapes numérotées, le déroulement du phénomène observé, en utilisant le vocabulaire scientifique approprié (pseudopodes, phagosome, phagolysosome).
\item Schématise l'étape (c) en annotant : membrane plasmique, cytoplasme, noyau, vésicule contenant la bactérie.
\item Nomme l'organite dont proviennent les enzymes responsables de la digestion de la bactérie et précise l'origine de ces enzymes (organite de synthèse).
\item La phagocytose est un processus qui consomme de l'énergie : cite l'organite qui la fournit et justifie son abondance dans les macrophages.
\item Les résidus de digestion sont ensuite rejetés hors de la cellule : par quel mécanisme ? Définis-le brièvement.
\item Déduis de cette étude deux arguments montrant l'importance des échanges cellulaires dans la vie des animaux.
\end{enumerate}
\end{exercice}

\newpage

\coursec{Corrigés des exercices}

\begin{corrige}[Exercice 1 — QCM : la cellule et ses organites]
\textbf{1. Réponse c)} Le réticulum endoplasmique, dont la membrane mesure environ 5~nm d'épaisseur, est invisible au microscope optique (pouvoir de résolution $\approx \num{0,2}~\mu$m). Le noyau, la mitochondrie (sous forme de granulations) et la membrane plasmique (limite de la cellule) sont observables en microscopie optique.

\textbf{2. Réponse b)} Les ribosomes assurent la traduction de l'ARN messager : c'est le siège de la \cle{synthèse des protéines}. La respiration cellulaire a lieu dans les mitochondries et la digestion intracellulaire dans les lysosomes.

\textbf{3. Réponse c)} L'appareil de Golgi reçoit les protéines synthétisées dans le réticulum endoplasmique granuleux, les modifie (glycosylation), les trie et les conditionne dans des vésicules de sécrétion.

\textbf{4. Réponse b)} Le transport actif s'effectue \emph{contre} le gradient de concentration ; il exige donc une dépense d'énergie (hydrolyse de l'ATP) et l'intervention de protéines membranaires transporteuses (pompes).
\end{corrige}

\begin{corrige}[Exercice 2 — Vrai ou faux justifié]
\begin{enumerate}
\item \textbf{Faux.} L'eau distillée est hypotonique par rapport au cytoplasme de l'hématie : l'eau entre dans la cellule par osmose, l'hématie gonfle puis éclate (hémolyse). C'est en milieu hypertonique qu'elle se rétracte (crénation).
\item \textbf{Faux.} La diffusion simple est un transport passif : elle s'effectue dans le sens du gradient électrochimique, sans dépense d'énergie par la cellule.
\item \textbf{Vrai.} Les lysosomes sont des vésicules contenant des enzymes hydrolytiques (hydrolases acides) qui digèrent les macromolécules et les structures étrangères.
\item \textbf{Vrai.} C'est la définition même de l'osmose : le solvant (l'eau) traverse la membrane semi-perméable du compartiment le moins concentré (hypotonique) vers le plus concentré (hypertonique) en soluté.
\item \textbf{Vrai.} L'exocytose est l'expulsion de substances enfermées dans une vésicule qui fusionne avec la membrane plasmique et déverse son contenu dans le milieu extracellulaire.
\item \textbf{Faux.} La pompe Na$^+$/K$^+$ est une ATPase : elle hydrolyse l'ATP (3 Na$^+$ expulsés contre 2 K$^+$ absorbés par ATP hydrolysé). C'est un transport actif.
\end{enumerate}
\end{corrige}

\begin{corrige}[Exercice 3 — Définitions]
\begin{enumerate}
\item \textbf{Osmose} : déplacement de molécules de solvant (eau) à travers une membrane semi-perméable, du milieu le moins concentré vers le milieu le plus concentré en soluté, jusqu'à l'équilibre.
\item \textbf{Dialyse} : diffusion des petites molécules dissoutes (cristalloïdes) à travers une membrane perméable, du milieu le plus concentré vers le milieu le moins concentré, les grosses molécules (colloïdes) étant retenues.
\item \textbf{Diffusion facilitée} : transport passif (sans ATP) d'une substance dans le sens de son gradient de concentration, grâce à une protéine membranaire spécifique (canal ou transporteur).
\item \textbf{Endocytose} : entrée de particules ou de macromolécules dans la cellule par invagination de la membrane plasmique, qui forme une vésicule intracytoplasmique ; phénomène actif consommant de l'ATP.
\item \textbf{Solution isotonique} : solution dont la concentration en solutés (et donc la pression osmotique) est égale à celle du milieu intracellulaire ; elle ne provoque aucun mouvement net d'eau (ex. NaCl à \num{0,9}~\% pour les hématies).
\end{enumerate}
\end{corrige}

\begin{corrige}[Exercice 4 — Calculs directs : pression osmotique]
\begin{enumerate}
\item Solution de saccharose : $\pi = R \times T \times C = \num{0,082} \times 300 \times \num{0,30} = \num{7,38}$~atm, soit $\pi \approx \num{7,4}$~atm.
\item Solution de glucose : $\pi = \num{0,082} \times 300 \times \num{0,15} = \num{3,69}$~atm, soit $\pi \approx \num{3,7}$~atm.
\item La solution de saccharose (\num{7,38}~atm) est \textbf{hypertonique} par rapport à la solution de glucose (\num{3,69}~atm), car sa concentration molaire est plus élevée ($\num{0,30} > \num{0,15}$~mol$\cdot$L$^{-1}$) : à température égale, la pression osmotique est proportionnelle à la concentration molaire.
\end{enumerate}
\end{corrige}

\begin{attention}[Points de vigilance]
Ne pas oublier de convertir la température en kelvins ($T(\text{K}) = t(^\circ\text{C}) + 273$) et d'utiliser la concentration \emph{molaire} (mol$\cdot$L$^{-1}$), jamais la concentration massique, dans la loi de Van't Hoff.
\end{attention}

\begin{corrige}[Exercice 5 — Schéma d'une cellule animale]
\begin{enumerate}
\item Le schéma doit présenter une cellule de forme arrondie, sans paroi ni chloroplaste, avec les huit structures demandées, dessinées à l'échelle relative correcte (noyau volumineux, organites ponctués dans le cytoplasme) et légendées par des traits horizontaux.

\begin{popfigure}
\begin{tikzpicture}[scale=1, font=\small]
\draw[fill=popcream, draw=popdark, thick] (0,0) ellipse (4.6 and 3.2);
\draw[draw=popink, thick] (0,0) ellipse (4.75 and 3.35);
\draw[fill=popblueL, draw=popblue, thick] (-0.8,0.2) circle (1.25);
\draw[fill=popblue, draw=popblue] (-0.8,0.2) circle (0.35);
\draw[fill=poporangeL, draw=poporange] (2.6,1.4) ellipse (0.75 and 0.35);
\draw[draw=poporange] (2.6,1.4) ellipse (0.55 and 0.22);
\draw[fill=poporangeL, draw=poporange] (-2.9,-1.6) ellipse (0.7 and 0.32);
\draw[draw=poporange] (-2.9,-1.6) ellipse (0.5 and 0.2);
\draw[fill=popgreenL, draw=popgreen, thick] (2.4,-1.5) ellipse (0.9 and 0.55);
\draw[draw=popgreen] (1.75,-1.5) -- (3.05,-1.5);
\draw[draw=popgreen] (1.85,-1.3) -- (2.95,-1.3);
\draw[draw=popgreen] (1.85,-1.7) -- (2.95,-1.7);
\foreach \y in {0.9,1.15,1.4,1.65} \draw[draw=poppurple] (-3.6,\y) -- (-1.9,\y);
\foreach \x in {-3.4,-3.0,-2.6,-2.2} \foreach \y in {0.9,1.15,1.4,1.65} \fill[poppurple] (\x,\y) circle (0.045);
\draw[draw=popgold] (0.9,2.1) arc (200:340:0.7);
\draw[draw=popgold] (0.9,2.3) arc (200:340:0.85);
\foreach \p in {(1.2,-2.3),(0.4,1.9),(-2.2,0.9),(3.4,0.3)} \fill[popdark] \p circle (0.06);
\draw[->, thick, poppink] (-1.9,1.3) -- (0.2,2.0);
\draw[->, thick, poppink] (0.2,2.0) -- (1.9,2.15);
\draw[->, thick, poppink] (1.9,2.15) -- (3.6,2.6);
\draw[->, thick, poppink] (3.6,2.6) -- (4.9,2.75);
\legende{Schéma d'une cellule animale (cellule pancréatique) en microscopie électronique : 1. membrane plasmique (double contour) ; 2. cytoplasme (fond clair) ; 3. noyau délimité par la membrane nucléaire ; 4. nucléole ; 5. mitochondries (ovales à crêtes internes) ; 6. appareil de Golgi (saccules empilés, à droite en bas) ; 7. réticulum endoplasmique granuleux (canaux parsemés de ribosomes, en haut à gauche) ; 8. réticulum endoplasmique lisse (canaux lisses, en haut) ; 9. ribosomes libres (points noirs). Flèches : trajet d'une protéine de sécrétion.}
\end{tikzpicture}
\end{popfigure}

\item Le trajet de la protéine de sécrétion (flèches sur le schéma) : \textbf{ribosomes du REG} $\rightarrow$ \textbf{lumière du réticulum endoplasmique granuleux} $\rightarrow$ \textbf{appareil de Golgi} (maturation, conditionnement) $\rightarrow$ \textbf{vésicule de sécrétion} $\rightarrow$ \textbf{fusion avec la membrane plasmique} $\rightarrow$ milieu extracellulaire.
\item La protéine quitte la cellule par \cle{exocytose}.
\end{enumerate}
\end{corrige}

\begin{attention}[Points de vigilance]
Une cellule animale ne possède \textbf{ni paroi, ni chloroplaste, ni grande vacuole centrale} : leur présence sur le schéma est une erreur éliminatoire au Bac. Les légendes doivent être écrites horizontalement, reliées aux structures par des traits fins qui ne se croisent pas.
\end{attention}

\begin{corrige}[Exercice 6 — Expérience de Dutrochet]
\begin{enumerate}
\item Schéma du dispositif :

\begin{popfigure}
\begin{tikzpicture}[font=\small]
\draw[thick] (0,0) -- (0,4.5) (1.2,0) -- (1.2,4.5);
\draw[thick, fill=popblueL] (0,0) -- (0,2.6) -- (1.2,2.6) -- (1.2,0);
\draw[thick, fill=popblueL] (-0.9,-2.2) -- (-0.9,0) -- (2.1,0) -- (2.1,-2.2) -- cycle;
\draw[thick, poporange, pattern=north east lines, pattern color=poporange] (-0.9,-2.2) -- (2.1,-2.2) -- (2.1,-2.45) -- (-0.9,-2.45) -- cycle;
\draw[thick] (-2.6,-3.1) -- (-2.6,-1.2) -- (3.8,-1.2) -- (3.8,-3.1);
\draw[fill=popgreenL, opacity=0.6] (-2.6,-3.1) rectangle (3.8,-1.5);
\draw[->, popink, thick] (-2.2,-1.7) -- (-1.4,-2.15);
\draw[->, popink, thick] (3.4,-1.7) -- (2.6,-2.15);
\node at (6.2,3.4) {tube capillaire};
\draw[->] (5.2,3.3) -- (0.9,3.3);
\node at (6.4,1.9) {solution de saccharose};
\draw[->] (5.0,1.8) -- (0.7,1.4);
\node at (6.3,-2.35) {membrane de vessie};
\draw[->] (4.9,-2.3) -- (2.2,-2.3);
\node at (-4.3,-1.75) {eau distillée};
\draw[->] (-3.2,-1.75) -- (-2.5,-1.9);
\legende{Osmomètre de Dutrochet : la cloche contenant la solution de saccharose est fermée par une membrane semi-perméable (perméable à l'eau, imperméable au saccharose) et plongée dans l'eau distillée.}
\end{tikzpicture}
\end{popfigure}

\item Le niveau du liquide \textbf{monte} dans le tube capillaire. En effet, la membrane de vessie est perméable à l'eau mais pas au saccharose : l'eau distillée (milieu hypotonique) pénètre dans la cloche (milieu hypertonique) par osmose, ce qui augmente le volume de liquide dans la cloche.
\item Les molécules d'eau, animées d'un mouvement brownien, frappent la membrane des deux côtés ; mais du côté de la solution, une partie des molécules d'eau est retenue par hydratation autour des molécules de saccharose. Le flux d'eau est donc plus important de l'eau distillée vers la solution que dans le sens inverse : il en résulte un \textbf{flux net d'eau vers la cloche}.
\item La montée s'arrête lorsque la \textbf{pression hydrostatique} exercée par la colonne de liquide dans le tube devient égale à la \textbf{pression osmotique} de la solution : les deux flux d'eau s'équilibrent alors (équilibre dynamique).
\item La pression osmotique étant proportionnelle à la concentration (loi de Van't Hoff), une solution deux fois plus concentrée développe une pression osmotique deux fois plus grande : la hauteur finale de la colonne sera environ \textbf{deux fois plus élevée} (et la montée initiale plus rapide).
\end{enumerate}
\end{corrige}

\begin{corrige}[Exercice 7 — Hématies et tonicité des milieux]
\begin{enumerate}
\item Aspect des hématies :

\begin{popfigure}
\begin{tikzpicture}[font=\small]
\draw[draw=poppink, thick, fill=poppinkL!10] (0,0) circle (1.0);
\node at (0,-1.5) {A : gonflée, éclatée};
\node at (0,-1.95) {(hémolyse)};
\draw[fill=poppinkL, draw=poppink, thick] (4,0) circle (0.9);
\draw[fill=white, draw=poppink] (4,0) circle (0.35);
\node at (4,-1.5) {B : normale};
\node at (4,-1.95) {(disque biconcave)};
\draw[fill=poppinkL, draw=poppink, thick, decorate, decoration={zigzag, segment length=4, amplitude=1.2}] (8,0) circle (0.7);
\node at (8,-1.5) {C : rétractée,};
\node at (8,-1.95) {sphérée (crénation)};
\legende{Aspect des hématies après 15 min : en A (NaCl 0,3 \%), l'hématie gonfle puis éclate (fantôme) ; en B (NaCl 0,9 \%), elle garde sa forme normale ; en C (NaCl 3 \%), elle se rétracte et se déforme.}
\end{tikzpicture}
\end{popfigure}

\item Solution A : \textbf{hypotonique} ; solution B : \textbf{isotonique} ; solution C : \textbf{hypertonique} (par rapport au cytoplasme des hématies).
\item Dans la solution A, le milieu extracellulaire est moins concentré que le cytoplasme : l'eau \textbf{entre} dans l'hématie par osmose ; la cellule, dépourvue de paroi, gonfle puis éclate (hémolyse). Dans la solution C, le milieu est plus concentré : l'eau \textbf{sort} de l'hématie, qui se déshydrate, se rétracte et prend un aspect crénelé (crénation).
\item Le sérum physiologique à \num{0,9}~\% est \textbf{isotonique} par rapport au plasma et aux cellules sanguines : injecté au malade, il ne provoque ni entrée ni sortie nette d'eau, donc ni hémolyse ni crénation des hématies. Il permet de réhydrater ou de diluer un médicament sans perturber l'équilibre hydrique des cellules.
\end{enumerate}
\end{corrige}

\begin{corrige}[Exercice 8 — Tableau comparatif des échanges cellulaires]
\begin{center}
\begin{tabularx}{\linewidth}{|l|X|X|X|X|}
\hline
\rowcolor{popblueL} \textbf{Mode de transport} & \textbf{Substance transportée} & \textbf{Sens du déplacement} & \textbf{Énergie (ATP)} & \textbf{Exemple physiologique} \\
\hline
Osmose & Eau (solvant) & Du milieu hypotonique vers le milieu hypertonique & Non & Réabsorption d'eau au niveau rénal \\
\hline
Diffusion simple & Petites molécules liposolubles ou gaz (O$_2$, CO$_2$) & Dans le sens du gradient de concentration & Non & Échanges O$_2$/CO$_2$ au niveau des alvéoles pulmonaires \\
\hline
Diffusion facilitée & Molécules hydrophiles (glucose, ions) via une protéine & Dans le sens du gradient de concentration & Non & Entrée du glucose dans l'hématie \\
\hline
Transport actif & Ions, molécules (via pompes) & Contre le gradient de concentration & Oui & Pompe Na$^+$/K$^+$ ; absorption du K$^+$ par les racines \\
\hline
Endocytose & Particules, macromolécules & Du milieu extracellulaire vers le cytoplasme (vésicule) & Oui & Phagocytose d'une bactérie par un macrophage \\
\hline
Exocytose & Macromolécules, résidus & Du cytoplasme vers le milieu extracellulaire & Oui & Sécrétion d'insuline par le pancréas \\
\hline
\end{tabularx}
\end{center}
\end{corrige}

\begin{corrige}[Exercice 9 — Type Bac : l'expérience de Pfeiffer et les facteurs de l'osmose]
\begin{enumerate}
\item Conditions de l'osmose : (i) deux milieux de concentrations différentes en soluté ; (ii) séparés par une \textbf{membrane semi-perméable} (perméable au solvant, imperméable au soluté).
\item La courbe $\pi = f(C)$ est une \textbf{droite passant par l'origine} (points alignés : (0~;~0), (17,1~;~1,2), \ldots, (85,5~;~6,0)).
\item Une droite passant par l'origine traduit une \textbf{proportionnalité} : $\pi = k \times C$, avec $k = \dfrac{1,2}{17,1} \approx \num{0,070}$~atm$\cdot$L$\cdot$g$^{-1}$. La pression osmotique est proportionnelle à la concentration massique (à température constante).
\item Conversion : $C_{\text{mol}} = \dfrac{C}{M} = \dfrac{C}{342}$.

\begin{center}
\begin{tabularx}{\linewidth}{|l|*{5}{X|}}
\hline
\rowcolor{popblueL} $C$ (g$\cdot$L$^{-1}$) & 17,1 & 34,2 & 51,3 & 68,4 & 85,5 \\
\hline
$C_{\text{mol}}$ (mol$\cdot$L$^{-1}$) & 0,050 & 0,100 & 0,150 & 0,200 & 0,250 \\
\hline
$\pi$ (atm) & 1,2 & 2,4 & 3,6 & 4,8 & 6,0 \\
\hline
$\pi / C_{\text{mol}}$ (L$\cdot$atm$\cdot$mol$^{-1}$) & 24,0 & 24,0 & 24,0 & 24,0 & 24,0 \\
\hline
\end{tabularx}
\end{center}

Le rapport $\dfrac{\pi}{C_{\text{mol}}}$ est \textbf{constant} et vaut 24,0~L$\cdot$atm$\cdot$mol$^{-1}$. Or $\pi = R T C_{\text{mol}}$, donc $R \times T = 24,0$~L$\cdot$atm$\cdot$mol$^{-1}$. Valeur théorique : $\num{0,082} \times 293 = \num{24,03}$~L$\cdot$atm$\cdot$mol$^{-1}$. Les deux valeurs coïncident (écart $< 0,2$~\%), ce qui valide la loi de Van't Hoff.
\item Lecture graphique : pour $\pi = 3,0$~atm, on lit $C \approx 43$~g$\cdot$L$^{-1}$. Par le calcul : $C = \dfrac{\pi}{k} = \dfrac{3,0}{\num{0,070}} \approx \num{42,8}$~g$\cdot$L$^{-1}$ (soit $C_{\text{mol}} = \dfrac{3,0}{24,0} = \num{0,125}$~mol$\cdot$L$^{-1}$, puis $\num{0,125} \times 342 = \num{42,75}$~g$\cdot$L$^{-1}$).
\item La pression osmotique d'une solution dépend de : la \textbf{concentration molaire} du soluté (proportionnalité) et de la \textbf{température absolue} (proportionnalité) ; elle ne dépend pas de la nature du soluté (loi de Van't Hoff : $\pi = RTC$).
\end{enumerate}
\textbf{Barème indicatif (10 pts)} : conditions 1 pt ; courbe soignée (axes, échelles, titre) 2 pts ; relation 1 pt ; tableau de conversion et constat 2 pts ; calcul de $RT$ et comparaison 1,5 pt ; détermination de X 1,5 pt ; conclusion 1 pt.
\end{corrige}

\begin{attention}[Points de vigilance]
La courbe doit être tracée avec des axes gradués, légendés (grandeur + unité) et un titre. Pour la question 4, exiger la conversion massique $\rightarrow$ molaire \emph{avant} le calcul du rapport : utiliser directement les g$\cdot$L$^{-1}$ est l'erreur classique.
\end{attention}

\begin{corrige}[Exercice 10 — Type Bac : absorption du potassium par des cellules racinaires]
\begin{enumerate}
\item Courbes :

\begin{popfigure}
\begin{tikzpicture}
\begin{axis}[width=0.95\linewidth, height=7cm, xlabel={Temps (min)}, ylabel={K$^+$ absorbé ($\mu$mol$\cdot$g$^{-1}$)}, xmin=0, xmax=130, ymin=0, ymax=16, xtick={0,20,40,60,90,120}, legend pos=north west, legend style={font=\scriptsize}, grid=both, grid style={popline!40}]
\addplot[color=popblue, thick, mark=*] coordinates {(0,0)(20,4)(40,7.5)(60,10)(90,12.5)(120,14)};
\addlegendentry{Lot 1 (témoin, \SI{25}{\celsius}, air)}
\addplot[color=poporange, thick, mark=square*] coordinates {(0,0)(20,1.2)(40,1.6)(60,1.7)(90,1.7)(120,1.7)};
\addlegendentry{Lot 2 (KCN)}
\addplot[color=popgreen, thick, mark=triangle*] coordinates {(0,0)(20,1)(40,1.8)(60,2.3)(90,2.8)(120,3)};
\addlegendentry{Lot 3 (\SI{4}{\celsius})}
\end{axis}
\end{tikzpicture}
\legende{Cinétique d'absorption des ions K$^+$ par des racines de maïs dans trois conditions expérimentales.}
\end{popfigure}

\item Le lot 2 (KCN) absorbe beaucoup moins de K$^+$ que le témoin et atteint un plateau très bas (1,7~$\mu$mol$\cdot$g$^{-1}$). Le cyanure \textbf{bloque la respiration cellulaire} (inhibition de la chaîne respiratoire mitochondriale), donc la production d'ATP. L'effondrement de l'absorption en présence de KCN montre que l'absorption du K$^+$ \textbf{exige de l'énergie} : c'est un phénomène \textbf{actif}.
\item À \SI{4}{\celsius} (lot 3), l'absorption est fortement ralentie par rapport à \SI{25}{\celsius}. Le facteur mis en évidence est la \textbf{température} : le froid ralentit l'activité enzymatique (dont les ATPases des pompes) et la respiration cellulaire, donc diminue la fourniture d'ATP et la vitesse du transport.
\item Concentration extracellulaire : \num{0,5}~mmol$\cdot$L$^{-1}$ ; concentration intracellulaire : \num{80}~mmol$\cdot$L$^{-1}$, soit \textbf{160 fois plus} à l'intérieur. Les ions K$^+$ passent donc du milieu le moins concentré vers le milieu le plus concentré : l'absorption s'effectue \textbf{contre le gradient de concentration}.
\item Il s'agit d'un \textbf{transport actif}. Trois conditions nécessaires : (i) une \textbf{protéine transporteuse} membranaire spécifique (pompe) ; (ii) une \textbf{dépense d'énergie} fournie par l'hydrolyse de l'ATP ; (iii) un déplacement \textbf{contre le gradient} de concentration (électrochimique).
\item La faible absorption résiduelle du lot 2 (1,7~$\mu$mol$\cdot$g$^{-1}$, indépendante de l'ATP) correspond à la \textbf{diffusion} (passive, éventuellement facilitée) : elle s'effectue dans le sens du gradient électrochimique, sans énergie, et s'arrête vite car le gradient s'équilibre rapidement (plateau).
\end{enumerate}
\textbf{Barème indicatif (10 pts)} : courbes (axes, échelles, légende) 2,5 pts ; comparaison lot 1/lot 2 et rôle du KCN 2 pts ; effet de la température 1,5 pt ; démonstration du contre-gradient 1,5 pt ; type de transport et conditions 1,5 pt ; transport résiduel 1 pt.
\end{corrige}

\begin{attention}[Points de vigilance]
Citer explicitement le rôle du KCN (inhibiteur de la respiration $\Rightarrow$ chute d'ATP) avant de conclure au caractère actif. Pour la question 4, la comparaison chiffrée des deux concentrations ($\num{80} \gg \num{0,5}$~mmol$\cdot$L$^{-1}$) est exigée : affirmer « contre le gradient » sans justification numérique ne rapporte pas tous les points.
\end{attention}

\begin{corrige}[Exercice 11 — Type Bac : phagocytose et défense de l'organisme]
\begin{enumerate}
\item Le phénomène illustré est la \cle{phagocytose}, qui appartient à l'\textbf{endocytose} (endocytose de particules solides, par opposition à la pinocytose qui concerne des gouttelettes liquides).
\item Déroulement :
\begin{enumerate}
\item \textbf{Adhésion} : la bactérie se fixe à la surface de la membrane plasmique du macrophage (reconnaissance par des récepteurs).
\item \textbf{Émission de pseudopodes} : le macrophage émet des prolongements cytoplasmiques qui entourent progressivement la bactérie.
\item \textbf{Formation du phagosome} : les pseudopodes fusionnent ; la bactérie se retrouve enfermée dans une vésicule intracytoplasmique, le \textbf{phagosome}.
\item \textbf{Digestion} : le phagosome fusionne avec des lysosomes pour former un \textbf{phagolysosome}, dont les enzymes hydrolytiques digèrent la bactérie.
\end{enumerate}
\item Schéma de l'étape (c) :

\begin{popfigure}
\begin{tikzpicture}[font=\small]
\draw[fill=popcream, draw=popdark, thick] (0,0) ellipse (4.2 and 2.9);
\draw[draw=popink, thick] (0,0) ellipse (4.35 and 3.05);
\draw[fill=popblueL, draw=popblue, thick] (-1.6,0.4) circle (1.0);
\draw[fill=popblue] (-1.6,0.4) circle (0.28);
\draw[fill=popgreenL, draw=popgreen, thick] (1.8,-0.6) circle (0.95);
\draw[fill=poporange, draw=popdark] (1.8,-0.6) circle (0.4);
\foreach \a in {0,60,120,180,240,300} \draw[popdark, thick] ({1.8+0.4*cos(\a)},{-0.6+0.4*sin(\a)}) -- ({1.8+0.62*cos(\a)},{-0.6+0.62*sin(\a)});
\node at (6.6,2.6) {membrane plasmique};
\draw[->] (5.4,2.5) -- (4.0,2.35);
\node at (-5.6,1.4) {noyau};
\draw[->] (-4.6,1.3) -- (-2.4,0.7);
\node at (6.3,-1.5) {phagosome contenant la bactérie};
\draw[->] (5.0,-1.35) -- (2.6,-0.85);
\node at (-5.6,-1.8) {cytoplasme};
\draw[->] (-4.5,-1.7) -- (-2.6,-1.3);
\legende{Étape (c) de la phagocytose : la bactérie (bâtonnet orange) est enfermée dans le phagosome, vésicule intracytoplasmique délimitée par une membrane issue de la membrane plasmique.}
\end{tikzpicture}
\end{popfigure}

\item Les enzymes digestives proviennent des \textbf{lysosomes}. Ce sont des protéines synthétisées par les \textbf{ribosomes du réticulum endoplasmique granuleux}, puis mises en réserve dans les lysosomes après maturation par l'appareil de Golgi.
\item L'énergie est fournie par les \textbf{mitochondries} (production d'ATP par respiration cellulaire). Les macrophages, qui réalisent des phagocytoses fréquentes et énergivores (déplacement, formation de pseudopodes, fusion de vésicules), possèdent donc des mitochondries \textbf{abondantes}.
\item Les résidus sont rejetés par \textbf{exocytose} : la vésicule contenant les déchets fusionne avec la membrane plasmique et déverse son contenu dans le milieu extracellulaire.
\item Deux arguments sur l'importance des échanges cellulaires :
\begin{itemize}
\item Ils assurent la \textbf{défense de l'organisme} : l'endocytose (phagocytose) permet l'élimination des agents pathogènes, ce qui est vital lors d'infections comme le paludisme.
\item Ils permettent le \textbf{renouvellement et l'épuration du milieu intérieur} : l'exocytose évacue les résidus de digestion, maintenant l'homéostasie cellulaire (on peut aussi citer la sécrétion d'anticorps ou d'hormones).
\end{itemize}
\end{enumerate}
\textbf{Barème indicatif (10 pts)} : identification du phénomène 1 pt ; quatre étapes avec vocabulaire exigé 2,5 pts ; schéma annoté 2 pts ; lysosomes et origine des enzymes 1,5 pt ; mitochondrie et justification 1 pt ; exocytose 1 pt ; deux arguments 1 pt.
\end{corrige}

\begin{attention}[Points de vigilance]
Le vocabulaire scientifique est exigé : \emph{pseudopodes}, \emph{phagosome}, \emph{phagolysosome} doivent apparaître dans la description. Ne pas confondre phagocytose (particules solides) et pinocytose (liquides), ni endocytose (entrée) et exocytose (sortie).
\end{attention}

\newpage

\coursec{Fiche bilan}

\begin{popfigure}
\begin{tikzpicture}[mindmap, grow cyclic,
  every node/.style={concept, text=white, font=\small\bfseries},
  root concept/.append style={concept color=popink, font=\large\bfseries},
  level 1/.append style={level distance=3.6cm, sibling angle=60},
  level 1 concept/.append style={font=\footnotesize\bfseries},
  scale=0.92, transform shape]
\node[root concept]{La cellule animale et ses échanges}
  child[concept color=popblue]{ node[align=center]{Organisation générale\\MO et ME} }
  child[concept color=popgreen]{ node[align=center]{Organites\\et leurs rôles} }
  child[concept color=poporange]{ node[align=center]{Échanges d'eau\\osmose -- dialyse} }
  child[concept color=poppurple]{ node[align=center]{Transport passif\\diffusion simple\\diffusion facilitée} }
  child[concept color=poppink]{ node[align=center]{Transport actif\\pompe Na$^+$/K$^+$} }
  child[concept color=popgold]{ node[align=center]{Endocytose\\Exocytose} };
\end{tikzpicture}
\legende{Carte mentale de la leçon : autour de la cellule animale (nœud central), six branches résument les savoirs essentiels --- organisation, organites, osmose et dialyse, transports passif et actif, endocytose et exocytose.}
\end{popfigure}

\begin{bilan}
\textbf{1. Définitions clés à connaître par c\oe ur}
\begin{itemize}
  \item \cle{Cellule} : unité structurale et fonctionnelle de tout être vivant ; chez l'animal, elle comprend la \cle{membrane plasmique}, le \cle{cytoplasme} et le \cle{noyau}.
  \item \cle{Osmose} : diffusion de l'eau (solvant) à travers une membrane semi-perméable, du milieu \textbf{hypotonique} (moins concentré en solutés) vers le milieu \textbf{hypertonique} (plus concentré).
  \item \cle{Dialyse} : diffusion des solutés à travers une membrane perméable, du milieu le plus concentré vers le milieu le moins concentré.
  \item \cle{Diffusion simple} : passage de molécules (O$_2$, CO$_2$, lipides) à travers la bicouche lipidique, selon le gradient de concentration, sans protéine ni ATP.
  \item \cle{Diffusion facilitée} : passage de molécules hydrophiles (glucose, ions) via des \textbf{protéines transporteurs} (canaux, perméases), selon le gradient, sans ATP.
  \item \cle{Transport actif} : passage \textbf{contre le gradient} de concentration, nécessitant une pompe membranaire et de l'\textbf{ATP} (ex. pompe Na$^+$/K$^+$ : 3 Na$^+$ sortent, 2 K$^+$ entrent par ATP hydrolysé).
  \item \cle{Endocytose} : entrée de particules par invagination membranaire (phagocytose : solides ; pinocytose : liquides).
  \item \cle{Exocytose} : expulsion de substances par fusion d'une vésicule avec la membrane plasmique (ex. sécrétion d'insuline).
\end{itemize}

\textbf{2. Organites et rôles (tableau à maîtriser)}
\begin{center}
\begin{tabularx}{\linewidth}{|l|X|}\hline
\rowcolor{popblueL} \textbf{Organite} & \textbf{Rôle essentiel} \\ \hline
Membrane plasmique & Enveloppe sélective ; siège de tous les échanges \\ \hline
Noyau & Contient l'ADN ; commande les activités cellulaires \\ \hline
Mitochondrie & Respiration cellulaire : production d'ATP \\ \hline
Réticulum endoplasmique granuleux & Synthèse des protéines (ribosomes fixés) \\ \hline
Réticulum endoplasmique lisse & Synthèse des lipides, détoxication \\ \hline
Appareil de Golgi & Maturation, conditionnement et sécrétion \\ \hline
Lysosome & Digestion intracellulaire (enzymes hydrolasiques) \\ \hline
Ribosome & Assemblage des protéines \\ \hline
Centrioles & Organisation de la division cellulaire \\ \hline
\end{tabularx}
\end{center}

\textbf{3. Expériences fondamentales (exigibles au Bac)}
\begin{itemize}
  \item \textbf{Expérience de Dutrochet} : une vessie de porc remplie d'eau sucrée, plongée dans l'eau pure, se gonfle $\Rightarrow$ l'eau entre par osmose vers le milieu le plus concentré.
  \item \textbf{Expérience de Pfeiffer} (osmomètre) : la montée du liquide dans le tube est proportionnelle à la concentration du soluté $\Rightarrow$ mise en évidence de la \cle{pression osmotique}.
  \item \textbf{Hématies} : en milieu hypotonique $\rightarrow$ turgescence puis hémolyse ; en milieu hypertonique $\rightarrow$ plasmolyse (crénation) ; en milieu isotonique ($\approx$ \SI{9}{g/L} de NaCl) $\rightarrow$ état normal.
\end{itemize}

\textbf{4. Méthodes express}
\begin{itemize}
  \item Prévoir le sens d'un échange d'eau : comparer les concentrations des deux milieux $\rightarrow$ l'eau va toujours vers le milieu \textbf{hypertonique}.
  \item Distinguer passif/actif : deux questions --- \emph{le transport suit-il le gradient ?} \emph{consomme-t-il de l'ATP ?} Oui/Non $\rightarrow$ passif ; Non/Oui $\rightarrow$ actif.
  \item Schéma de cellule : toujours légender (membrane, cytoplasme, noyau, organites) et indiquer le grossissement.
\end{itemize}
\end{bilan}

\begin{aretenir}[Checklist avant le Bac]
\begin{itemize}
  \item $\square$ Je sais schématiser et légender une cellule animale vue en MO et en ME.
  \item $\square$ Je connais le rôle de chaque organite (mitochondrie, REG, REL, Golgi, lysosome, ribosome).
  \item $\square$ Je sais définir l'osmose et la dialyse sans les confondre.
  \item $\square$ Je sais interpréter l'expérience de Dutrochet (vessie qui se gonfle).
  \item $\square$ Je sais interpréter l'expérience de Pfeiffer (osmomètre, pression osmotique).
  \item $\square$ Je prévois le comportement d'une hématie en milieu hypo-, iso- ou hypertonique.
  \item $\square$ Je distingue diffusion simple et diffusion facilitée (rôle des protéines transporteurs).
  \item $\square$ Je sais que le transport actif va contre le gradient et consomme de l'ATP (pompe Na$^+$/K$^+$).
  \item $\square$ Je différencie phagocytose, pinocytose et exocytose, avec un exemple chacun.
  \item $\square$ Je sais citer l'importance des échanges cellulaires : nutrition, respiration, excrétion, défense immunitaire.
  \item $\square$ Je rédige mes interprétations en trois temps : observation $\rightarrow$ hypothèse $\rightarrow$ conclusion.
  \item $\square$ Je vérifie que chaque schéma rendu est légendé et titré.
\end{itemize}
\end{aretenir}

\findecours

\end{document}
